% Options for packages loaded elsewhere
\PassOptionsToPackage{unicode}{hyperref}
\PassOptionsToPackage{hyphens}{url}
\PassOptionsToPackage{dvipsnames,svgnames,x11names}{xcolor}
\documentclass[
  11pt,
  a4paper,
]{article}
\usepackage{xcolor}
\usepackage[margin=0.75in]{geometry}
\usepackage{amsmath,amssymb}
\setcounter{secnumdepth}{-\maxdimen} % remove section numbering
\usepackage{iftex}
\ifPDFTeX
  \usepackage[T1]{fontenc}
  \usepackage[utf8]{inputenc}
  \usepackage{textcomp} % provide euro and other symbols
\else % if luatex or xetex
  \usepackage{unicode-math} % this also loads fontspec
  \defaultfontfeatures{Scale=MatchLowercase}
  \defaultfontfeatures[\rmfamily]{Ligatures=TeX,Scale=1}
\fi
\usepackage{lmodern}
\ifPDFTeX\else
  % xetex/luatex font selection
  \setmainfont[]{Times New Roman}
  \setsansfont[]{Helvetica}
\fi
% Use upquote if available, for straight quotes in verbatim environments
\IfFileExists{upquote.sty}{\usepackage{upquote}}{}
\IfFileExists{microtype.sty}{% use microtype if available
  \usepackage[]{microtype}
  \UseMicrotypeSet[protrusion]{basicmath} % disable protrusion for tt fonts
}{}
\makeatletter
\@ifundefined{KOMAClassName}{% if non-KOMA class
  \IfFileExists{parskip.sty}{%
    \usepackage{parskip}
  }{% else
    \setlength{\parindent}{0pt}
    \setlength{\parskip}{6pt plus 2pt minus 1pt}}
}{% if KOMA class
  \KOMAoptions{parskip=half}}
\makeatother
\usepackage{longtable,booktabs,array}
\newcounter{none} % for unnumbered tables
\usepackage{calc} % for calculating minipage widths
% Correct order of tables after \paragraph or \subparagraph
\usepackage{etoolbox}
\makeatletter
\patchcmd\longtable{\par}{\if@noskipsec\mbox{}\fi\par}{}{}
\makeatother
% Allow footnotes in longtable head/foot
\IfFileExists{footnotehyper.sty}{\usepackage{footnotehyper}}{\usepackage{footnote}}
\makesavenoteenv{longtable}
\setlength{\emergencystretch}{3em} % prevent overfull lines
\providecommand{\tightlist}{%
  \setlength{\itemsep}{0pt}\setlength{\parskip}{0pt}}
\usepackage[]{natbib}
\bibliographystyle{plainnat}
\usepackage{fontspec}
\usepackage{geometry}
\usepackage{fancyhdr}
\usepackage{xcolor}
\usepackage{natbib}
\usepackage{gb4e}
\noautomath
\newcommand{\tdim}[1]{\textit{#1}}
\newcommand{\tdimword}[1]{\textit{#1}}
\let\oldeachwordone\eachwordone
\renewcommand{\eachwordone}{\oldeachwordone\hspace{0.3em}}
\let\oldeachwordtwo\eachwordtwo
\renewcommand{\eachwordtwo}{\oldeachwordtwo\hspace{0.3em}}
\renewcommand{\textsc}[1]{{\footnotesize\MakeUppercase{#1}}}
\geometry{margin=0.75in}
\setmainfont{Times New Roman}
\setsansfont{Helvetica}
\pagestyle{fancy}
\fancyhf{}
\fancyfoot[C]{\thepage}
\setcitestyle{authoryear,round,semicolon}
\newcommand{\frontmattersection}[1]{\section*{#1}\addcontentsline{toc}{section}{#1}}
\newcommand{\reviewwarning}[1]{\par\medskip\noindent\textbf{#1}\par\medskip}
\newcommand{\reviewwarninginline}[1]{\textbf{#1}\par\smallskip}
\newcounter{reviewchapter}
\newcommand{\reviewchapterbreak}{\stepcounter{reviewchapter}\setcounter{xnumi}{0}}
\renewcommand{\thexnumi}{\arabic{reviewchapter}.\arabic{xnumi}}
\exewidth{(4.99)}
\usepackage{bookmark}
\IfFileExists{xurl.sty}{\usepackage{xurl}}{} % add URL line breaks if available
\urlstyle{same}
\hypersetup{
  colorlinks=true,
  linkcolor={Maroon},
  filecolor={Maroon},
  citecolor={Blue},
  urlcolor={Blue},
  pdfcreator={LaTeX via pandoc}}

\author{}
\date{}

\begin{document}

{
\setcounter{tocdepth}{3}
\tableofcontents
}
\frontmattersection{Review preview status}

This is a review preview, not a finished grammar. It is assembled from
first-pass publication-review slices and is controlled by
\texttt{output/publication\_review/whole\_grammar\_coverage\_checkpoint\_after\_transitivity.md},
\texttt{output/publication\_review/whole\_grammar\_coverage\_checkpoint\_after\_reduplication.md},
\texttt{output/publication\_review/whole\_grammar\_coverage\_audit.md},
\texttt{docs/SKELETON\_GRAMMAR.md},
\texttt{docs/grammar/GRAMMAR\_SOURCE\_INVENTORY.md}, and
\texttt{PROGRESS.md}.

\texttt{output/publication\_review/review\_notes\_transitivity.md}
brought the transitivity packet to review-note maturity, and the
post-transitivity checkpoint now treats the packet set as stable enough
for a review preview assembled from actual slice prose. This document is
not a new grammar slice, not a dictionary slice, and not a human-review
packet. It is intended to help human review and direct editing, not to
certify completion.

Many sections are deliberately narrow. Missing or blocked domains are
marked explicitly. The PDF built from this assembly is a review preview
PDF, not a final publication PDF.

\frontmattersection{PDF/build status}

This preview is reproducible from committed sources with
\texttt{python3\ scripts/assemble\_publication\_review\_preview.py}. The
script writes
\texttt{output/publication\_review/assembled\_grammar\_review\_preview.md},
generates
\texttt{output/publication\_review/assembled\_grammar\_review\_preview.tex}
through Pandoc plus natbib/BibTeX citation processing, and compiles
\texttt{output/publication\_review/assembled\_grammar\_review\_preview.pdf}
with XeLaTeX while routing publication-review example blocks through the
shared analyzer and gb4e interlinear machinery.

The assembly reuses current repository conventions where practical: the
repository bibliography in \texttt{literature/bibliography.bib}, XeLaTeX
compilation and the \texttt{Times\ New\ Roman} / \texttt{Helvetica} font
pair already used in \texttt{scripts/export\_interlinear.py}, the shared
Bible/analyzer helpers in \texttt{scripts/interlinear\_latex.py}, and
the same 0.75-inch page-margin convention for generated TeX output.

\section*{Abbreviations}
\addcontentsline{toc}{section}{Abbreviations}

\subsection*{Standard Leipzig Glossing Abbreviations}
\begin{tabular}{@{}ll@{\hspace{2em}}ll@{}}
\textsc{1} & first person & \textsc{imp} & imperative \\
\textsc{2} & second person & \textsc{incl} & inclusive \\
\textsc{3} & third person & \textsc{ins} & instrumental \\
\textsc{abil} & abilitative & \textsc{intens} & intensive \\
\textsc{abl} & ablative & \textsc{irr} & irrealis \\
\textsc{all} & allative & \textsc{iter} & iterative \\
\textsc{appl} & applicative & \textsc{loc} & locative \\
\textsc{caus} & causative & \textsc{neg} & negative \\
\textsc{com} & comitative & \textsc{nmlz} & nominalizer \\
\textsc{compl} & completive & \textsc{pfv} & perfective \\
\textsc{cont} & continuative & \textsc{pl} & plural \\
\textsc{cvb} & converb & \textsc{poss} & possessive \\
\textsc{dat} & dative & \textsc{prosp} & prospective \\
\textsc{dimin} & diminutive & \textsc{prox} & proximal \\
\textsc{erg} & ergative & \textsc{q} & question particle \\
\textsc{excl} & exclusive & \textsc{recip} & reciprocal \\
\textsc{exp} & experiential & \textsc{refl} & reflexive \\
\textsc{gen} & genitive & \textsc{rel} & relativizer \\
\textsc{hab} & habitual & \textsc{res} & resultative \\
\textsc{horiz} & horizontal & \textsc{sg} & singular \\
\textsc{imm} & immediate & \textsc{top} & topic \\
\end{tabular}

\vspace{1em}
\subsection*{Tedim Chin-Specific Conventions}
\begin{tabular}{@{}ll@{}}
\textsc{1sg$\rightarrow$3} & first singular acting on third person \\
\textsc{2$\rightarrow$1} & second person acting on first person \\
\textsc{3$\rightarrow$1} & third person acting on first person \\
\textsc{ag} & agent (nominalizer) \\
\textsc{aug} & augmentative (big/great) \\
\textsc{dist} & distal/anaphoric demonstrative \\
\textsc{down} & downward directional \\
\textsc{hab.cont} & habitual continuative \\
\textsc{i} & verb form I (citation form) \\
\textsc{ii} & verb form II (dependent form) \\
\textsc{more} & comparative (more X) \\
\textsc{neg.abil} & negative abilitative (cannot) \\
\textsc{prox} & proximal demonstrative \\
\textsc{toward} & goal directional \\
\textsc{up} & upward directional \\
\end{tabular}

\vspace{1em}
\noindent\textbf{Verb Forms:} Tedim Chin verbs have two conjugation forms.
\textsc{i} (Form I) is the citation/independent form; \textsc{ii} (Form II)
appears in dependent clauses and certain constructions.

\frontmattersection{Known narrow-slice limitations}

\begin{itemize}
\tightlist
\item
  VP structure / suffix stacking: currently anchored by
  \texttt{bawlzoding}.
\item
  derivation / valency: currently anchored by \texttt{-sak}.
\item
  prefix/agreement: currently anchored by \texttt{kanei\ /\ kainn}.
\item
  clause linkage: currently anchored by \texttt{ciangin}.
\item
  nominalization: currently anchored by \texttt{-na\ /\ bawlna}.
\item
  NP structure / possession: currently anchored by \texttt{hih\ mite},
  \texttt{mi\ khat}, \texttt{mi\ khempeuh}.
\item
  noun domain: currently anchored by \texttt{gam} and
  \texttt{aksi\ /\ aksi-te}.
\item
  reduplication: currently anchored by \texttt{mahmah\ /\ taktak}, with
  \texttt{peuhpeuh} secondary.
\item
  transitivity: currently anchored by \texttt{sih\ /\ suak} versus
  \texttt{hawl\ /\ en}.
\end{itemize}

\frontmattersection{Major unresolved domains}

\begin{itemize}
\tightlist
\item
  {[}MAJOR GAP: phonology/tone remains blocked or theory-heavy.{]}
\item
  {[}MAJOR GAP: verb paradigms remain report-backed but not
  packet-shaped.{]}
\item
  {[}MAJOR GAP: broader discourse remains partly surfaced and
  boundary-heavy.{]}
\item
  {[}MAJOR GAP: analyzer-gap topics remain cross-cutting blockers.{]}
\end{itemize}

Second-pass expansions such as \texttt{-pih}, \texttt{ki-}, hong-/kong-,
switch reference, relative clauses, transparent compounds, wider
reduplication, and labile or ambitransitive transitivity remain outside
this first-pass assembled review preview.

\reviewchapterbreak

\section{Phonology and tone}\label{phonology-and-tone}

{[}MAJOR GAP: phonology/tone remains blocked or theory-heavy.{]}

The controlling checkpoints and audit still treat phonology/tone as
blocked or theory-heavy, so no publication-review grammar slice is
inlined here yet.

\reviewchapterbreak

\section{Deixis, pronouns, and nominal
domain}\label{deixis-pronouns-and-nominal-domain}

\subsection{Demonstratives / deixis}\label{demonstratives-deixis}

\emph{Source slice:
\texttt{output/publication\_review/grammar\_demonstratives\_print\_slice.md}}

\subsubsection{Scope}\label{scope}

This account offers a short treatment of Tedim demonstratives and
deixis. It is intentionally narrow. It focuses on a small set of
manually checked Bible examples and on the interaction between the core
demonstratives \tdim{hih} and \tdim{tua}, their plural forms, and a few
strong constructional extensions. It does not attempt a full treatment
of copulas, sentence-final particles, directionals, or broad discourse
structure.

\subsubsection{Overview of the demonstrative
system}\label{overview-of-the-demonstrative-system}

The current Bible corpus and the literature agree on a compact
demonstrative core: proximal \tdim{hih} and distal \tdim{tua}
\citep{henderson1965, zamngaihcing2017}. In Bible usage, however,
\tdim{tua} is not merely a spatial ``that''. It is also a strongly
anaphoric and discourse-linking form, especially in narrative and
temporal sequencing. The most economical description is therefore a
system with \tdim{hih} as the core proximal demonstrative and \tdim{tua}
as the core distal/anaphoric demonstrative.

The plural forms \tdim{hihte} and \tdim{tuate} behave transparently as
DEM + \tdim{-te}, and the Bible also gives strong constructional
material such as \tdim{hih bangin}, \tdim{tua bangin},
\tdim{tua ciangin}, and \tdim{tua ahih ciangin}. One caution remains
explicit throughout this slice: Zam Ngaih Cing gives distal \tdim{huā},
while the Bible corpus overwhelmingly uses \tdim{tua} as the ordinary
distal/anaphoric form \citep{zamngaihcing2017}. The relation is
plausible, but it is not resolved here as a settled orthographic
identity.

\subsubsection{\texorpdfstring{Core demonstratives: \tdim{hih} and
\tdim{tua}}{Core demonstratives:  and }}\label{core-demonstratives-and}

\tdim{Hih} is the clearest proximal form in the present corpus. It works
both as an identificational demonstrative and as an adnominal
determiner.

\begin{exe}
\ex \label{ex:dem-hih}
\gll \tdimword{Híh} \tdimword{pèn} \tdimword{Adam'} \tdimword{suanlekhakte'} \tdimword{laibu} \tdimword{àhì} \tdimword{hì.} \\
     \textsc{prox} \textsc{top} \textsc{adam}.\textsc{poss} genealogy-\textsc{pl}.\textsc{poss} book be.\textsc{3sg} \textsc{decl} \\
\glt 'This is the book of the generations of Adam.' (Genesis 5:1)
\end{exe}

Genesis 5:1 is especially useful because it shows \tdim{hih} as the
unmarked proximal form in identificational prose. Genesis 9:12
\tdim{Hih pen ... thuciamna lim ahi hi} confirms the same pattern, while
Exodus 32:9 \tdim{Hih mite ka mu zo hi} shows \tdim{hih} before a noun
phrase in ordinary adnominal use.

\tdim{Tua} belongs to the same system, but its range in the Bible is
broader. It clearly marks distal or previously activated referents, and
it also extends naturally into anaphoric and discourse-linking uses.

\begin{exe}
\ex \label{ex:dem-tua}
\gll \tdimword{Pasian} \tdimword{ìn,} \tdimword{tuite'} \tdimword{lai-zang-àh} \tdimword{vàn} \tdimword{kuum-pì} \tdimword{ōm} \tdimword{hèn} \tdimword{lá,} \tdimword{túa} \tdimword{vàn} \tdimword{kuum-pì} \tdimword{ìn} \tdimword{tùi} \tdimword{lè} \tdimword{tùi} \tdimword{kī-khen-sàk} \tdimword{hèn} \tdimword{cí} \tdimword{hì.} \\
     God \textsc{erg} water-\textsc{pl}.\textsc{poss} middle-side-\textsc{loc} sky bow-big exist \textsc{juss} take \textsc{dist} sky bow-big \textsc{erg} water and water \textsc{refl}-divide-\textsc{caus} \textsc{juss} say \textsc{decl} \\
\glt 'And God said, Let there be a firmament in the midst of the waters, and let it divide the waters from the waters.' (Genesis 1:6)
\end{exe}

Genesis 1:6 is a clean adnominal example. Genesis 21:27
\tdim{tua mi nihte} is equally good for a more discourse-linked noun
phrase. The important descriptive point is that the corpus keeps using
\tdim{tua} where English may translate ``that'', ``the same'', or simply
a contextually given referent.

\subsubsection{\texorpdfstring{Plural forms: \tdim{hihte} and
\tdim{tuate}}{Plural forms:  and }}\label{plural-forms-and}

The plural demonstratives are straightforward. The Bible corpus strongly
supports \tdim{hihte} and \tdim{tuate} as transparent plural extensions
of the singular forms.

\begin{exe}
\ex \label{ex:dem-hihte}
\gll \tdimword{Israel} \tdimword{ìn} \tdimword{Josef'} \tdimword{tapa-tè} \tdimword{à} \tdimword{mùh} \tdimword{cìang-ìn} \tdimword{àmàh} \tdimword{ìn,} \tdimword{híh-tè} \tdimword{kua} \tdimword{àhì} \tdimword{hìam} \tdimword{à} \tdimword{cí} \tdimword{hì.} \\
     \textsc{israel} \textsc{erg} \textsc{josef}.\textsc{poss} son-\textsc{pl} \textsc{3sg} see.\textsc{ii} then-\textsc{erg} \textsc{3sg}.\textsc{pro} \textsc{erg} \textsc{hihte} who be.\textsc{3sg} \textsc{q} \textsc{3sg} say \textsc{decl} \\
\glt 'And Israel beheld Joseph's sons, and said, Who are these?' (Genesis 48:8)
\end{exe}

\begin{exe}
\ex \label{ex:dem-tuate}
\gll \tdimword{Mual} \tdimword{lìan-tè} \tdimword{tùngàh} \tdimword{tùi} \tdimword{om-tò} \tdimword{lái} \tdimword{à,} \tdimword{túa-tè} \tdimword{tùngàh} \tdimword{pì} \tdimword{sàwm-nìh} \tdimword{lè} \tdimword{nìh} \tdimword{thuk-ìn} \tdimword{tuumcip} \tdimword{hì.} \\
     mountain big-\textsc{pl} on-\textsc{loc} water be-over midst \textsc{3sg} \textsc{dist}-\textsc{pl} on-\textsc{loc} grandmother ten-two and two deep-\textsc{inst} water \textsc{decl} \\
\glt 'Fifteen cubits upward did the waters prevail; and the mountains were covered.' (Genesis 7:20)
\end{exe}

Genesis 10:20 \tdim{Hihte pen ... Ham’ suanlekhakte ahi hi} shows the
proximal plural in topic-like prose, while Genesis 2:19
\tdim{mipa in tuate bang a ci hiam} shows the distal plural as an
ordinary pronoun. Nothing in the present corpus suggests that these need
any more complex treatment than DEM + \tdim{-te}.

\subsubsection{Adnominal and pronominal
uses}\label{adnominal-and-pronominal-uses}

The same forms serve both as adnominal modifiers and as independent
pronouns. Adnominal use is easy to see in Exodus 32:9 \tdim{Hih mite},
Genesis 1:6 \tdim{tua van kuumpi}, and Genesis 21:27
\tdim{tua mi nihte}. Pronominal use is equally clear in Genesis 5:1
\tdim{Hih pen}, Genesis 48:8 \tdim{Hihte kua ahi hiam?}, and Genesis
2:19 \tdim{tuate bang a ci hiam}.

The Bible also shows productive matter/topic phrases such as
\tdim{hih thu} and \tdim{tua thu}. In Genesis 24:9
\tdim{hih thu tawh kisai-in} points to the present matter under
discussion, while Genesis 6:6 \tdim{tua thu in ama lungsim dahsak hi}
and Exodus 2:15 \tdim{Faro in tua thu a zak ciangin} show anaphoric
reference to something just reported. These are important for the
grammar chapter, even if they do not have to be first-round dictionary
headwords.

\subsubsection{Discourse and temporal
deixis}\label{discourse-and-temporal-deixis}

This is where \tdim{tua} most clearly exceeds a merely spatial
description. In Bible narrative, \tdim{tua} is a major discourse and
temporal linker.

\begin{exe}
\ex \label{ex:dem-tua-ciangin}
\gll \tdimword{Pasian} \tdimword{ìn,} \tdimword{khuavak} \tdimword{ōm} \tdimword{hèn} \tdimword{cí} \tdimword{hì;} \tdimword{túa} \tdimword{cìang-ìn} \tdimword{khuavak} \tdimword{ōm} \tdimword{pah} \tdimword{hì.} \\
     God \textsc{erg} light exist \textsc{juss} say \textsc{decl} \textsc{dist} then-\textsc{erg} light exist do.so \textsc{decl} \\
\glt 'And God said, Let there be light: and there was light.' (Genesis 1:3)
\end{exe}

\tdim{Tua ciangin} works naturally as ``then / at that time'' in event
sequencing. Genesis 5:5 repeats the same pattern in obituary-style
narrative: \tdim{tua ciangin amah si hi}.

\tdim{Tua ahih ciangin} is even more explicitly a discourse bridge,
linking a prior state or event to the next clause.

\begin{exe}
\ex \label{ex:dem-tua-ahih-ciangin}
\gll \tdimword{Túa} \tdimword{àhìh} \tdimword{cìang-ìn} \tdimword{Topa} \tdimword{Pasian} \tdimword{ìn} \tdimword{mipa} \tdimword{ihmut} \tdimword{sùak} \tdimword{màhmàh} \tdimword{sàk} \tdimword{à,} \tdimword{àmà} \tdimword{ihmut} \tdimword{kal-ìn} \tdimword{à} \tdimword{nak-gùh-tè} \tdimword{khàt} \tdimword{lá-ìn} \tdimword{túa} \tdimword{mùn} \tdimword{pèn} \tdimword{sá-tak} \tdimword{tàwh} \tdimword{à} \tdimword{dimsak} \tdimword{hì.} \\
     \textsc{dist} be.\textsc{3sg}.\textsc{rel} then-\textsc{erg} Lord God \textsc{erg} man deep.sleep wither very \textsc{caus} \textsc{3sg} \textsc{3sg}.\textsc{poss} deep.sleep middle-\textsc{cvb} \textsc{3sg} nose-bone-\textsc{pl} one take-\textsc{erg} \textsc{dist} place \textsc{top} flesh-piece \textsc{com} \textsc{3sg} fill \textsc{decl} \\
\glt 'And the LORD God caused a deep sleep to fall upon Adam, and he slept: and he took one of his ribs, and closed up the flesh instead thereof.' (Genesis 2:21)
\end{exe}

Genesis 1:21 and Genesis 2:19 confirm that \tdim{tua ahih ciangin} is
one of the strongest discourse-transition frames in the present corpus.
It is better understood as a fixed demonstrative-temporal construction
than as a purely compositional spatial phrase.

\subsubsection{\texorpdfstring{Manner constructions with
\tdim{bangin}}{Manner constructions with }}\label{manner-constructions-with}

The Bible also gives strong manner and discourse constructions built
from the demonstratives plus \tdim{bangin}.

\begin{exe}
\ex \label{ex:dem-hih-bangin}
\gll \tdimword{àmaute} \tdimword{hìlh-ìn,} \tdimword{nòtè} \tdimword{ìn} \tdimword{kà} \tdimword{topa} \tdimword{Esau} \tdimword{kìang-àh} \tdimword{híh} \tdimword{bàng-ìn} \tdimword{nà} \tdimword{cí} \tdimword{dìng} \tdimword{ùh} \tdimword{hì:} \tdimword{Nà} \tdimword{nàsèm-pà} \tdimword{Jakob} \tdimword{ìn,} \tdimword{laban'} \tdimword{kìang-àh} \tdimword{peem-tà-ìn} \tdimword{tù} \tdimword{cìang} \tdimword{dòng} \tdimword{kà} \tdimword{ōm} \tdimword{hì.} \\
     \textsc{3pl}.\textsc{pro} teach-\textsc{cvb} \textsc{2pl}.\textsc{pro} \textsc{erg} \textsc{1sg} lord \textsc{esau} beside-\textsc{loc} \textsc{prox} like-\textsc{erg} \textsc{2sg} say \textsc{irr} 2/3PL \textsc{decl} \textsc{2sg} servant-male \textsc{jakob} \textsc{erg} \textsc{laban}.\textsc{poss} beside-\textsc{loc} widow.son-\textsc{pfv}-\textsc{erg} now then until \textsc{1sg} exist \textsc{decl} \\
\glt 'And he commanded them, saying, Thus shall ye speak unto my lord Esau; Thy servant Jacob saith thus, I have sojourned with Laban, and stayed there until now:' (Genesis 32:4)
\end{exe}

\begin{exe}
\ex \label{ex:dem-tua-bangin}
\gll \tdimword{Topa} \tdimword{ìn} \tdimword{túa} \tdimword{bàng-ìn} \tdimword{Egypt} \tdimword{mite'} \tdimword{khutsung} \tdimword{pàn-ìn} \tdimword{túa} \tdimword{nì-ìn} \tdimword{Israel-té} \tdimword{honkhia} \tdimword{à,} \tdimword{tùi-pì} \tdimword{gei-àh} \tdimword{Egypt} \tdimword{mì} \tdimword{à} \tdimword{sí-tè,} \tdimword{Israel} \tdimword{mì-tè} \tdimword{ìn} \tdimword{mú} \tdimword{ùh} \tdimword{hì.} \\
     Lord \textsc{erg} \textsc{dist} like-\textsc{erg} \textsc{egypt} person-\textsc{pl}.\textsc{poss} palm \textsc{abl}-\textsc{erg} \textsc{dist} day-\textsc{erg} \textsc{israel}-\textsc{te} deliver \textsc{3sg} water-big edge-\textsc{loc} \textsc{egypt} person \textsc{3sg} die-\textsc{pl} \textsc{israel} person-\textsc{pl} \textsc{erg} see.\textsc{i} 2/3PL \textsc{decl} \\
\glt 'Thus the LORD saved Israel that day out of the hand of the Egyptians; and Israel saw the Egyptians dead upon the sea shore.' (Exodus 14:30)
\end{exe}

\tdim{Hih bangin} is especially good for direct instruction or quoted
manner, as also seen in Genesis 42:18 \tdim{hih bangin gamta un}.
\tdim{Tua bangin} often resumes or summarizes an already described
event, as in Genesis 50:21
\tdim{Amah in tua bangin ... amaute a hehnem hi}. These patterns belong
naturally in a demonstratives slice because the deictic base still
contributes to how the construction organizes discourse.

\subsubsection{\texorpdfstring{Deferred forms: \tdim{hi} and
\tdim{hih ciangin}}{Deferred forms:  and }}\label{deferred-forms-and}

\tdim{Hi} should be deferred. Exact-token \tdim{hi} is far too entangled
with copular, auxiliary, and sentence-final material to serve as a safe
demonstrative headword in this packet. Genesis 48:18
\tdim{Hi lo hi, pa aw; hih pen a suak masa ahi hi} shows how quickly the
shorter form is absorbed into other clause types. The right place for
\tdim{hi} is later work on copula or sentence-final particles, not this
first demonstratives packet.

\tdim{Hih ciangin} should also be deferred. The current dossier showed
that many apparent hits involve verbal \tdim{hih} ``do/be thus'' rather
than demonstrative \tdim{hih}. Genesis 18:10, the old generated-report
example, does not contain exact demonstrative \tdim{hih ciangin}. For
the present slice, \tdim{tua ciangin} is strong enough to represent the
temporal side of the system on its own.

\subsubsection{Editorial summary}\label{editorial-summary}

The current evidence supports a compact but stable demonstrative system.
\tdim{Hih} is the core proximal demonstrative. \tdim{Tua} is the core
distal/anaphoric demonstrative and should not be reduced to a merely
spatial ``that''. \tdim{Hihte} and \tdim{tuate} are transparent plural
forms, and the constructional extensions \tdim{hih bangin},
\tdim{tua bangin}, \tdim{tua ciangin}, and \tdim{tua ahih ciangin} are
all well supported by manually checked Bible examples.

Just as important, the chapter has to preserve its cautions. \tdim{Hi}
is not yet safe as a demonstrative headword, \tdim{hih ciangin} is not
yet a stable temporal counterpart to \tdim{tua ciangin}, and the
relation between Bible \tdim{tua} and the literature's \tdim{huā} is
plausible but unresolved. With those limits kept explicit,
demonstratives and deixis make a good narrow chapter for the next Tedim
packet.

\subsection{Pronouns / clusivity}\label{pronouns-clusivity}

\emph{Source slice:
\texttt{output/publication\_review/grammar\_pronouns\_print\_slice.md}}

\subsubsection{Scope}\label{scope-1}

This account presents a short print-facing chapter on Tedim personal
pronouns and closely related pronominal marking. It covers independent
personal pronouns, first-person plural forms and clusivity, possessive
prefixes, emphatic forms in \tdim{-mah}, reflexive and reciprocal
marking with \tdim{ki-}, and a cautious first treatment of \tdim{hong-}
and \tdim{kong-}. It does not attempt a full account of demonstratives,
interrogatives, quantifiers, TAM, or the wider verbal agreement system.

\subsubsection{Personal pronouns}\label{personal-pronouns}

Earlier descriptions agree on a basic personal-pronoun system with
distinct first, second, and third persons and with plural forms for at
least the first and second persons
\citetext{\citealp[32-33]{henderson1965}; \citealp[sec.~3.2.1]{zamngaihcing2017}}.
Henderson explicitly lists \tdim{kei}, \tdim{nang}, \tdim{amah},
\texttt{ei/eite}, \texttt{ko/kote}, and \texttt{no/note}, while Zam
Ngaih Cing confirms the broader person-and-number system and the absence
of gender in the third person
\citetext{\citealp[32]{henderson1965}; \citealp[sec.~3.2.1]{zamngaihcing2017}}.
For print purposes, the following paradigm is stable enough to use:

{\def\LTcaptype{none} % do not increment counter
\begin{longtable}[]{@{}lll@{}}
\toprule\noalign{}
Person & Singular & Plural \\
\midrule\noalign{}
\endhead
\bottomrule\noalign{}
\endlastfoot
1 & \tdim{kei} & \tdim{eite}, \tdim{kote} \\
2 & \tdim{nang} & \tdim{note} \\
3 & \tdim{amah} & \tdim{amaute} \\
\end{longtable}
}

\begin{exe}
\ex \label{ex:pro-amah}
\gll \tdimword{Mipa} \tdimword{ìn} \tdimword{à} \tdimword{zi'} \tdimword{mīn} \tdimword{pèn} \tdimword{Eve,} \tdimword{cí} \tdimword{hì.} \tdimword{Bàng} \tdimword{hàng} \tdimword{hìam} \tdimword{cìh} \tdimword{lèh} \tdimword{àmàh} \tdimword{pèn} \tdimword{mì-hìng} \tdimword{khèmpeuh'} \tdimword{nù} \tdimword{àhì} \tdimword{hì.} \\
     man \textsc{erg} \textsc{3sg} wife.\textsc{poss} name \textsc{top} \textsc{eve} say \textsc{decl} like reason \textsc{q} say.\textsc{nom} and \textsc{3sg}.\textsc{pro} \textsc{top} person-kind all.\textsc{poss} female be.\textsc{3sg} \textsc{decl} \\
\glt 'And Adam called his wife's name Eve; because she was the mother of all living.' (Genesis 3:20)
\end{exe}

\begin{exe}
\ex \label{ex:pro-note}
\gll \tdimword{Ēn} \tdimword{ùn,} \tdimword{nòtè} \tdimword{lè} \tdimword{nòtè'} \tdimword{khìt} \tdimword{à} \tdimword{nà} \tdimword{suanlekhak-tè} \tdimword{ùh} \tdimword{lè,} \\
     look \textsc{imp}.\textsc{pl} \textsc{2pl}.\textsc{pro} and \textsc{2pl}.\textsc{pro}.\textsc{poss} \textsc{seq} \textsc{3sg} \textsc{2sg} genealogy-\textsc{pl} 2/3PL and \\
\glt 'And I, behold, I establish my covenant with you, and with your seed after you;' (Genesis 9:9)
\end{exe}

These examples are ordinary independent pronouns rather than bound
agreement markers. \tdim{amah} functions as a free third-person form,
while \tdim{note} shows the equally straightforward second-person
plural. The printed chapter can therefore begin from free pronouns and
only then move to the more tightly bound prefixal system.

\subsubsection{First-person plural forms and
clusivity}\label{first-person-plural-forms-and-clusivity}

The main editorial problem in this slice is not whether Tedim has
clusivity, but how far the current evidence supports global labels for
the two first-person plural series. Henderson clearly distinguishes
\texttt{ei/eite} from \texttt{ko/kote} through pronominal-concord
prefixes, pairing the former with \tdim{i-} and the latter with
\tdim{ka-} \citep[32-33]{henderson1965}. Zam Ngaih Cing likewise treats
clusivity as a real feature of Tedim person marking, even though her
presentation of the plural forms is not identical in surface detail
\citetext{\citealp[sec.~3.2.1]{zamngaihcing2017}; \citealp[sec.~3.2.2]{zamngaihcing2017}}.
The separate clusivity dossier for this slice shows that sampled Bible
dialogue contexts strongly support \texttt{ko/kote} as exclusive, but
that \texttt{ei/eite} has both clear inclusive uses and less
straightforward uses in the current Bible evidence. Comparative Sukte is
useful here mainly as a contrast, since Singh does not describe the same
inclusive/exclusive opposition for Sukte
\citep[sec.~4.6.1]{sukte_grammar}.

\begin{exe}
\ex \label{ex:pro-eite}
\gll \tdimword{Túa} \tdimword{cìang-ìn} \tdimword{Abram} \tdimword{ìn} \tdimword{Lot'} \tdimword{kìang-àh,} \tdimword{nàng} \tdimword{lè} \tdimword{kèi'} \tdimword{kī-kal,} \tdimword{nàng'} \tdimword{gán-cíng-tè} \tdimword{lè} \tdimword{kèi'} \tdimword{gán-cíng-te'} \tdimword{kī-kal-àh} \tdimword{kī-tot-ná} \tdimword{ōm-sàk} \tdimword{kèi} \tdimword{nì.} \tdimword{Bàng} \tdimword{hàng} \tdimword{hìam} \tdimword{cìh} \tdimword{lèh} \tdimword{eite} \tdimword{beh} \tdimword{khàt} \tdimword{ì-hí} \tdimword{hì"} \\
     \textsc{dist} then-\textsc{erg} \textsc{abram} \textsc{erg} \textsc{lot}.\textsc{poss} beside-\textsc{loc} \textsc{2sg}.\textsc{pro} and \textsc{neg}.\textsc{emph}.\textsc{poss} \textsc{refl}-middle \textsc{2sg}.\textsc{pro}.\textsc{poss} cattle-care-\textsc{pl} and \textsc{neg}.\textsc{emph}.\textsc{poss} cattle-care-\textsc{pl}.\textsc{poss} \textsc{refl}-middle-\textsc{loc} \textsc{refl}-cut-\textsc{nmlz} exist-\textsc{caus} \textsc{neg} day like reason \textsc{q} say.\textsc{nom} and \textsc{1pl}.\textsc{pro} tribe one \textsc{1pl}.\textsc{incl}-be \textsc{decl} \\
\glt 'And Abram said unto Lot, Let there be no strife, I pray thee, between me and thee, and between my herdmen and thy herdmen; for we be brethren.' (Genesis 13:8)
\end{exe}

\begin{exe}
\ex \label{ex:pro-kote}
\gll \tdimword{Kòtè} \tdimword{tàwh} \tdimword{kī-ten-ná} \tdimword{hòng} \tdimword{bāwl} \tdimword{ùn.} \tdimword{Nà} \tdimword{tanu-tè} \tdimword{ùh} \tdimword{kòtè'} \tdimword{tùngàh} \tdimword{hòng} \tdimword{pía} \tdimword{ùn-lá,} \tdimword{no} \tdimword{à} \tdimword{dìng-ìn} \tdimword{kà} \tdimword{tanu-tè} \tdimword{ùh} \tdimword{lá} \tdimword{ùn.} \\
     \textsc{1pl}.\textsc{pro} \textsc{com} \textsc{refl}-marry-\textsc{nmlz} \textsc{3→1} make \textsc{imp}.\textsc{pl} \textsc{2sg} daughter-\textsc{pl} 2/3PL \textsc{1pl}.\textsc{pro}.\textsc{poss} on-\textsc{loc} \textsc{3→1} give \textsc{pl}.\textsc{imp}-and young \textsc{3sg} \textsc{irr}-\textsc{erg} \textsc{1sg} daughter-\textsc{pl} 2/3PL take \textsc{imp}.\textsc{pl} \\
\glt 'And make ye marriages with us, and give your daughters unto us, and take our daughters unto you.' (Genesis 34:9)
\end{exe}

Genesis 13:8 shows a genuinely inclusive use of \tdim{eite}, since Abram
explicitly includes Lot in the relevant group. It does not prove that
every \tdim{eite} token is inclusive. Genesis 34:9, by contrast, is a
strong diagnostic for exclusive \tdim{kote}, since Hamor addresses
Jacob's family from a distinct in-group. The present slice therefore
treats \tdim{kote} as exclusive and leaves the exact global status of
\tdim{eite} under review.\footnote{The earlier report-level reversal
  \tdim{kote} inclusive versus \tdim{eite} exclusive has now been
  removed. The safer editorial position is narrower: \texttt{ko/kote}
  can already be treated as exclusive, but \texttt{ei/eite} still shows
  mixed Bible-corpus behavior and should remain under review rather than
  receiving a single global label.}

\subsubsection{Possessive prefixes}\label{possessive-prefixes}

The same person-marking forms that appear in pronominal concord also
appear before nouns as possessive prefixes. Henderson treats \tdim{ka-},
\tdim{na-}, \tdim{a-}, and \tdim{i-} primarily as pronominal concord
prefixes, while Zam Ngaih Cing foregrounds their possessive use in noun
phrases
\citetext{\citealp[32-33]{henderson1965}; \citealp[sec.~3.2.2]{zamngaihcing2017}; \citealp[sec.~3.3.4.1.1]{zamngaihcing2017}}.
For a printed grammar, the safest description is that the forms are
shared across possessive and agreement environments, but are easiest to
present first in their nominal use.

\begin{exe}
\ex \label{ex:poss-na}
\gll \tdimword{Kua'} \tdimword{tá-nù} \tdimword{nà} \tdimword{hì} \tdimword{hìam,} \tdimword{hòng} \tdimword{gēn} \tdimword{ìn.} \tdimword{Nà} \tdimword{pa'} \tdimword{ínn-àh} \tdimword{kòtè'} \tdimword{giah} \tdimword{nàdìng} \tdimword{à} \tdimword{áwng} \tdimword{dìng} \tdimword{hìam?} \tdimword{à} \tdimword{cí} \tdimword{hì.} \\
     who.\textsc{poss} child-female \textsc{2sg} \textsc{decl} \textsc{q} \textsc{3→1} speak \textsc{erg} \textsc{2sg} male.\textsc{poss} house-\textsc{loc} \textsc{1pl}.\textsc{pro}.\textsc{poss} camp \textsc{purp} \textsc{3sg} open \textsc{irr} \textsc{q} \textsc{3sg} say \textsc{decl} \\
\glt 'Whose daughter art thou? tell me, I pray thee: is there room in thy father's house for us to lodge in?' (Genesis 24:23)
\end{exe}

\begin{exe}
\ex \label{ex:poss-a}
\gll \tdimword{Mipa} \tdimword{ìn} \tdimword{à} \tdimword{zi'} \tdimword{mīn} \tdimword{pèn} \tdimword{Eve,} \tdimword{cí} \tdimword{hì.} \\
     man \textsc{erg} \textsc{3sg} wife.\textsc{poss} name \textsc{top} \textsc{eve} say \textsc{decl} \\
\glt 'And the man called his wife's name Eve.' (Genesis 3:20)
\end{exe}

The same pattern is visible in \tdim{ka pa' inn} `my father's house' and
in the \tdim{i-} prefix that Henderson associates with the
\texttt{ei/eite} series \citep[32-33]{henderson1965}. Singh's Sukte
comparison is also helpful here, since it shows cognate person-marking
prefixes in possessive use even though the wider systems are not
identical \citep[sec.~4.5.4]{sukte_grammar}. Plural verbal markers such
as \tdim{-uh} belong to the later agreement chapter rather than to the
pronoun inventory itself.

\subsubsection{\texorpdfstring{Emphatic pronouns in
\tdim{-mah}}{Emphatic pronouns in }}\label{emphatic-pronouns-in}

Emphatic pronouns are formed by adding \tdim{-mah} to a personal-pronoun
base
\citetext{\citealp[32]{henderson1965}; \citealp[sec.~3.2.6]{zamngaihcing2017}}.
In print, it is best to treat these forms neither as an unrelated
lexical series nor as mere stylistic variants, but as a productive
emphatic pattern built on the ordinary pronoun paradigm.

\begin{exe}
\ex \label{ex:emph-keimah}
\gll \tdimword{Kain} \tdimword{ìn} \tdimword{Topa'} \tdimword{tùngàh,} \tdimword{Kèimah} \tdimword{gìm} \tdimword{hòng} \tdimword{kī-pìak-ná,} \tdimword{kà} \tdimword{thùak} \tdimword{zàwh} \tdimword{dìng} \tdimword{hì} \tdimword{lò} \tdimword{hì.} \\
     \textsc{kain} \textsc{erg} Lord.\textsc{poss} on-\textsc{loc} \textsc{1sg}.\textsc{pro} suffering \textsc{3→1} \textsc{refl}-give.to-\textsc{nmlz} \textsc{1sg} suffer able \textsc{irr} \textsc{decl} \textsc{neg} \textsc{decl} \\
\glt 'And Cain said unto the LORD, My punishment is greater than I can bear.' (Genesis 4:13)
\end{exe}

\begin{exe}
\ex \label{ex:emph-nangmah}
\gll \tdimword{Abimelek} \tdimword{ìn,} \tdimword{Ēn} \tdimword{ìn,} \tdimword{kà} \tdimword{léi-tāng,} \tdimword{nà} \tdimword{mái-àh} \tdimword{ōm} \tdimword{hì.} \tdimword{Nangmah} \tdimword{ìn} \tdimword{hòih} \tdimword{nà} \tdimword{sàk-ná} \tdimword{mùn-àh} \tdimword{tèng} \tdimword{ìn,} \tdimword{à} \tdimword{cí} \tdimword{hì.} \\
     \textsc{abimelek} \textsc{erg} look \textsc{erg} \textsc{1sg} land-earth \textsc{2sg} face-\textsc{loc} exist \textsc{decl} \textsc{2sg}.\textsc{pro} \textsc{erg} good \textsc{2sg} cause-\textsc{nmlz} place-\textsc{loc} remain \textsc{erg} \textsc{3sg} say \textsc{decl} \\
\glt 'And Abimelech said, Behold, my land is before thee: dwell where it pleaseth thee.' (Genesis 20:15)
\end{exe}

The same element also appears in negative-indefinite forms such as
\tdim{kuamah} `no one' and \tdim{bangmah} `nothing'
\citep[sec.~3.2.7]{zamngaihcing2017}. Those items belong more naturally
to a later section on indefinites and interrogative-based forms, but the
connection should already be noted here.

\subsubsection{\texorpdfstring{Reflexive and reciprocal marking with
\tdim{ki-}}{Reflexive and reciprocal marking with }}\label{reflexive-and-reciprocal-marking-with}

Zam Ngaih Cing describes reflexive pronouns as repeated pronouns linked
by \tdim{leh} \citep[sec.~3.2.5]{zamngaihcing2017}. The biblical corpus
nevertheless makes it clear that a print grammar also needs to
acknowledge verbal \tdim{ki-}, since many reader-facing examples of
reflexive or reciprocal meaning are expressed through \tdim{ki-} forms
rather than through a free reflexive pronoun alone.

\begin{exe}
\ex \label{ex:refl-ki}
\gll \tdimword{Túa} \tdimword{thù} \tdimword{hàng-ìn} \tdimword{pasal} \tdimword{ìn} \tdimword{à} \tdimword{nù} \tdimword{lè} \tdimword{à} \tdimword{pà} \tdimword{nusia-ìn} \tdimword{à} \tdimword{zi} \tdimword{tàwh} \tdimword{kī-gawm} \tdimword{à,} \tdimword{àmau} \tdimword{tegel} \tdimword{pum} \tdimword{khàt} \tdimword{à} \tdimword{bàng} \tdimword{ùh} \tdimword{hì.} \\
     \textsc{dist} word because-\textsc{erg} husband \textsc{erg} \textsc{3sg} female and \textsc{3sg} father abandon-\textsc{erg} \textsc{3sg} wife \textsc{com} \textsc{refl}-seize \textsc{3sg} \textsc{3pl}.\textsc{pro} both body one \textsc{3sg} like 2/3PL \textsc{decl} \\
\glt 'Therefore shall a man leave his father and his mother, and shall cleave unto his wife: and they shall be one flesh.' (Genesis 2:24)
\end{exe}

For print purposes, \tdim{ki-} can already be described as a productive
reflexive or reciprocal verbal prefix. The remaining caution is lexical
rather than structural: some surface \tdim{ki-} forms are lexicalized
stems, so not every \tdim{ki-} word should be taken automatically as
transparent reflexive morphology.

\subsubsection{Pronominal prefixes and inverse/directional
marking}\label{pronominal-prefixes-and-inversedirectional-marking}

Henderson's discussion of pronominal concord is still the best starting
point for the bound prefix system. She pairs \tdim{kei} with \tdim{ka-},
\tdim{nang} with \tdim{na-}, the \texttt{ei/eite} series with \tdim{i-},
the \texttt{ko/kote} series with \tdim{ka-}, \texttt{no/note} with
\tdim{na-}, and all other nominals with \tdim{a-}
\citep[32-33]{henderson1965}. Zam Ngaih Cing's discussion of nominal
prefixes confirms the same basic person-marking inventory in noun
phrases \citep[sec.~3.3.4.1.1]{zamngaihcing2017}. A printed chapter on
pronouns does not need the whole verbal paradigm yet, but it does need
to note that the independent pronouns belong to a wider system of bound
person marking.

The harder question concerns \tdim{hong-} and \tdim{kong-}. The
literature treats them as participant-oriented preverbal prefixes with
directional or inverse-like behavior, and Otsuka's causative discussion
shows that first- and second-person objects are structurally important
to their distribution
\citetext{\citealp[sec.~5.8.1.3]{zamngaihcing2017}; \citealp{otsuka_causative}}.
The current corpus outputs, however, do not yet yield a stable automatic
example set for \tdim{hong-}, so the present slice uses only narrow,
manually checked illustrations.\footnote{The present editorial problem
  is not whether these forms exist, but how narrowly they should be
  defined in a first print chapter. The current wording stays close to
  cases that can be read directly from the verse context without relying
  on a fully solved backend analysis.}

\begin{exe}
\ex \label{ex:hong-prefix}
\gll \tdimword{Àmàh} \tdimword{ìn} \tdimword{à} \tdimword{sàng-ná} \tdimword{pàn-ìn} \tdimword{à} \tdimword{khùt} \tdimword{tàwh} \tdimword{zā-mīn} \tdimword{kèi} \tdimword{hòng} \tdimword{lá} \tdimword{à,} \tdimword{tùi} \tdimword{thuk-pì} \tdimword{pàn-ìn} \tdimword{kèimah} \tdimword{hòng} \tdimword{kà-ì-khìa} \tdimword{hì.} \\
     \textsc{3sg}.\textsc{pro} \textsc{erg} \textsc{3sg} high-\textsc{nmlz} \textsc{abl}-\textsc{erg} \textsc{3sg} hand \textsc{com} hundred-thousand \textsc{neg} \textsc{3→1} take \textsc{3sg} water deep-big \textsc{abl}-\textsc{erg} \textsc{1sg}.\textsc{pro} \textsc{3→1} \textsc{1sg}-go-exit \textsc{decl} \\
\glt 'He sent from above, he took me, He drew me out of many waters.' (Psalms 18:16)
\end{exe}

\begin{exe}
\ex \label{ex:kong-prefix}
\gll \tdimword{Faro} \tdimword{ìn} \tdimword{Josef'} \tdimword{kìang-àh,} \tdimword{Egypt} \tdimword{gàm} \tdimword{khèmpeuh} \tdimword{à} \tdimword{ùk} \tdimword{dìng-ìn} \tdimword{nàng} \tdimword{kòng} \tdimword{kòih} \tdimword{khin-zò} \tdimword{hì,} \tdimword{à} \tdimword{cí} \tdimword{hì.} \\
     \textsc{faro} \textsc{erg} \textsc{josef}.\textsc{poss} beside-\textsc{loc} \textsc{egypt} land all \textsc{3sg} rule \textsc{irr}-\textsc{erg} \textsc{2sg}.\textsc{pro} \textsc{1sg→3} put \textsc{compl}-\textsc{compl} \textsc{decl} \textsc{3sg} say \textsc{decl} \\
\glt 'And Pharaoh said unto Joseph, See, I have set thee over all the land of Egypt.' (Genesis 41:41)
\end{exe}

These examples are enough to justify a cautious print description:
\tdim{hong-} and \tdim{kong-} belong with person-sensitive preverbal
marking, but the exact boundary between inverse, venitive, benefactive,
and wider directional readings still needs fuller chapter-level review.
They should therefore appear in the dictionary slice and grammar prose,
but under an explicitly provisional analysis rather than as a closed
paradigm.

\subsubsection{Editorial summary}\label{editorial-summary-1}

This slice now supports a real draft chapter on pronouns and pronominal
marking. Independent pronouns, emphatic forms in \tdim{-mah}, and
possessive prefixes are straightforward enough for print. \tdim{ki-} is
also clear enough to present as reflexive or reciprocal verbal marking,
provided lexicalized forms remain a review checkpoint. The main
unresolved issue is not the existence of clusivity or of
\texttt{hong-/kong-}, but how far the current report layer can be
trusted to label those patterns automatically without closer editorial
control.

\subsubsection{References}\label{references}

\subsection{NP structure / possession}\label{np-structure-possession}

\emph{Source slice:
\texttt{output/publication\_review/grammar\_np\_possession\_print\_slice.md}}

\subsubsection{Editorial scope}\label{editorial-scope}

This is the first narrow NP structure / possession grammar slice. It is
controlled by
\texttt{output/publication\_review/candidates\_np\_possession.tsv} and
\texttt{output/publication\_review/dossier\_np\_possession\_scope.md}.
Supporting/background evidence comes from
\texttt{docs/grammar/reports/03-noun-06-np-structure.md},
\texttt{docs/grammar/reports/04-np-07-possession.md},
\texttt{docs/grammar/lit-reviews/04-np-07-possession-lit.md}, and
\texttt{docs/grammar/morphemes/01-prefixes.md}.

This is not a full noun-phrase chapter, not a full possession chapter,
not a full prefix/agreement chapter, and not a full case or relator
chapter. It also stays narrow against
\texttt{output/publication\_review/review\_notes\_prefix\_agreement.md},
\texttt{output/publication\_review/review\_notes\_pronouns.md},
\texttt{output/publication\_review/review\_notes\_case\_marking.md},
\texttt{output/publication\_review/review\_notes\_relators\_postpositions.md},
\texttt{output/publication\_review/review\_notes\_nominalization.md},
and \texttt{tests/test\_prefix\_agr\_poss.py}.

The present slice therefore covers only the first safe NP-ordering
claim. No dictionary slice exists yet for NP structure / possession,
because this packet is still establishing a controlled
structural/syntactic claim rather than a lexical layer. The packet now
properly proceeds through NP structure / possession review notes rather
than a dictionary slice.

\subsubsection{Basic NP ordering}\label{basic-np-ordering}

The safest current claim in the packet is a small NP-ordering
observation rather than a possession claim.

\tdim{hih mite} is the demonstrative-before-noun anchor. The controlled
segmentation is \tdim{hih mi-te}, and the controlled gloss is
\tdim{PROX person-PL}.

\tdim{mi khat} is the head-noun plus numeral anchor. The controlled
gloss is \tdim{person one}.

\tdim{mi khempeuh} is the head-noun plus quantifier anchor. The
controlled segmentation is \tdim{mi khem-peuh}, and the controlled gloss
is \tdim{person all}.

Taken together, these rows support a conservative NP-ordering statement:
demonstratives can precede the head noun, while numerals and
quantifier-like modifiers can follow the head noun. That is enough for a
first print-facing claim, but it is not yet a full theory of all NP
modifiers, adjectival ordering, recursive noun-phrase structure, or case
closure.

\subsubsection{Why possession is not yet the first
slice}\label{why-possession-is-not-yet-the-first-slice}

Possession and possessor-possessed structure are also visible in the
candidate layer through \tdim{ka pa}, \tdim{Topa' inn}, and
\tdim{a pa' inn}.

Those rows stay outside the first print-facing claim because they still
interact with prefix/agreement routing, pronouns, apostrophe or genitive
analysis, and broader possessor-possessed theory. \tdim{ka pa} is the
safest possession row if a later possession sub-scope is chosen, but it
is still more boundary-heavy than the three clean NP-order anchors.

\subsubsection{Boundary material}\label{boundary-material}

The rest of the NP structure / possession packet stays outside the first
grammar slice because each row is still dominated by another unresolved
boundary.

\tdim{ka pa}, \tdim{Topa' inn}, and \tdim{a pa' inn} stay outside
because possession and possessor-possessed structure still interact with
prefix/agreement routing, pronouns, apostrophe or genitive analysis, and
broader possession theory.

\tdim{Topa' tungah} stays outside because possessive NP plus relator or
case material remains shared with
\texttt{review\_notes\_case\_marking.md} and
\texttt{review\_notes\_relators\_postpositions.md}.

\tdim{ka suahna leitang} stays outside because nominalized noun-headed
material remains shared with \texttt{review\_notes\_nominalization.md}.

isolated \tdim{a}, \tdim{ka}, or \tdim{na} prefix surfaces stay outside
because they are analyzer-noisy away from a controlled nominal host.

Pronoun-led possessor rows such as \tdim{amah a pa} stay outside because
they still sit between NP structure, possession, and the completed
pronoun packet.

Tone-marked or literature-only genitive claims such as \tdim{-á} stay
outside because they are not yet tied tightly enough to corpus-backed
first-slice anchors.

report-only counts stay outside because attestation alone does not make
a row safe for the first print-facing claim.

Any broad recursive possession chapter claim stays outside because this
packet is not yet ready to generalize from one safe NP-ordering slice to
the whole noun-domain architecture.

\subsubsection{Safe first-slice claim}\label{safe-first-slice-claim}

At the current slice maturity level, the safest NP structure /
possession claim is that Tedim has candidate-controlled evidence for
basic NP ordering: demonstratives can precede the noun, while numerals
and quantifier-like modifiers can follow the noun. Possession and
possessor-possessed structures remain candidate-layer or boundary
material.

That claim is deliberately smaller than a full noun-phrase chapter,
smaller than a full possession chapter, smaller than a full
prefix/agreement chapter, and smaller than a full case or relator
chapter.

\subsubsection{Recommended next step}\label{recommended-next-step}

This packet now properly proceeds through NP structure / possession
review notes rather than a dictionary slice, because this packet is
structural/syntactic rather than lexical.

If later work returns to possession after review notes, it should be a
separate narrow possession sub-scope led by \tdim{ka pa}, not a broad
possession chapter.

\subsection{Noun domain}\label{noun-domain}

\emph{Source slice:
\texttt{output/publication\_review/grammar\_noun\_domain\_print\_slice.md}}

\subsubsection{Editorial scope}\label{editorial-scope-1}

This is the first narrow noun-domain grammar slice. It is controlled by
\texttt{output/publication\_review/candidates\_noun\_domain.tsv} and
\texttt{output/publication\_review/dossier\_noun\_domain\_scope.md}.
Supporting/background evidence comes from
\texttt{docs/grammar/reports/03-noun-01-simple.md},
\texttt{docs/grammar/reports/03-noun-02-compounds.md},
\texttt{docs/grammar/reports/03-noun-03-proper.md},
\texttt{docs/grammar/compound\_transparency\_audit.md}, and
\texttt{docs/grammar/opaque\_lexemes.md}.

This is not a full noun chapter, not a compound-noun chapter, not a
proper-noun chapter, and not a dictionary slice. It also stays narrow
against
\texttt{output/publication\_review/review\_notes\_np\_possession.md},
\texttt{output/publication\_review/review\_notes\_nominalization.md},
\texttt{output/publication\_review/review\_notes\_relators\_postpositions.md},
\texttt{output/publication\_review/review\_notes\_case\_marking.md}, and
\texttt{output/publication\_review/review\_notes\_pronouns.md}.

The present slice therefore covers only the first safe simple-noun-stem
claim. No dictionary slice exists for the noun domain, because this
packet is still establishing a grammar-facing noun-domain foundation
rather than a lexical layer.

\subsubsection{Simple noun stems}\label{simple-noun-stems}

\tdim{gam} is the main simple free noun stem anchor.

The controlled report-backed forms are \tdim{gam}, \tdim{gam-te},
\tdim{gam-'}, \tdim{gam-in}, \tdim{gam-ah}, and \tdim{gam-te-ah}.

Taken together, these forms show a free noun stem that can host ordinary
plural and case-like marking. That is enough for a first print-facing
claim, but it is not a full noun inflection chapter, not a full case
chapter, and not a full noun chapter.

\tdim{aksi / aksi-te} is the supporting plural row. It shows ordinary
plural marking on a second simple noun without broadening the slice into
a larger noun-domain inventory.

\subsubsection{Why compounds are not yet the first
slice}\label{why-compounds-are-not-yet-the-first-slice}

\tdim{minam} and \tdim{thugen} are the candidate-layer transparent
compounds in the current packet.

They stay outside the first print-facing claim because compounds still
require transparency and lexicalization decisions before they can safely
lead grammar prose.

\tdim{singnai} and \tdim{sanggam} are the main boundary rows here.
\tdim{sanggam} stays outside because it is opaque or lexicalized, while
\tdim{singnai} stays outside because it is a transparency-problem row
rather than a clean transparent-compound anchor.

\subsubsection{Why proper nouns are not yet the first
slice}\label{why-proper-nouns-are-not-yet-the-first-slice}

\tdim{Abraham} is the clean proper-noun candidate, and \tdim{Topa} is
title-like boundary material.

Proper nouns stay outside the first print-facing claim because they are
more lexical-inventory-like than the simple noun anchors, and
\tdim{Topa} also remains entangled with noun-domain and NP-possession
boundary questions.

\subsubsection{Boundary material}\label{boundary-material-1}

\tdim{minam} and \tdim{thugen} stay outside the first grammar slice
because transparent compounds still require lexicalization and
transparency decisions before they can lead noun-domain grammar prose.

\tdim{singnai} and \tdim{lamethuai} stay outside because they remain
transparency-problem rows rather than first-slice anchors.

\tdim{sanggam} stays outside because it is opaque or lexicalized rather
than a safe transparent-compound example.

\tdim{kholhna} stays outside because it is nominalization-boundary
material rather than a clean simple noun stem.

\tdim{Abraham} and \tdim{Topa} stay outside because proper nouns remain
more lexical-inventory-like than the simple noun anchors, while
\tdim{Topa} is also title-like boundary material.

\tdim{Topa' inn} or broader possessor syntax stay outside because they
belong with NP structure / possession rather than the first noun-domain
slice.

Pronoun-led possessors or person-head material stay outside because they
still sit between noun-domain work, pronouns, and NP structure.

relator/postposition or case-dominated noun rows stay outside because
their main value belongs with
\texttt{review\_notes\_relators\_postpositions.md} and
\texttt{review\_notes\_case\_marking.md}.

analyzer-noisy, report-only, or count-only noun-domain claims stay
outside because they do not yet produce safe first-slice anchors.

Any broad noun chapter claim stays outside because this packet is not
yet ready to generalize from one safe simple-noun-stem slice to the
whole noun domain.

\subsubsection{Safe first-slice claim}\label{safe-first-slice-claim-1}

At the current slice maturity level, the safest noun-domain claim is
that Tedim has candidate-controlled evidence for simple free noun stems
that can host ordinary plural and case-like marking, with \tdim{gam} as
the main anchor and \tdim{aksi / aksi-te} as supporting plural evidence.
Compound nouns and proper nouns remain candidate-layer or boundary
material.

That claim is deliberately smaller than a full noun chapter, smaller
than a compound-noun chapter, smaller than a proper-noun chapter, and
smaller than a dictionary slice.

\subsubsection{Recommended next step}\label{recommended-next-step-1}

This grammar slice is now paired with
\texttt{output/publication\_review/review\_notes\_noun\_domain.md}, so
the packet is ready for human review at its current simple-noun-stem
slice maturity level.

If more noun-domain work is chosen after review notes, the next
sub-scope should be transparent compounds led by \tdim{minam} and
\tdim{thugen}, not opaque compounds or proper nouns.

\subsection{Case marking}\label{case-marking}

\emph{Source slice:
\texttt{output/publication\_review/grammar\_case\_marking\_print\_slice.md}}

\subsubsection{Scope}\label{scope-2}

This review slice presents a short draft chapter on nominal case marking
and closely related postpositional constructions in Tedim Chin. It is
intended as an editorial model for later sections of the grammar, not as
a full treatment of nominal morphology. The discussion concentrates on
ergative \tdim{-in}, locative \tdim{-ah}, ablative and source forms
\tdim{-pan} and \tdim{-panin}, and comitative \tdim{-tawh}. A final
section notes the role of relator nouns, since the clearest spatial
examples regularly involve stems such as \tdim{lak}, \tdim{sung},
\tdim{kiang}, and \tdim{laizang}.

The present print-facing examples are controlled by
\texttt{candidates\_case\_marking.tsv} and interpreted in
\texttt{dossier\_case\_marking.md}. This keeps the slice narrow: the
packet now prints candidate-backed anchors and explicit caveats rather
than broadening into a new full chapter or a broad automatic case-marker
survey.

\subsubsection{Case marking in outline}\label{case-marking-in-outline}

Earlier descriptions agree that Tedim marks nominal relations after the
noun phrase, but they differ in how they analyze those forms. Henderson
describes them structurally as phrase-final or post-nominal particles
rather than as a semantic case system
\citetext{\citealp[59]{henderson1965}; \citealp[104]{henderson1965}}.
Zam Ngaih Cing, by contrast, offers a seven-case analysis with explicit
ergative, locative, ablative, and comitative categories
\citep[sec.~3.3.3.3]{zamngaihcing2017}. Comparative Sukte evidence
points in the same direction for several core markers, even though Singh
labels the cognate \tdim{-in} as nominative rather than ergative
\citep[sec.~4.5.1]{sukte_grammar}.

For a printed grammar, the most useful generalization is therefore
twofold. First, the language is well described by a case-marking
analysis in the modern sense. Second, the most natural corpus examples
often combine case markers with relational nominal stems, so a final
chapter should not force an artificial boundary between ``case
suffixes'' and ``spatial nouns''.

\subsubsection{\texorpdfstring{Ergative
\tdim{-in}}{Ergative }}\label{ergative}

The ergative marker \tdim{-in} marks the transitive subject in Tedim
Chin \citep[sec.~3.3.3.3.1]{zamngaihcing2017}. Henderson does not use
the label \emph{ergative}, but her treatment of phrase-final nominal
particles is compatible with the later analysis
\citep[59]{henderson1965}. Otsuka likewise assumes ergative \tdim{-in}
in discussing causers in causative constructions
\citep{otsuka_causative}. Comparative Sukte retains a cognate \tdim{-in}
marker, though Singh describes it as nominative
\citep[sec.~4.5.1]{sukte_grammar}.

In this slice, Genesis 4:3 remains the accepted print anchor because
\texttt{candidates\_case\_marking.tsv} and
\texttt{dossier\_case\_marking.md} both treat \tdim{Kain in} as the
cleanest candidate-backed ergative window. Raw \tdim{-in} extraction is
unsafe: rows such as \tdim{ciangin} are conjunctional or other non-case
material, so they must not be promoted as ergative case examples.

\begin{exe}
\ex \label{ex:erg-in}
\gll \tdimword{Túa} \tdimword{hùn} \tdimword{sùng-ìn} \tdimword{Kain} \tdimword{ìn} \tdimword{léi} \tdimword{pà-ná} \tdimword{piang} \tdimword{gah} \tdimword{pìak-ná} \tdimword{Topa'} \tdimword{tùngàh} \tdimword{pái-pìh} \tdimword{à,} \\
     \textsc{dist} time inside-\textsc{erg} \textsc{kain} \textsc{erg} buy father-\textsc{nmlz} be.born branch give.to-\textsc{nmlz} Lord.\textsc{poss} on-\textsc{loc} go-\textsc{appl} \textsc{3sg} \\
\glt 'And in process of time it came to pass, that Cain brought of the fruit of the ground an offering unto the LORD.' (Genesis 4:3)
\end{exe}

This example is sufficient for a draft chapter because it shows the
essential pattern without relying on doubtful segmentation. The agent
phrase \tdim{Kain in} is simple, animate, and immediately followed by a
clearly transitive predicate. At the same time, the section remains
cautious: one accepted candidate-backed example is enough to print the
analysis, but additional \tdim{-in} rows should still be filtered
against ambiguity controls such as \tdim{ciangin} before they are
treated as nominal case evidence.

\subsubsection{\texorpdfstring{Locative
\tdim{-ah}}{Locative }}\label{locative}

The locative \tdim{-ah} marks location
\citep[sec.~3.3.3.3.3]{zamngaihcing2017}. Henderson already treats
\tdim{-ah} as a locative post-nominal particle and discusses its
phonological behavior in locative nominal figures
\citep[54-56]{henderson1965}. The current candidate layer, however,
distinguishes plain noun-plus-locative evidence such as \tdim{khua-ah}
from relator-noun-plus-case constructions such as \tdim{laizangah},
\tdim{vantungah}, \tdim{kiangah}, \tdim{sungah}, and \tdim{tungah}. The
printed description should therefore be slightly richer than a simple
gloss ``at/in'', while still keeping those two subtypes distinct.

\begin{exe}
\ex \label{ex:loc-ah}
\gll \tdimword{Pasian} \tdimword{ìn,} \tdimword{tuite'} \tdimword{lai-zang-àh} \tdimword{vàn} \tdimword{kuum-pì} \tdimword{ōm} \tdimword{hèn} \tdimword{lá,} \tdimword{túa} \tdimword{vàn} \tdimword{kuum-pì} \tdimword{ìn} \tdimword{tùi} \tdimword{lè} \tdimword{tùi} \tdimword{kī-khen-sàk} \tdimword{hèn} \tdimword{cí} \tdimword{hì.} \\
     God \textsc{erg} water-\textsc{pl}.\textsc{poss} middle-side-\textsc{loc} sky bow-big exist \textsc{juss} take \textsc{dist} sky bow-big \textsc{erg} water and water \textsc{refl}-divide-\textsc{caus} \textsc{juss} say \textsc{decl} \\
\glt 'And God said, Let there be a firmament in the midst of the waters, and let it divide the waters from the waters.' (Genesis 1:6)
\end{exe}

\begin{exe}
\ex \label{ex:loc-ah-2}
\gll \tdimword{léi-tùng} \tdimword{khùa} \tdimword{à} \tdimword{vak-sàk} \tdimword{dìng-ìn} \tdimword{vantung-àh} \tdimword{khuavak} \tdimword{hì} \tdimword{ùh} \tdimword{hèn} \tdimword{cí} \tdimword{hì.} \\
     land-on town \textsc{3sg} shine-\textsc{caus} \textsc{irr}-\textsc{erg} heaven-\textsc{loc} light \textsc{decl} 2/3PL \textsc{juss} say \textsc{decl} \\
\glt 'and let them be for lights in the firmament of the heaven to give light upon the earth' (Genesis 1:15)
\end{exe}

The first example shows especially clearly that \tdim{-ah} often
completes a larger relational expression rather than attaching only to a
simple lexical noun. The noun \tdim{laizang} contributes the spatial
geometry `middle, interior region', and \tdim{-ah} marks that whole
phrase as locative. The second printed example is still best treated in
the same relator-noun-plus-case domain, even if \tdim{vantung} looks
less obviously relational than \tdim{laizang}. The plain
noun-plus-locative control in the candidate layer is \tdim{khua-ah}, not
one of these relator rows.

The analyzer export also needs to be read cautiously here. Some locative
and relator rows surface with \tdim{pos\_span=FUNC}, even though the
grammar still treats the base as nominal or relational. That export
label is useful metadata, but it is not decisive enough to collapse
noun-like and relator-like rows into a single undifferentiated suffix
list.

\subsubsection{\texorpdfstring{Directional/allative
\tdim{-a}}{Directional/allative }}\label{directionalallative}

Directional or allative \tdim{-a} remains deferred in the current
packet. The candidate layer keeps \tdim{-a} separate from \tdim{-ah},
and the dossier is explicit that the present export does not yet
distinguish directional \tdim{-a} cleanly from pronominal or other
functional \tdim{a} tokens. For that reason, the current slice does
\textbf{not} collapse \tdim{-a} into \tdim{-ah}, and it does not yet
print a dedicated \tdim{-a} example.

\subsubsection{\texorpdfstring{Source marking: \tdim{-pan} and
\tdim{-panin}}{Source marking:  and }}\label{source-marking-and}

Zam Ngaih Cing describes \tdim{-pan} as the ablative marker indicating
source or point of departure \citep[sec.~3.3.3.3.5]{zamngaihcing2017}.
Comparative Sukte has a cognate \tdim{-pan}, which strengthens the
historical plausibility of the analysis \citep{sukte_grammar}. Henderson
does not present an ablative category in the modern sense, but the
source-marking behavior is nevertheless compatible with her structural
account of post-nominal particles \citep[104]{henderson1965}.

\begin{exe}
\ex \label{ex:pan}
\gll \tdimword{Túa} \tdimword{àhìh} \tdimword{màn-ìn} \tdimword{híh} \tdimword{thù-khàm-te} \tdimword{lak-pàn} \tdimword{à} \tdimword{néu-pèn-tè} \tdimword{khàt} \tdimword{bèk} \tdimword{nà-ngawn} \tdimword{zúi-khá} \tdimword{lò-ìn,} \tdimword{midang-tè} \tdimword{ìn} \tdimword{zòng} \tdimword{à} \tdimword{zùih} \tdimword{lòh} \tdimword{nàdìngìn} \tdimword{à} \tdimword{gēn} \tdimword{mì-tè} \tdimword{pèn} \tdimword{vàn-tùng} \tdimword{kī-ukna} \tdimword{sùng-àh} \tdimword{mì} \tdimword{néu-pèn} \tdimword{hì} \tdimword{dìng} \tdimword{ùh} \tdimword{hì.} \\
     \textsc{dist} be.\textsc{3sg}.\textsc{rel} reason-\textsc{erg} \textsc{prox} law-\textsc{pl} midst-\textsc{abl} \textsc{3sg} small-most-\textsc{pl} one only \textsc{2sg}-own follow-go \textsc{neg}-\textsc{erg} other.person-\textsc{pl} \textsc{erg} also \textsc{3sg} follow-\textsc{nom} \textsc{neg}-\textsc{nom} \textsc{purp}.\textsc{erg} \textsc{3sg} speak person-\textsc{pl} \textsc{top} sky-on \textsc{refl}-rule.\textsc{nmlz} inside-\textsc{loc} person small-\textsc{top} \textsc{decl} \textsc{irr} 2/3PL \textsc{decl} \\
\glt 'Whosoever therefore shall break one of these least commandments, and shall teach men so, he shall be called the least in the kingdom of heaven.' (Matthew 5:19)
\end{exe}

This example is worth keeping because it shows a natural source
construction built on a relator noun. The source relation is real, but
it is not expressed through a bare noun alone: \tdim{lak} `among, in the
midst of' combines with \tdim{-pan} to yield a source phrase meaning
`from among'. The current candidate layer therefore treats \tdim{lakpan}
as source marking on a relator noun, not merely as a bare suffix
example.

The related form \tdim{-panin} is common in the corpus and clearly
belongs to the same source-marking domain, especially after relator
nouns and other spatial expressions. What remains uncertain is not its
source meaning, but its exact structural status in every occurrence.
Some examples look like straightforward source phrases, while others
behave more like extended or tightly fused source-marking expressions.

\begin{exe}
\ex \label{ex:panin}
\gll \tdimword{Topa} \tdimword{ìn} \tdimword{Abram} \tdimword{kìang-àh,} \tdimword{nà} \tdimword{gàm,} \tdimword{nà} \tdimword{beh,} \tdimword{lè} \tdimword{nà} \tdimword{pa'} \tdimword{ínn} \tdimword{pàn-ìn} \tdimword{nàng} \tdimword{kòng} \tdimword{làh} \tdimword{dìng} \tdimword{gàm-àh} \tdimword{pái} \tdimword{ìn} \\
     Lord \textsc{erg} \textsc{abram} beside-\textsc{loc} \textsc{2sg} land \textsc{2sg} tribe and \textsc{2sg} male.\textsc{poss} house \textsc{abl}-\textsc{erg} \textsc{2sg}.\textsc{pro} \textsc{1sg→3} lamp \textsc{irr} land-\textsc{loc} go \textsc{erg} \\
\glt 'Now the LORD had said unto Abram, Get thee out of thy country, and from thy kindred, and from thy father’s house.' (Genesis 12:1)
\end{exe}

For that reason, a printed grammar can already gloss \tdim{-panin} as a
source form related to \tdim{-pan}, but should avoid a stronger
compositional claim until the full range of corpus contexts has been
reviewed more closely. The analyzer segmentation \tdim{pan-in} is useful
evidence for the packet, but it is not by itself a final structural
analysis.

\subsubsection{\texorpdfstring{Comitative
\tdim{-tawh}}{Comitative }}\label{comitative}

The comitative \tdim{-tawh} is described in the literature as the marker
of accompaniment \citep[sec.~3.3.3.3.6]{zamngaihcing2017}. Henderson
also lists \tdim{tawh} among the post-nominal particles
\citep[104]{henderson1965}. In the biblical corpus, however, the most
accessible examples are semantically broader than simple human
accompaniment. A printed chapter should therefore distinguish the core
comitative use from its material and instrumental extensions rather than
treating all examples as equivalent.\footnote{The currently available
  biblical examples for \tdim{-tawh} are semantically good enough for
  review, but they do not all illustrate the most prototypical
  accompaniment reading. That is why the present slice pairs one
  accompaniment example with one extension example rather than treating
  Genesis 2:7 or Genesis 2:21 as if they were ordinary companion
  phrases.}

\begin{exe}
\ex \label{ex:tawh-accompaniment}
\gll \tdimword{Tangvalte'} \tdimword{nèk-sá-tè} \tdimword{lè} \tdimword{kèi} \tdimword{tàwh} \tdimword{hòng} \tdimword{kúan} \tdimword{mite'} \tdimword{tành} \tdimword{dìng-tè} \tdimword{lobuang} \tdimword{bangmah} \tdimword{kà} \tdimword{lá} \tdimword{kèi} \tdimword{dìng} \tdimword{hì.} \\
     youth-\textsc{pl}.\textsc{poss} eat.\textsc{ii}-flesh-\textsc{pl} and \textsc{1sg}.\textsc{pro} \textsc{com} \textsc{3→1} family person-\textsc{pl}.\textsc{poss} spread stand-\textsc{pl} disturb nothing \textsc{1sg} take \textsc{neg} \textsc{irr} \textsc{decl} \\
\glt 'save only that which the young men have eaten, and the portion of the men which went with me' (Genesis 14:24)
\end{exe}

\begin{exe}
\ex \label{ex:tawh-extension}
\gll \tdimword{Topa} \tdimword{Pasian} \tdimword{ìn} \tdimword{léi-lak} \tdimword{pà-ná} \tdimword{leivui} \tdimword{tàwh} \tdimword{mì-hìng} \tdimword{bāwl} \tdimword{à,} \tdimword{à} \tdimword{nak} \tdimword{sùng-àh} \tdimword{nuntak-ná} \tdimword{hu} \tdimword{sàng} \tdimword{sùk} \tdimword{hì.} \\
     Lord God \textsc{erg} dust father-\textsc{nmlz} dust \textsc{com} person-kind make \textsc{3sg} \textsc{3sg} nose inside-\textsc{loc} live-\textsc{nmlz} breath high collapse \textsc{decl} \\
\glt 'And the LORD God formed man of the dust of the ground, and breathed into his nostrils the breath of life.' (Genesis 2:7)
\end{exe}

Genesis 14:24 illustrates the core accompaniment use well enough for
print: \tdim{kei tawh} marks the men as companions who went with Abram.
Genesis 2:7, by contrast, shows a broader use in which \tdim{-tawh}
marks material or means. That wider range should remain explicit in the
grammar and stay split in the same way it is split in the candidate
layer. The comitative meaning remains central, but the marker extends
into associated material and instrument-like readings in ways that the
English gloss \emph{with} only partly captures.

\subsubsection{Relator nouns}\label{relator-nouns}

The locative and source examples already show that Tedim spatial grammar
is not exhausted by a list of case markers. Relator nouns such as
\tdim{lak}, \tdim{sung}, \tdim{kiang}, \tdim{tung}, \tdim{laizang}, and
\tdim{vantung} regularly host locative and ablative marking, and they do
so at very high frequency in the candidate-backed packet. Henderson's
structural treatment makes room for this by analyzing many such forms in
terms of nominal figures rather than sharply separated case suffixes
\citep[59]{henderson1965}. Zam Ngaih Cing's case system, meanwhile,
makes the semantic contribution of the markers clearer
\citep[sec.~3.3.3.3]{zamngaihcing2017}.

For a final chapter, the best solution will probably be to treat the
system in two layers: first the case markers themselves, then the class
of relational nouns that commonly host them. The present review slice
stops short of a full relator-noun section, but it already makes clear
that the two layers belong in the same part of the grammar and should
not be flattened into a bare suffix list.

\subsubsection{Editorial summary}\label{editorial-summary-2}

This slice now supports a modest but genuine draft chapter. Ergative
\tdim{-in} can be described with a manually confirmed biblical example;
\tdim{-ah} and \tdim{-pan} are already straightforward to illustrate;
\tdim{-panin} can be described conservatively as a related source form;
and \tdim{-tawh} can be presented with one clear accompaniment example
plus one extension example. The remaining work is not to broaden the
chapter, but to decide how fully the final grammar should integrate
relator nouns into the same printed discussion.

\subsubsection{References}\label{references-1}

\begin{itemize}
\tightlist
\item
  Henderson, Eugénie J. A. 1965. \emph{Tiddim Chin: A Descriptive
  Analysis of Two Texts}. London: Oxford University Press.
\item
  Otsuka, Kosei. 2009. ``Causative and benefactive suffix -sàk in Tiddim
  Chin.'' \emph{思言: 東京外国語大学記述言語学論集} 5: 3-24.
\item
  Singh, L. S. 2018. \emph{A Descriptive Grammar of Sukte (Salhte)}.
  Imphal: Grassroot Publications.
\item
  Zam Ngaih Cing. 2017. \emph{A Descriptive Grammar of Tedim Chin}. PhD
  dissertation, North-Eastern Hill University, Shillong.
\end{itemize}

\subsection{Relators / postpositions}\label{relators-postpositions}

\emph{Source slice:
\texttt{output/publication\_review/grammar\_relators\_postpositions\_print\_slice.md}}

\subsubsection{Editorial scope}\label{editorial-scope-2}

This is the first narrow grammar slice for Tedim relators /
postpositions. It is controlled by
\texttt{candidates\_relators\_postpositions.tsv} and
\texttt{dossier\_relators\_postpositions\_scope.md}. The generated
reports \texttt{docs/grammar/reports/03-noun-04-relators.md} and
\texttt{docs/grammar/reports/03-noun-05-postpositions.md} remain
discovery or background sources only.

The case-marking packet is the boundary control for this slice:
\texttt{candidates\_case\_marking.tsv},
\texttt{dossier\_case\_marking.md},
\texttt{grammar\_case\_marking\_print\_slice.md},
\texttt{dictionary\_case\_markers\_print\_slice.md}, and
\texttt{review\_notes\_case\_marking.md} remain the main print-facing
treatment of case marking. This slice is therefore not a rewrite of the
case-marking packet. The relators/postpositions dictionary slice now
exists at \texttt{dictionary\_relators\_postpositions\_print\_slice.md},
and review notes now exist at
\texttt{review\_notes\_relators\_postpositions.md}.

\subsubsection{Relator nouns as relational
hosts}\label{relator-nouns-as-relational-hosts}

The first-slice relator-noun anchors are \tdim{kiang}, \tdim{lak},
\tdim{sung}, \tdim{tung}, and cautiously \tdim{pualam}. These are best
treated as relational nouns or relational stems that can host locative
or source marking rather than as simple bare case suffixes.

\tdim{Kiang} is the clearest `beside / near / at the side of' anchor.
\tdim{Lak} is the clearest `among / amid / between' anchor. \tdim{Sung}
is the clearest `inside / within / among' anchor. \tdim{Tung} is usable
for `on / above / upon', but it stays slightly closer to case-marking
territory than the first three anchors and should therefore keep an
explicit caveat. \tdim{Pualam} is usable for `outside / exterior side',
but it should remain the most cautious positive relator-noun anchor in
the first slice because it still looks somewhat more lexical-nominal
than \tdim{kiang}, \tdim{lak}, \tdim{sung}, or \tdim{tung}.

The point of this slice is not to deny the role of \tdim{-ah} or
\tdim{-pan}. It is to keep visible that forms such as \tdim{kiang},
\tdim{lak}, \tdim{sung}, \tdim{tung}, and \tdim{pualam} contribute
relational-host structure before locative or source marking is added.
They are not simply bare case suffixes in disguise.

\subsubsection{Separate and relator-hosted
postpositions}\label{separate-and-relator-hosted-postpositions}

The first-slice postpositional material stays deliberately narrow.
\tdim{Pan} is usable only as a separate or relator-hosted source
postposition, as in \tdim{kiang pan} or \tdim{lak pan}. \tdim{Panin} is
usable only as a source form with structural caution. \tdim{Tawh} is
usable only as a separate accompaniment or associative postposition, and
it must keep its case-marking boundary caveat visible.

This slice therefore treats \tdim{pan}, \tdim{panin}, and \tdim{tawh}
only where they are clearly separate or clearly relator-hosted. It does
\textbf{not} replace the existing case-marking treatment of \tdim{-ah},
\tdim{-in}, \tdim{-pan}, \tdim{-panin}, or \tdim{-tawh}. The goal is to
show how relator nouns and separate postpositions meet in the same
constructions without reopening the settled case-marking packet.

\tdim{Tawhin} remains deferred or boundary-only in the first slice. The
current evidence is sparse and strongly instrument-like, so it is not
yet a clean postpositional anchor parallel to \tdim{pan}, \tdim{panin},
or \tdim{tawh}.

\subsubsection{Case-marking boundary}\label{case-marking-boundary}

Attached or fused-looking rows such as \tdim{kiangah}, \tdim{sungah},
\tdim{tungah}, \tdim{lakpan}, and similar forms are shared boundary
territory with the case-marking packet. They are useful here because
they show the relational-host side of the system, but they are not a
reason to reopen the case-marking analysis or to treat every
attached-looking form as an independent postposition.

This is the main value of the current grammar slice. The case-marking
packet already explains the marker side of the system. The
relators/postpositions slice adds value by explaining the
relational-host side: why \tdim{kiang}, \tdim{lak}, \tdim{sung},
\tdim{tung}, and \tdim{pualam} matter, and why separate or
relator-hosted \tdim{pan}, \tdim{panin}, and \tdim{tawh} should not be
flattened into a broad inventory of suffix-like strings.

\subsubsection{Deferred and boundary
material}\label{deferred-and-boundary-material}

Several items stay outside the first grammar slice as core claims.

\begin{itemize}
\tightlist
\item
  \tdim{nuai} remains lower-confidence boundary material rather than a
  first-slice relator-noun anchor.
\item
  \tdim{mai} remains lower-confidence boundary material because it stays
  close to lexical noun or body-part uses.
\item
  \tdim{tawhin} remains deferred or boundary-only because the current
  evidence is sparse and instrument-like.
\item
  Raw report counts are not evidence for the slice.
\item
  \tdim{kipan} and \tdim{kipanin} rows involving \tdim{ki} remain
  boundary material rather than clean independent postpositions.
\item
  Attached or fused-looking forms should not be promoted here as
  independent postpositions just because they contain \tdim{pan},
  \tdim{panin}, or locative material.
\end{itemize}

\subsubsection{Recommended next step}\label{recommended-next-step-2}

With \texttt{review\_notes\_relators\_postpositions.md} now added, the
packet is ready for human review at the current slice maturity level.
Later work should remain narrow and case-boundary-controlled rather than
broadening into a new case-marking chapter.

\subsection{Numerals}\label{numerals}

\emph{Source slice:
\texttt{output/publication\_review/grammar\_numerals\_print\_slice.md}}

\subsubsection{Scope}\label{scope-3}

This is a short print-facing draft section on numerals in Tedim Chin,
controlled by \texttt{candidates\_numerals.tsv} and
\texttt{dossier\_numerals.md}. It covers only a small candidate-backed
set: basic cardinal counting phrases, compound tens, one ordinal form,
one occurrence-counting or multiplicative form, one large-number phrase
with caveats, and explicit ambiguity controls for \tdim{kua} and
\tdim{khat}.

It does not yet attempt a full numeral paradigm, a full classifier
system, a full account of distributive reduplication, quantifiers, or
generated-report frequency tables. Dictionary and review-note slices
have not yet begun.

\subsubsection{Numerals in outline}\label{numerals-in-outline}

The current candidate-backed packet supports a narrow generalization.
Simple cardinals can follow a counted noun in examples such as
\tdim{kum nih} and \tdim{ni sagih}. Compound tens are represented by
\tdim{sawmkua}, ordinal formation is represented by \tdim{nihna},
occurrence-counting is represented by \tdim{sawmvei}, and large biblical
number expressions are visible but need caveats. Across the packet,
\tdim{kua} and \tdim{khat} require constructional controls rather than
raw-string treatment.

\subsubsection{Basic cardinal counting
phrases}\label{basic-cardinal-counting-phrases}

The cleanest basic counting phrases in the current packet are Genesis
11:10 and Genesis 7:10.

\begin{exe}
\ex \label{ex:num-kum-nih}
\gll \tdimword{kúm} \tdimword{nìh} \\
     year two \\
\glt 'two years'
\end{exe}

\begin{exe}
\ex \label{ex:num-ni-sagih}
\gll \tdimword{nì} \tdimword{sagih} \\
     day seven \\
\glt 'seven days'
\end{exe}

These are clean post-nominal counting phrases. The printed claim should
stay modest: they show candidate-backed noun-plus-numeral counting, not
a full word-order typology for every numeral construction in the
language.

\subsubsection{Compound tens}\label{compound-tens}

Genesis 5:9 gives the clearest current compound-ten example:

\begin{exe}
\ex \label{ex:num-sawmkua}
\gll \tdimword{kúm} \tdimword{sàwm-kua} \\
     year ten-nine \\
\glt 'ninety years'
\end{exe}

\tdim{Sawmkua} matters for two reasons. First, it shows a compound-ten
structure. Second, it is the strongest current numeral-side control for
\tdim{kua = nine}. The dossier also keeps one analyzer/export caveat
visible: the segmentation and gloss support the compound reading, but
the lemma and POS layer is flattened (\tdim{kum | sawm}, \tdim{N | N}).
That caveat should stay in the prose rather than being silently
normalized away.

\subsubsection{Ordinals}\label{ordinals}

Genesis 7:11 supplies the current ordinal anchor:

\begin{exe}
\ex \label{ex:num-nihna}
\gll \tdimword{nìh-ná} \\
     two-\textsc{nmlz} \\
\glt 'second'
\end{exe}

\tdim{Nihna} is the current candidate-backed ordinal example. The
generated report mentions \tdim{masa} `first', but \tdim{masa} is not
promoted in this slice, and this section does not attempt a full ordinal
paradigm. The dossier also records that \tdim{pos\_span = N} in the
export; that is a label caveat, not a reason to reject the ordinal
analysis.

\subsubsection{Occurrence counting / multiplicative
form}\label{occurrence-counting-multiplicative-form}

Genesis 31:7 keeps one compact occurrence-counting form in view:

\begin{exe}
\ex \label{ex:num-sawmvei}
\gll \tdimword{sàwm-véi} \\
     ten-times \\
\glt 'ten times'
\end{exe}

The generated report paraphrases this as \tdim{vei sawm}, but the
candidate layer uses export-backed \tdim{sawmvei}, and that fused form
should control the present slice. This row is enough to keep
occurrence-counting visible, but it is not enough to build a full
classifier system.

\subsubsection{Large-number phrase}\label{large-number-phrase}

Genesis 5:27 gives one usable large-number phrase, but only with
explicit caveats:

\begin{exe}
\ex \label{ex:num-large}
\gll \tdimword{kúm} \tdimword{zā-kua} \tdimword{lè} \tdimword{kúm} \tdimword{sàwm-gùk} \tdimword{lè} \tdimword{kua} \\
     year hundred-nine and year ten-six and who \\
\glt 'nine hundred sixty and nine years' (Genesis 5:27)
\end{exe}

This construction is numeral, not interrogative. The dossier is
nevertheless explicit about the export caveats: the analyzer compresses
\tdim{za-kua}, and the final \tdim{kua} is glossed as \tdim{who} in the
candidate TSV even though the constructional context is numeral. The
print gloss therefore keeps \tdim{nine} reader-facing while marking the
export gloss in the example line. This makes the row usable with caveat,
but not a fully polished anchor for a larger large-number section.

\subsubsection{\texorpdfstring{\tdim{Kua}
ambiguity}{ ambiguity}}\label{ambiguity}

\tdim{Kua} needs an explicit control note in the printed grammar. In
constructionally numeral contexts such as \tdim{sawmkua} and the Genesis
5:27 large-number phrase, \tdim{kua} is numeral \tdim{nine}. In
interrogative clauses, however, \tdim{kua} is also \tdim{who}.

The blocked control in the numerals packet is Genesis 48:8:

\begin{quote}
Hihte kua ahi hiam?
\end{quote}

That row belongs to the interrogatives packet, not to numerals. Future
numerals prose must therefore not use raw \tdim{kua} hits as numeral
evidence.

\subsubsection{\texorpdfstring{\tdim{Khat} and the numeral/indefinite
boundary}{ and the numeral/indefinite boundary}}\label{and-the-numeralindefinite-boundary}

Genesis 32:24 gives a useful but boundary-marked example:

\begin{exe}
\ex \label{ex:num-mi-khat}
\gll \tdimword{mì} \tdimword{khàt} \\
     person one \\
\glt 'a man' / 'one person'
\end{exe}

\tdim{Mi khat} is analyzer-backed and worth keeping visible, but its
English context ``a man'' shows why it should not be treated as an
uncomplicated bare numeral \tdim{one} example. The packet therefore
keeps it as numeral/indefinite boundary evidence and does not start a
quantifiers retrofit here.

\subsubsection{Distributive reduplication
deferred}\label{distributive-reduplication-deferred}

The generated report claims distributive \tdim{sagih sagih} for `by
sevens', but the candidate layer keeps that material deferred. In the
current analyzer export window, Genesis 7:2 preserves only a single
\tdim{sagih} token. Distributive reduplication therefore remains not
print-ready in this slice, and \tdim{sagih sagih} should not be printed
as though it were already analyzer-backed.

\subsubsection{Editorial summary}\label{editorial-summary-3}

This slice now safely supports six modest claims: noun-plus-numeral
counting examples are printable; compound-ten \tdim{sawmkua} is
candidate-backed; ordinal \tdim{nihna} is candidate-backed;
occurrence-counting \tdim{sawmvei} remains visible with a fused-export
caveat; one large-number phrase is usable with explicit analyzer
caveats; and the packet has explicit controls for \tdim{kua} and
\tdim{khat}.

What remains deferred is equally important: raw frequency counts, full
numeral paradigms, \tdim{masa}, distributive reduplication, a full
classifier system, quantifier overlap, and raw \tdim{kua} matching. The
next step after this grammar slice is the dictionary print slice, while
review-note work has not yet begun.

\subsection{Quantifiers}\label{quantifiers}

\emph{Source slice:
\texttt{output/publication\_review/grammar\_quantifiers\_print\_slice.md}}

\subsubsection{Scope}\label{scope-4}

This is a short print-facing draft section on quantifiers in Tedim Chin,
controlled by \texttt{candidates\_quantifiers.tsv} and
\texttt{dossier\_quantifiers.md}.

It covers only a small candidate-backed set: universal \tdim{khempeuh};
partitive or existential \tdim{pawlkhat} with caveat; \tdim{khat} as
numeral/indefinite boundary evidence; negative-licensed quantifiers
\tdim{kuamah} and \tdim{bangmah}; degree or quantity \tdim{tampi tak};
comparative and intensifier edge rows with \tdim{zaw} and \tdim{mahmah};
deferred \tdim{peuhpeuh} and \tdim{tawm}; and blocked bang-family
false-friend material.

It does not yet attempt a full quantifier system, a full universal or
distributive system, a full indefinite or partitive account, a full
negative-quantifier chapter, a full degree/intensifier/comparative
chapter, or generated-report frequency tables. Dictionary and
review-note slices have not yet begun.

\subsubsection{Quantifiers in outline}\label{quantifiers-in-outline}

The current candidate-backed packet supports a narrow generalization.
\tdim{Khempeuh} currently provides the safest universal-quantifier
anchor. \tdim{Pawlkhat} is useful, but the accepted example remains
partitive or alternative-grouping evidence rather than a plain bare
\tdim{some}. \tdim{Khat} remains boundary evidence shared with numerals.
\tdim{Kuamah} and \tdim{bangmah} are usable only in negative-licensed
contexts. \tdim{Tampi tak} is the main degree or quantity row.
\tdim{Zaw} and \tdim{mahmah} are edge rows only, and \tdim{peuhpeuh}
plus \tdim{tawm} remain visible but deferred.

\subsubsection{\texorpdfstring{Universal quantifier
\tdim{khempeuh}}{Universal quantifier }}\label{universal-quantifier}

Genesis 2:1 supplies the current candidate-backed universal anchor:

\begin{exe}
\ex \label{ex:quant-khempeuh}
\gll \tdimword{vàn-tùng} \tdimword{léi-tùng} \tdimword{lè} \tdimword{à} \tdimword{sùng-à} \tdimword{ōm-tè} \tdimword{khèmpeuh} \\
     sky-on land-on and \textsc{3sg} inside-\textsc{3sg} exist-\textsc{pl} all \\
\glt 'the heavens and the earth, and all that was in them' (Genesis 2:1)
\end{exe}

\tdim{Khempeuh} is the current accepted universal anchor. The example is
a scoped noun phrase, not a license for raw \tdim{khempeuh} harvesting
or for generated-report count claims. The printed claim should therefore
stay modest: the current packet supports one clear universal anchor, not
a complete universal-quantifier system.

\subsubsection{\texorpdfstring{Partitive / existential
\tdim{pawlkhat}}{Partitive / existential }}\label{partitive-existential}

Genesis 32:8 gives the current accepted-with-caveat \tdim{pawlkhat} row:

\begin{exe}
\ex \label{ex:quant-pawlkhat}
\gll \tdimword{pàwl-khàt} \\
     some-one \\
\glt 'one group' / 'one company'
\end{exe}

\tdim{Pawlkhat} is useful, but the context is partitive or alternative
grouping (``one company \ldots{} the other company''), not an
uncomplicated bare \tdim{some}. The dossier is also explicit that the
noisy opening \tdim{Pawlkhatah} token should not be substituted for the
clean later \tdim{pawlkhat} control token.

\subsubsection{\texorpdfstring{\tdim{Khat} and the numeral/indefinite
boundary}{ and the numeral/indefinite boundary}}\label{and-the-numeralindefinite-boundary-1}

Genesis 32:24 provides the current boundary row:

\begin{exe}
\ex \label{ex:quant-mi-khat}
\gll \tdimword{mì} \tdimword{khàt} \\
     person one \\
\glt 'a man' / 'one person'
\end{exe}

\tdim{Mi khat} is reused from the numerals packet as boundary evidence.
It prevents quantifiers from silently absorbing numeral \tdim{khat} as
an article-like quantifier. The present slice therefore does not treat
\tdim{khat} as an uncomplicated quantifier anchor, and it does not
reopen the numerals packet here.

\subsubsection{Negative quantifiers and negation
overlap}\label{negative-quantifiers-and-negation-overlap}

The current packet prints \tdim{kuamah} and \tdim{bangmah} only in
negative-licensed clauses.

\begin{exe}
\ex \label{ex:quant-kuamah}
\gll \tdimword{kuamah} \tdimword{mú} \tdimword{lò} \\
     nobody see.\textsc{i} \textsc{neg} \\
\glt 'he saw nobody' / 'he saw no man'
\end{exe}

\begin{exe}
\ex \label{ex:quant-bangmah-neg}
\gll \tdimword{bangmah} \tdimword{ōm} \tdimword{lò} \tdimword{hì} \\
     nothing exist \textsc{neg} \textsc{decl} \\
\glt 'there is nothing' / 'nothing exists'
\end{exe}

These rows are accepted with caveat only in negative-licensed contexts.
This grammar slice should therefore cross-reference, not reopen, the
stabilized negation packet. Raw \tdim{kuamah} or \tdim{bangmah} hits
should not be treated as print evidence unless the clause is checked and
the negative licensing is explicit.

\subsubsection{Bang-family false
friends}\label{bang-family-false-friends}

The packet also keeps one blocked control visible:

\begin{quote}
Blocked control: \tdim{tua bangmah hi-in}
\end{quote}

This Exodus 27:11 row is not ordinary negative-quantifier evidence. Its
function in the packet is to prevent quantifiers from absorbing
\tdim{bangmah} outside clear negative licensing. That control also
respects the stabilized interrogatives packet, where bang-family false
friends had to be handled separately from ordinary question material.

\subsubsection{\texorpdfstring{Degree / quantity
\tdim{tampi}}{Degree / quantity }}\label{degree-quantity}

Genesis 17:2 gives the current degree or quantity anchor:

\begin{exe}
\ex \label{ex:quant-tampi}
\gll \tdimword{tampi} \tdimword{tak} \\
     many truly \\
\glt 'exceedingly' / 'very many'
\end{exe}

\tdim{Tampi tak} is enough to keep degree or quantity material visible
in the first slice. It should not, however, launch a full adjective or
adverb chapter.

\subsubsection{Comparative and intensifier edge
rows}\label{comparative-and-intensifier-edge-rows}

The packet keeps one comparative and one intensifier row visible, but
only as edge material:

\begin{exe}
\ex \label{ex:quant-zaw-edge}
\gll \tdimword{vang-lìan} \tdimword{zàw} \\
     power-big more \\
\glt 'more powerful' / 'mightier'
\end{exe}

\begin{exe}
\ex \label{ex:quant-mahmah-edge}
\gll \tdimword{hau} \tdimword{màhmàh} \\
     rich very \\
\glt 'very rich'
\end{exe}

\tdim{Zaw} and \tdim{mahmah} are useful boundary evidence. They are not
the basis for a full comparison or intensifier chapter in this slice,
and they should remain caveated edge rows rather than prompts for a
broad degree-modification analysis.

\subsubsection{Deferred material}\label{deferred-material}

Two visible rows remain deferred.

First, \tdim{peuhpeuh} is represented by \tdim{mi peuhpeuh}, but that
row remains deferred because the example behaves more like free-choice
\tdim{whoever / any person} material than like settled
distributive-universal evidence. It is therefore not print-ready in the
present slice.

Second, \tdim{tawm} remains deferred because the current export glosses
it as \tdim{produce}, and the low-quantity reading is still too noisy
for print promotion. The current slice also leaves fuller
universal/distributive, comparative, degree/intensifier, and
negative-quantifier systems deferred. Coordinators and sentence-final
particles remain outside scope.

\subsubsection{Editorial summary}\label{editorial-summary-4}

This slice safely supports six modest claims: \tdim{khempeuh} is the
current universal anchor; \tdim{pawlkhat} is usable as partitive or
existential evidence with caveat; \tdim{khat} is visible only as
numeral/indefinite boundary evidence; \tdim{kuamah} and \tdim{bangmah}
are usable as negative-licensed quantifier evidence; \tdim{tampi tak} is
the current degree/quantity anchor; and \tdim{zaw} plus \tdim{mahmah}
remain caveated edge rows.

What remains deferred is equally important: raw frequency counts,
\tdim{peuhpeuh} as settled distributive-universal evidence, \tdim{tawm}
as settled low-quantity evidence, \tdim{bangmah} outside clear negative
licensing, \tdim{khat} as an uncomplicated article-like quantifier, a
full degree/intensifier/comparative chapter, and any move into
coordinators or sentence-final particles. The next step after this
grammar slice is the dictionary print slice, while review-note work has
not yet begun.

\reviewchapterbreak

\section{Predicate structure and verbal
morphology}\label{predicate-structure-and-verbal-morphology}

\subsection{Stem alternation}\label{stem-alternation}

\emph{Source slice:
\texttt{output/publication\_review/grammar\_stem\_alternation\_print\_slice.md}}

\subsubsection{Verb-stem alternation}\label{verb-stem-alternation}

This file is now a \textbf{draft argument plan}, not the final polished
grammar section. The eventual prose should not be organized only by
print status. It should discuss \textbf{both syntactic contexts and
individual verb pairs}, in that order:

\begin{enumerate}
\def\labelenumi{\arabic{enumi}.}
\tightlist
\item
  system-wide Form I / Form II distribution by syntactic context;
\item
  the strongest illustrative pairs;
\item
  pair-by-pair caveated and difficult evidence;
\item
  one-sided, control, and rejected material.
\end{enumerate}

Earlier descriptions agree that Tedim has a two-form verbal system, even
though the terminology differs. Henderson describes Form I and Form II,
while Zam Ngaih Cing speaks of Stem 1 and Stem 2
\citep{henderson1965, zamngaihcing2017}. The present review packet keeps
those descriptive traditions in view, but it now separates coverage,
promotion, and quotation safety much more explicitly than the earlier
broad analyzer table did.

\paragraph{The Form I / Form II
contrast}\label{the-form-i-form-ii-contrast}

The opening section of the eventual grammar should introduce the system
with the strongest lexical pairs, not with the widest possible
inventory.

\begin{itemize}
\tightlist
\item
  \textbf{Core starting point}: \tdim{mu \textasciitilde{} muh},
  \tdim{ne \textasciitilde{} nek}, \tdim{nei \textasciitilde{} neih}.
\item
  \textbf{Basic descriptive claim}: Form I is especially clear in
  ordinary finite predication and many main-clause uses. Form II is
  especially visible in dependent, temporal, purposive, attributive,
  nominalized, and other constructionally bound environments.
\item
  \textbf{Immediate caution}: do not reduce the system to slogans such
  as \emph{Form II = subordinate}, \emph{Form II = negative}, \emph{Form
  I = finite}, or \emph{Form II = nominalized only}.
\end{itemize}

The argument here should stay structural and distributional before it
becomes lexical and pair-specific.

The eventual section will likely need five editorial surfaces:

\begin{itemize}
\tightlist
\item
  a \textbf{small core showcase table};
\item
  a \textbf{larger promoted-pair inventory table};
\item
  a \textbf{pair-by-pair discussion section};
\item
  a \textbf{one-sided / same-form / functional coverage table};
\item
  a \textbf{blocked or analyzer-noise table / appendix paragraph}.
\end{itemize}

\paragraph{Distribution by syntactic
context}\label{distribution-by-syntactic-context}

The main organizing file for this part of the write-up is
\texttt{output/publication\_review/stem\_alternation\_syntactic\_context\_matrix.tsv}.
The matrix is for organizing claims and subsection order; the citation
shortlist remains the source of quotation candidates. The prose should
work through the following contexts in order, using the matrix rather
than treating all environments as equivalent evidence.

\subparagraph{Finite and main-clause
uses}\label{finite-and-main-clause-uses}

Start with ordinary finite predication. This is where Form I is
clearest, even though some lexical pairs still show real Form II matrix
tokens.

\subparagraph{Imperatives and
directives}\label{imperatives-and-directives}

Use imperatives and directives to reinforce the ordinary clause-force
profile of Form I, while noting that some lexical pairs still show
marked Form II directive material.

\subparagraph{Negative clauses}\label{negative-clauses}

Negative clauses belong in the argument, but they are \textbf{not} a
single diagnostic. The prose should say that negative environments
matter, not that Form II simply equals negation.

\subparagraph{\texorpdfstring{Dependent temporal clauses with
\tdim{ciangin}}{Dependent temporal clauses with }}\label{dependent-temporal-clauses-with}

This is one of the strongest places to show Form II in clearly dependent
syntax. It should be one of the central distributional subsections.

\subparagraph{\texorpdfstring{Temporal and nominal \tdim{ni-in}
contexts}{Temporal and nominal  contexts}}\label{temporal-and-nominal-contexts}

Use \tdim{ni-in} as a second temporal construction so the grammar does
not make \tdim{ciangin} carry the whole dependent-clause argument.

\subparagraph{\texorpdfstring{Clause-linking with
\tdim{kipan}}{Clause-linking with }}\label{clause-linking-with}

This subsection is especially important for pairs such as
\tdim{nusia \textasciitilde{} nusiat}, where the best Form II evidence
is constructionally bound and source-linking rather than a neat finite
contrast.

\subparagraph{\texorpdfstring{Purposive
\tdim{nadingin}}{Purposive }}\label{purposive}

Purposive or irrealis-heavy \tdim{ding / nadingin} clauses are one of
the clearest places where Form II becomes prominent across several
promoted pairs.

\subparagraph{\texorpdfstring{Nominalized
\tdim{-na}}{Nominalized }}\label{nominalized}

Nominalized material should be discussed directly, but with an explicit
filter against lexicalized nouns and compounds.

\subparagraph{\texorpdfstring{Attributive and relative \tdim{mi}
contexts}{Attributive and relative  contexts}}\label{attributive-and-relative-contexts}

These contexts belong with the broader non-final distribution. They
matter, but they should not be turned into a mechanical rule.

\subparagraph{Possessed and genitive attributive
contexts}\label{possessed-and-genitive-attributive-contexts}

This section should stay compact and should probably center on
\tdim{nei \textasciitilde{} neih}, since that pair gives the clearest
evidence here.

\subparagraph{Modal and ability uses}\label{modal-and-ability-uses}

This subsection should explicitly center on
\tdim{thei \textasciitilde{} theih}, since the pair is real but
constructionally special.

\subparagraph{Quotative and say-complement
contexts}\label{quotative-and-say-complement-contexts}

Discuss this environment mainly to explain why functional material such
as \tdim{ci \textasciitilde{} cih} should not be merged blindly into the
lexical-verb showcase table.

\subparagraph{Derived, causative, compound, and lexicalized
material}\label{derived-causative-compound-and-lexicalized-material}

This should be an explicit filter section, not a positive-evidence
section. The grammar needs to show why \tdim{-sak} material, compounds,
lexicalized families, and review-bucket rows cannot simply be counted as
ordinary stem evidence.

\paragraph{Core showcase pairs}\label{core-showcase-pairs}

The core pair subsection should remain narrow and should illustrate the
system with the strongest three pairs:

\begin{itemize}
\item
  \tdim{mu \textasciitilde{} muh}
\item
  \tdim{ne \textasciitilde{} nek}
\item
  \tdim{nei \textasciitilde{} neih}
\end{itemize}

These pairs should do the main pedagogical work for the Form I / Form II
contrast.

This should become the \textbf{small core showcase table}, not the full
promoted inventory.

\paragraph{Promoted caveated pairs}\label{promoted-caveated-pairs}

The next section should move to the broader but still promoted
inventory:

\begin{itemize}
\item
  \tdim{za \textasciitilde{} zak}
\item
  \tdim{pia \textasciitilde{} piak}
\item
  \tdim{nusia \textasciitilde{} nusiat}
\item
  \tdim{bia \textasciitilde{} biak}
\item
  \tdim{thei \textasciitilde{} theih}
\item
  \tdim{piang \textasciitilde{} pian}
\item
  \tdim{zui \textasciitilde{} zuih}
\item
  \tdim{khial \textasciitilde{} khialh}
\item
  \tdim{kia \textasciitilde{} kiak}
\item
  \tdim{sawlkhia \textasciitilde{} sawlkhiat}
\end{itemize}

Each pair should get a short paragraph or table note answering two
questions:

\begin{enumerate}
\def\labelenumi{\arabic{enumi}.}
\tightlist
\item
  \textbf{What does the evidence positively show?}
\item
  \textbf{What is the specific caveat?}
\end{enumerate}

The point is not to collapse them into one generic ``caveated examples''
bucket, but to say what kind of Form II evidence each pair contributes.

This should likely become the \textbf{larger promoted-pair inventory
table}, with short editorial notes rather than long prose for every row.

\paragraph{Difficult but grammatically important
pairs}\label{difficult-but-grammatically-important-pairs}

This section should keep the difficult cases in the argument rather than
silently dropping them:

\begin{itemize}
\item
  \tdim{ngai \textasciitilde{} ngaih}
\item
  \tdim{pua \textasciitilde{} puak}
\item
  \tdim{pai \textasciitilde{} paih}
\item
  \tdim{tua \textasciitilde{} tuah}
\item
  \tdim{tua \textasciitilde{} tuak}
\end{itemize}

These are not simply rejected. They matter because they show
lexical-family contamination, shared-base problems, one-sided evidence,
and analyzer overgeneration.

\paragraph{One-sided Bible attestations and questionnaire
controls}\label{one-sided-bible-attestations-and-questionnaire-controls}

This section should separate two things that were too easily blurred in
earlier drafts:

\begin{enumerate}
\def\labelenumi{\arabic{enumi}.}
\tightlist
\item
  \textbf{One-sided or constructionally skewed lexical verbs} such as
  \tdim{bawl \textasciitilde{} bawlh},
  \tdim{dipkua \textasciitilde{} dipkuat},
  \tdim{gen \textasciitilde{} genh},
  \tdim{hawlkhia \textasciitilde{} hawlkhiat},
  \tdim{husia \textasciitilde{} husiat},
  \tdim{kho \textasciitilde{} khoh},
  \tdim{kido \textasciitilde{} kidot},
  \tdim{lua \textasciitilde{} luah}, \tdim{tu \textasciitilde{} tuh},
  \tdim{tuahpha \textasciitilde{} tuahphat}, and
  \tdim{vial \textasciitilde{} vialh}.
\item
  \textbf{Same-form questionnaire controls} such as
  \tdim{dawn \textasciitilde{} dawn},
  \tdim{hong \textasciitilde{} hong}, \tdim{om \textasciitilde{} om},
  \tdim{ci \textasciitilde{} ci}, \tdim{hi \textasciitilde{} hi},
  \tdim{bawl \textasciitilde{} bawl}, \tdim{zui \textasciitilde{} zui},
  \tdim{pai \textasciitilde{} pai}, and the other same-form rows now
  tracked in the lexical inventory.
\end{enumerate}

The prose should say clearly that same-form questionnaire rows are
useful controls, but they are not overt alternating pairs in the Bible
layer.

These rows are \textbf{not} discarded. They should stay visible in a
coverage table because the grammar needs to account for verbs flagged by
the literature, questionnaire, analyzer inventory, and Bible audit even
when only one side is cleanly attested in the current Bible layer.

\paragraph{Rejected/non-verbal/analyzer-noise
cases}\label{rejectednon-verbalanalyzer-noise-cases}

This section should explain why some apparent pairs stay blocked:

\begin{itemize}
\tightlist
\item
  nouns and nominal compounds;
\item
  lexicalized or compound families;
\item
  derivational \tdim{-sak} material;
\item
  homophones and shared-base contamination;
\item
  category mismatches;
\item
  analyzer overgeneration.
\end{itemize}

Representative blocked rows should include:

\begin{itemize}
\item
  \tdim{keu \textasciitilde{} keuh}
\item
  \tdim{khai \textasciitilde{} khaih}
\item
  \tdim{sia \textasciitilde{} siah}
\item
  \tdim{tan \textasciitilde{} tanh}
\item
  \tdim{mual \textasciitilde{} mualh}
\item
  \tdim{sum \textasciitilde{} sumh}
\item
  \tdim{thu \textasciitilde{} thuh}
\item
  \tdim{lampi \textasciitilde{} lampih}
\item
  \tdim{khua \textasciitilde{} khuat}
\item
  \tdim{gamla \textasciitilde{} gamlat}
\end{itemize}

\paragraph{How the evidence files should be
used}\label{how-the-evidence-files-should-be-used}

The final prose should keep the current evidence layers distinct:

\begin{itemize}
\tightlist
\item
  \texttt{output/publication\_review/stem\_alternation\_citation\_shortlist.tsv}
  = the \textbf{only quotation-safe layer} for printed examples.
\item
  \texttt{output/publication\_review/stem\_alternation\_syntactic\_context\_matrix.tsv}
  = the basis for the \textbf{context-by-context distributional
  argument} and for organizing claims by subsection.
\item
  \texttt{output/publication\_review/stem\_alternation\_pair\_discussion\_plan.tsv}
  = the basis for the \textbf{pair-by-pair discussion order and claims}.
\item
  \texttt{output/publication\_review/stem\_alternation\_lexical\_inventory.tsv}
  = the coverage layer for lexical category, promotion status, blockers,
  and source basis.
\item
  \texttt{output/publication\_review/stem\_alternation\_promotable\_examples.tsv}
  and
  \texttt{output/publication\_review/stem\_alternation\_manual\_promotion\_review.tsv}
  = editorial support files for promotion and example triage.
\item
  \texttt{output/publication\_review/stem\_alternation\_corpus\_audit.tsv}
  = generated local background evidence; useful for checking rows, but
  not a quotation-safe layer and not a tracked argument table.
\end{itemize}

\paragraph{Writing order for the eventual
section}\label{writing-order-for-the-eventual-section}

The actual grammar section should be drafted in this order:

\begin{enumerate}
\def\labelenumi{\arabic{enumi}.}
\tightlist
\item
  \textbf{System-wide syntactic distribution}
\item
  \textbf{Best-attested showcase pairs}
\item
  \textbf{Promoted and difficult pair-by-pair discussion}
\item
  \textbf{One-sided / same-form / functional coverage table}
\item
  \textbf{Blocked or analyzer-noise appendix paragraph}
\end{enumerate}

That is the architecture this packet should now support.

\paragraph{Next prose draft}\label{next-prose-draft}

The next commit should create a separate prose draft file rather than
overwrite this planning file, probably at
\texttt{output/publication\_review/grammar\_stem\_alternation\_section\_draft.md}.

That draft should be organized as:

\begin{enumerate}
\def\labelenumi{\arabic{enumi}.}
\tightlist
\item
  \textbf{System-wide syntactic distribution}
\item
  \textbf{Core showcase examples}
\item
  \textbf{Promoted-pair inventory}
\item
  \textbf{Pair-by-pair notes for promoted and difficult pairs}
\item
  \textbf{One-sided / same-form / functional coverage table}
\item
  \textbf{Blocked/noise appendix paragraph}
\end{enumerate}

\subsection{Verb paradigms}\label{verb-paradigms}

{[}MAJOR GAP: verb paradigms remain report-backed but not
packet-shaped.{]}

\texttt{docs/grammar/reports/05-verb-00-paradigm-tables.md} remains part
of the evidence base, but it has not yet been converted into a
review-note-stage packet with an assembled grammar slice.

\subsection{Prefix / agreement}\label{prefix-agreement}

\emph{Source slice:
\texttt{output/publication\_review/grammar\_prefix\_agreement\_print\_slice.md}}

\subsubsection{Editorial scope}\label{editorial-scope-3}

This is the first narrow prefix/agreement grammar slice for Tedim. It is
controlled by
\texttt{output/publication\_review/candidates\_prefix\_agreement.tsv}
and
\texttt{output/publication\_review/dossier\_prefix\_agreement\_scope.md}.
Supporting/background evidence comes from
\texttt{docs/grammar/reports/05-verb-03-agreement.md},
\texttt{docs/grammar/reports/04-np-07-possession.md},
\texttt{docs/grammar/morphemes/01-prefixes.md},
\texttt{docs/grammar/lit-reviews/04-np-07-possession-lit.md},
\texttt{docs/grammar/DISAMBIGUATION.md}, and the regression evidence in
\texttt{tests/test\_prefix\_agr\_poss.py}.

This is not a full agreement chapter, not a full possession chapter, not
a full object-prefix or inverse chapter, and not a rewrite of the
completed pronouns/clusivity packet. It also stays narrow against
\texttt{output/publication\_review/review\_notes\_pronouns.md},
\texttt{output/publication\_review/review\_notes\_derivation\_valency.md},
and
\texttt{output/publication\_review/review\_notes\_vp\_structure\_stacking.md}.

The present slice therefore covers only the agreement-versus-possession
routing contrast, with \tdim{kanei} as the clearest agreement anchor and
\tdim{kainn} as the clearest possessive-routing anchor. No dictionary
slice exists yet for prefix/agreement, because this packet is still
establishing a controlled routing claim rather than a lexical headword
layer. The packet now proceeds through review notes rather than through
a lexical headword layer.

\subsubsection{Agreement versus possession
routing}\label{agreement-versus-possession-routing}

The first safe prefix/agreement claim is a routing contrast. Before
verbs, the shared pronominal prefix family may be routed as agreement;
before nouns, the same family may be routed as possession.

\tdim{Kanei} and \tdim{kainn} are the core pair for that contrast.
\tdim{Kanei} keeps the prefix family on a verbal host, while
\tdim{kainn} keeps it on a nominal host.
\texttt{tests/test\_prefix\_agr\_poss.py} is the key regression control
here, because it explicitly requires verb-side AGR glossing to stay
distinct from noun-side POSS glossing.

That is enough for the first print-facing claim. The slice does not need
to resolve every larger prefix question before stating that host type
already controls a safe agreement-versus-possession routing contrast in
the current candidate layer.

\subsubsection{Agreement anchor: kanei}\label{agreement-anchor-kanei}

\tdim{Kanei} is the clearest verbal agreement anchor in the packet. The
candidate TSV marks it as the main AGR-side row, and
\texttt{tests/test\_prefix\_agr\_poss.py} protects the glossing as
\tdim{ka-nei} / \tdim{1SG-have}.

The grammar claim here is deliberately limited. At the current slice
maturity level, \tdim{kanei} supports only the routing statement that
the shared prefix family can surface as verbal agreement before a verb
host. This is strong enough for a narrow print slice, but still smaller
than a full agreement chapter or a full prefix paradigm.

\subsubsection{Possessive anchor: kainn}\label{possessive-anchor-kainn}

\tdim{Kainn} is the clearest possessive-routing anchor in the packet.
The candidate TSV marks it as the nominal counterpart to \tdim{kanei},
and \texttt{tests/test\_prefix\_agr\_poss.py} protects the glossing as
\tdim{ka-inn} / \tdim{1SG.POSS-house}.

The grammar claim again stays small. At the current slice maturity
level, \tdim{kainn} supports the routing statement that the same prefix
family can be analyzed as possessive before a noun host. This is enough
to justify a first print-facing contrast without pretending that the
project already has a full possession chapter or a full possessor-syntax
account.

\subsubsection{Why this is not just pronouns
again}\label{why-this-is-not-just-pronouns-again}

The completed pronouns/clusivity packet already handles independent
pronouns, clusivity, and the broader pronoun paradigm through
\texttt{output/publication\_review/review\_notes\_pronouns.md}.

This slice is doing something narrower. It is about prefix routing
across verbal and nominal hosts, not about reopening the
independent-pronoun paradigm. That is why \tdim{kanei} and \tdim{kainn}
are better first anchors here than \tdim{ipai}, \tdim{ko}, \tdim{ei}, or
any broader person-paradigm table.

\subsubsection{Boundary material}\label{boundary-material-2}

The rest of the candidate packet stays outside the first grammar slice
because each row is still dominated by another unresolved boundary.

\tdim{ainn} stays outside because the broader \tdim{a-} family overlaps
with verbal agreement, relativizer-like material, and other domains. It
is useful boundary evidence, but not the first clean routing anchor.

\tdim{ipai} stays outside because inclusive/exclusive \tdim{i-} material
belongs first to the completed pronouns/clusivity packet rather than to
this first routing slice.

\tdim{hongmu} and \tdim{kongmu} stay outside because object-prefix or
inverse-like material still needs a later dedicated sub-scope with
tighter directional and inverse controls.

\tdim{kipan} stays outside because \tdim{ki-} reflexive or middle
material remains boundary-only between prefix/agreement and
derivation/valency.

Apostrophe possession and broader possessor syntax also stay outside
because this is not a full possession chapter.

\subsubsection{Safe first-slice claim}\label{safe-first-slice-claim-2}

At the current slice maturity level, the safest prefix/agreement claim
is that Tedim has candidate-controlled evidence for routing a shared
pronominal prefix family differently by host type: \tdim{kanei} supports
verbal agreement routing, while \tdim{kainn} supports nominal possessive
routing.

That claim is deliberately smaller than a full agreement chapter,
smaller than a full possession chapter, smaller than a full
object-prefix or inverse chapter, and smaller than a rewritten pronoun
packet.

\subsubsection{Recommended next step}\label{recommended-next-step-3}

This packet now properly proceeds to prefix/agreement review notes
rather than to a dictionary slice, because it is a routing/analysis
packet rather than a lexical headword packet.

If the project later wants one more prefix step after review notes and
human review, the next sub-scope should be a separate hong-/kong-
object-prefix or inverse candidate expansion rather than a dictionary
layer.

\subsection{Transitivity}\label{transitivity}

\emph{Source slice:
\texttt{output/publication\_review/grammar\_transitivity\_print\_slice.md}}

\subsubsection{Editorial scope}\label{editorial-scope-4}

This is the first narrow transitivity grammar slice. It is controlled by
\texttt{output/publication\_review/candidates\_transitivity.tsv} and
\texttt{output/publication\_review/dossier\_transitivity\_scope.md}.
Supporting/background evidence comes from
\texttt{docs/grammar/reports/05-verb-12-transitivity.md}.

Boundary control comes from
\texttt{output/publication\_review/review\_notes\_derivation\_valency.md},
\texttt{output/publication\_review/review\_notes\_stem\_alternation.md},
\texttt{output/publication\_review/review\_notes\_prefix\_agreement.md},
\texttt{output/publication\_review/review\_notes\_vp\_structure\_stacking.md},
\texttt{output/publication\_review/review\_notes\_tam.md}, and
\texttt{output/publication\_review/review\_notes\_case\_marking.md}.

This is not a full valency chapter, not a full verb-class chapter, not a
dictionary slice, and not a full argument-structure account.

\subsubsection{Clean intransitive anchor:
sih}\label{clean-intransitive-anchor-sih}

\tdim{sih} is the clean intransitive anchor for the first slice. The
controlled form and gloss/function are \tdim{sih / die}.

The claim stays narrow.
\texttt{output/publication\_review/candidates\_transitivity.tsv} and
\texttt{output/publication\_review/dossier\_transitivity\_scope.md}
treat \tdim{sih} as the safest current intransitive anchor because it is
a simple ABS-led event row that does not immediately force
derivation/valency, prefix/agreement, or case-marking analysis into the
first print-facing claim.

\subsubsection{Supporting intransitive row:
suak}\label{supporting-intransitive-row-suak}

\tdim{suak} is supporting intransitive evidence rather than the sole
anchor. The controlled form and gloss/function are \tdim{suak / become}.

That makes \tdim{suak} useful as a second row for the intransitive side
of the contrast. It supports a narrow change-of-state intransitive
reading without forcing the slice to generalize over the whole
report-level intransitive inventory.

\subsubsection{Clean transitive anchor:
hawl}\label{clean-transitive-anchor-hawl}

\tdim{hawl} is the clean transitive anchor for the first slice. The
controlled form and gloss/function are \tdim{hawl / seek}.

The claim here also stays narrow. \tdim{hawl} is the safest current
transitive anchor because it gives the packet a simple transitive event
row without leaning on derivation-heavy material such as \tdim{piangsak}
or stem-family material such as \tdim{ngai / ngaih}. The report
percentages in \texttt{docs/grammar/reports/05-verb-12-transitivity.md}
are useful discovery evidence, but they are not treated here as
categorical proof.

\subsubsection{Supporting transitive row:
en}\label{supporting-transitive-row-en}

\tdim{en} is supporting transitive evidence rather than the leading
transitive claim. The controlled form and gloss/function are
\tdim{en / look.at}.

That makes \tdim{en} useful as secondary support for the transitive side
of the slice. It helps confirm the narrow contrast, but lower counts and
perception-verb semantics make it weaker than \tdim{hawl} as a
chapter-leading anchor.

\subsubsection{Why ambitransitive/labile material is not yet the first
slice}\label{why-ambitransitivelabile-material-is-not-yet-the-first-slice}

\tdim{mu / muh} remains visible boundary evidence for alternation or
labile-looking behavior. It matters because it shows that the
transitivity report is not exhausted by a clean binary class split.

It stays outside the first print-facing claim, however, because
\tdim{mu / muh} overlaps directly with Form I / Form II stem
alternation. The same caution applies to \tdim{za / zak},
\tdim{nei / neih}, and \tdim{ngai / ngaih}, which remain candidate-layer
or boundary material rather than anchors for this commit.

\subsubsection{Boundary material}\label{boundary-material-3}

The following stay outside the first grammar slice:

\begin{itemize}
\item
  \tdim{mu / muh}
\item
  \tdim{za / zak}
\item
  \tdim{nei / neih}
\item
  \tdim{ngai / ngaih}
\item
  \tdim{piangsak}
\item
  \tdim{pia}
\item
  \tdim{gen}
\item
  \tdim{tom}
\item
  \tdim{hong}
\item
  \tdim{ki}
\item
  \tdim{dawt}
\item
  \tdim{bei}
\item
  \tdim{pia(k)sak}
\item
  case-dominated rows
\item
  derivation-heavy rows
\item
  prefix/agreement-heavy rows
\item
  analyzer-noisy, lexicalized, report-only, or whole-system verb-class
  claims
\end{itemize}

\tdim{piangsak} remains outside because it is derivation/valency-heavy
and lexicalized-looking. \tdim{pia}, \tdim{gen}, and \tdim{tom} remain
outside because they are dominated by broader case-marking and
argument-structure interpretation. \tdim{hong} and \tdim{ki} remain
outside because their report behavior is too mixed for the first clean
intransitive claim. \tdim{dawt}, \tdim{bei}, and \tdim{pia(k)sak} remain
outside because they are too thin, too noisy, or too report-bound for
the first print-facing contrast.

\subsubsection{Safe first-slice claim}\label{safe-first-slice-claim-3}

At the current slice maturity level, the safest transitivity claim is
that Tedim has candidate-controlled evidence for a narrow
intransitive/transitive contrast, with \tdim{sih} and \tdim{suak}
supporting the intransitive side and \tdim{hawl}, with secondary support
from \tdim{en}, supporting the transitive side.

Alternation, labile behavior, stem alternation, derivation/valency,
prefix/agreement, case-marking, and whole-system verb-class claims
remain candidate-layer or boundary material.

\subsubsection{Recommended next step}\label{recommended-next-step-4}

After this grammar slice, the next step should be transitivity review
notes rather than a dictionary slice, because this packet is
grammar-facing and argument-structure-oriented rather than lexical.

If more transitivity work is chosen before review notes, the next
sub-scope should be \tdim{mu / muh} as a labile or stem-alternation
boundary packet, but not in this commit.

\subsection{VP structure / suffix
stacking}\label{vp-structure-suffix-stacking}

\emph{Source slice:
\texttt{output/publication\_review/grammar\_vp\_structure\_stacking\_print\_slice.md}}

\subsubsection{Editorial scope}\label{editorial-scope-5}

This is the first narrow VP structure / suffix stacking grammar slice
for Tedim. It is controlled by
\texttt{output/publication\_review/candidates\_vp\_structure\_stacking.tsv}
and
\texttt{output/publication\_review/dossier\_vp\_structure\_stacking\_scope.md}.
Supporting/background evidence comes from
\texttt{docs/grammar/reports/05-verb-02-vp-structure.md},
\texttt{docs/grammar/reports/05-verb-10-combinations.md}, and the
regression evidence in \texttt{tests/test\_vp\_slots.py}.

This is not a full VP chapter. It is not a rewrite of TAM, directionals,
negation, sentence-final particles, or relators/postpositions. Those
packet boundaries remain explicit through
\texttt{output/publication\_review/review\_notes\_tam.md},
\texttt{output/publication\_review/review\_notes\_directionals.md},
\texttt{output/publication\_review/review\_notes\_negation.md},
\texttt{output/publication\_review/review\_notes\_sentence\_final\_particles.md},
and
\texttt{output/publication\_review/review\_notes\_relators\_postpositions.md}.

The present slice therefore covers only the first safe suffix-stacking
claim: \tdim{bawlzoding} as the central print-usable-with-caveat anchor
for aspect plus irrealis stacking. No ordinary dictionary slice exists
for this packet because it is constructional rather than lexical.

\subsubsection{Baseline: completed single-suffix
packets}\label{baseline-completed-single-suffix-packets}

The packet starts from two completed single-suffix baselines that are
already owned elsewhere.

\tdim{bawlzo} remains the compact V+ASPECT baseline already owned by the
TAM packet. It shows that the repository already has a stable completive
anchor, but it should be used here only as baseline evidence rather than
reopened as new VP-structure prose.

\tdim{pokhia} remains the compact V+DIR baseline already owned by the
directionals packet. It shows that the repository already has a stable
post-stem directional anchor, but it should be used here only as
baseline evidence rather than widened into a new VP-slot claim.

These baseline rows matter because the first VP slice should not
redescribe already-completed packets. Their role here is only to show
that compact single-suffix material is already controlled before a
multi-suffix claim is added.

\subsubsection{First stacking anchor: aspect plus
irrealis}\label{first-stacking-anchor-aspect-plus-irrealis}

\tdim{bawlzoding} is the central first-slice stack. The candidate layer
marks it as \tdim{print\_usable\_with\_caveat}, and
\texttt{tests/test\_vp\_slots.py} already keeps it visible as
aspect-plus-modal regression evidence.

The safe print claim is deliberately narrow. \tdim{bawlzoding} supports
the ordering observation:

\begin{quote}
verb stem + completive/aspectual material + irrealis/modal material
\end{quote}

In other words, the current slice supports a small VP-structure claim
that aspectual material can precede irrealis/modal material inside a
multi-suffix verbal complex.

For the present packet, that means reading \tdim{bawlzoding} as a
constructional stack rather than over-trusting every analyzer label on
the row. The current analyzer gloss is noisy (\texttt{make-south-IRR}),
so the slice should not pretend that the middle gloss is already a
perfect semantic label. The point of the row is its suffix-stacking
evidence: the repository already has a regression-backed form in which a
verbal stem is followed by completive/aspectual material and then by
irrealis/modal material.

This is enough for a first grammar slice, but not enough for a full VP
slot template. The present packet therefore treats \tdim{bawlzoding} as
the clearest current anchor for aspect plus irrealis stacking and stops
there.

\subsubsection{Boundary and deferred
stacks}\label{boundary-and-deferred-stacks}

The rest of the visible stacks remain outside the first core slice
because each one is dominated by another packet boundary.

\tdim{khia-ta} is useful boundary evidence because it shows real
TAM/directional overlap, but it is not the first-slice core. The
completed TAM and directionals packets already keep this row visible as
overlap evidence, and this slice should not reopen either packet through
that row.

\tdim{ciahsakkik}, \tdim{bawlsakthei}, and \tdim{paikhiatsak} are real
multi-suffix complexes, but they are derivation/valency-heavy stacks.
Once \tdim{-sak} and other valency-changing material become central, the
next packet boundary is derivation/valency rather than narrow VP
stacking. These rows therefore stay deferred until the
derivation/valency packet is explicitly selected.

\tdim{khiathei ding om lo} is also real overlap evidence, but it is
TAM-negation overlap rather than a clean first VP-stack anchor. The
completed negation and TAM packets already control that territory, so it
should not be promoted as the model for core VP-stacking prose.

\tdim{dingin} remains clause-bound irrealis or subordination material.
It is important because it shows how quickly visible verbal stacking can
turn into clause-linkage evidence, but it should wait for a
subordination packet rather than being promoted here as simple VP suffix
stacking.

\subsubsection{Safe first-slice claim}\label{safe-first-slice-claim-4}

At the current slice maturity level, the safest VP-structure claim is
that Tedim permits at least some multi-suffix verbal complexes in which
aspectual material precedes irrealis/modal material, with
\tdim{bawlzoding} as the clearest current anchor.

That claim stays deliberately smaller than a full slot template. It does
not claim that every visible stack in the reports belongs to one solved
VP order, and it does not reopen completed TAM, directional, negation,
sentence-final, or relator/postposition packets.

\subsubsection{Recommended next step}\label{recommended-next-step-5}

After this grammar slice, the packet can move to review notes without
forcing an ordinary dictionary slice, because the safe first-slice claim
is constructional rather than lexical.

The next substantive packet after this narrow constructional slice
should therefore be derivation/valency candidate scoping rather than a
broader VP rewrite.

\subsection{TAM / aspect / modal}\label{tam-aspect-modal}

\emph{Source slice:
\texttt{output/publication\_review/grammar\_tam\_print\_slice.md}}

\subsubsection{Editorial scope}\label{editorial-scope-6}

This is the first narrow TAM / aspect / modal grammar slice for Tedim
Chin, controlled by \texttt{candidates\_tam.tsv} and
\texttt{dossier\_tam\_scope.md}. It is not a full TAM chapter.

The first-slice TAM anchors are limited to \tdim{-ngei}, \tdim{-gige},
\tdim{-zel}, \tdim{-ta}, \tdim{-zo}, \tdim{-kik}, \tdim{-ding}, and
\tdim{-thei}, represented here by \tdim{paingei}, \tdim{neigige},
\tdim{paizel}, \tdim{kilawmta}, \tdim{bawlzo}, \tdim{hongpaikik},
\tdim{omding}, and \tdim{bawlthei}. This slice therefore stays with
compact suffixal anchors already marked print-ready or print-usable in
the candidate TSV instead of widening into broad clause-structure,
sentence-final, directional, or VP-slot prose. The dictionary slice now
exists, but review-note work has not yet begun.

\subsubsection{Experiential and habitual
anchors}\label{experiential-and-habitual-anchors}

\tdim{Paingei} (\tdim{pai-ngei} -\textgreater{} \tdim{go-EXP}) is the
current experiential anchor. It supports the narrow claim that
\tdim{-ngei} can mark experiential meaning such as `have V-ed before'
when it stays tied to a compact verbal host.

\tdim{Neigige} (\tdim{nei-gige} -\textgreater{} \tdim{have-HAB}) is the
current habitual anchor. It supports a modest print claim that
\tdim{-gige} can mark habitual or regularly repeated action.

\tdim{Paizel} (\tdim{pai-zel} -\textgreater{} \tdim{go-HAB.CONT}) is the
current habitual-continuative anchor. It is usable, but it borders
broader continuative aspect and should stay construction-controlled. The
first TAM slice should therefore keep \tdim{-zel} tied to compact
anchors such as \tdim{paizel} rather than turning every
continuative-looking or lexicalized \tdim{zel} sequence into TAM
evidence.

\subsubsection{Compact aspectual
anchors}\label{compact-aspectual-anchors}

\tdim{Kilawmta} (\tdim{ki-lawm-ta} -\textgreater{}
\tdim{REFL-worthy-PFV}) is the current compact perfective anchor. It
supports a modest completed-event reading for \tdim{-ta}, but only with
sentence-final overlap caveat. The slice relies on compact suffixed
verbs, not bare \tdim{ta}, and it keeps sentence-final overlap material
such as \tdim{mangngilh ta hi} outside the core TAM evidence.

\tdim{Bawlzo} (\tdim{bawl-zo} -\textgreater{} \texttt{make-COMPL}) is
the current compact completive anchor. It supports a modest `finish
V-ing' or completive reading for \tdim{-zo}, but only with the
bare-\tdim{zo} and sentence-final overlap caveat kept explicit. The
slice therefore treats \tdim{bawlzo} as usable compact evidence while
refusing to turn bare \tdim{zo} into default TAM proof.

\tdim{Hongpaikik} (\tdim{hong-pai-kik} -\textgreater{}
\tdim{3→1-go-ITER}) is the current iterative anchor. It is usable as
iterative or `again' evidence, but it should not become a motion/return
chapter. The first slice keeps \tdim{-kik} narrow and
construction-backed rather than expanding it into every return- or
motion-flavored verbal sequence.

\subsubsection{Compact modal anchors}\label{compact-modal-anchors}

\tdim{Omding} (\tdim{om-ding} -\textgreater{} \tdim{exist-IRR}) is the
current compact \tdim{-ding} anchor. It supports a narrow irrealis /
future / modal claim, but only with the explicit dingin and clause-bound
caveat. \tdim{Dingin} and other clause-bound \tdim{-ding} material stay
outside the starting print slice, so \tdim{omding} remains the safer
anchor than report-style purposive or clause-linking material.

\tdim{Bawlthei} (\tdim{bawl-thei} -\textgreater{} \texttt{make-ABIL}) is
the current compact abilitative anchor. It supports a narrow `can / be
able' reading for \tdim{-thei}, but only with the
negation/irrealis-stack caveat kept explicit. Negative-modal strings
such as \tdim{khiathei ding om lo} remain overlap controls rather than
core evidence for this first slice.

\subsubsection{Deferred and overlap
material}\label{deferred-and-overlap-material}

The following items remain out of this first grammar slice:

\begin{itemize}
\tightlist
\item
  \tdim{pailai}, because \tdim{-lai} still overlaps lexical \tdim{lai} /
  \tdim{go-midst} material and is not yet a clean prospective anchor;
\item
  \tdim{dingin} and other clause-bound \tdim{-ding} material;
\item
  negative modal stacks such as \tdim{khiathei ding om lo};
\item
  sentence-final overlap such as \tdim{mangngilh ta hi};
\item
  directional/TAM stacking such as \tdim{khia-ta};
\item
  broader VP-slot stacks such as \tdim{bawlzoding} and
  \tdim{bawlsakthei};
\item
  report-summary items such as \tdim{-nawn} and \tdim{-khin} until
  cleaner anchors are selected.
\end{itemize}

These forms stay visible as overlap or deferred controls, but they do
not belong to the first compact TAM grammar slice.

\subsubsection{Recommended next step}\label{recommended-next-step-6}

After this grammar slice, the next step is the TAM dictionary print
slice drafted against the same \texttt{candidates\_tam.tsv} and
\texttt{dossier\_tam\_scope.md}; that slice now exists, and review-note
work has not yet begun. Broad TAM rewrite work remains out of scope.

\subsection{Directionals}\label{directionals}

\emph{Source slice:
\texttt{output/publication\_review/grammar\_directionals\_print\_slice.md}}

\subsubsection{Scope}\label{scope-5}

This is a short print-facing draft section on directionals in Tedim
Chin, controlled by \texttt{candidates\_directionals.tsv} and
\texttt{dossier\_directionals.md}.

It covers only a small candidate-backed set: outward \tdim{-khia},
anchored by \tdim{pokhia}; away \tdim{-khiat}, anchored by
\tdim{nawhkhiat}; nominalized \tdim{-khiat-na} boundary material,
represented by \tdim{hotkhiatna}; upward \tdim{-toh}, anchored by
\tdim{kilaktoh}; nominalized \tdim{-toh-na} boundary material,
represented by \tdim{kahtohna}; blocked comitative/accompany \tdim{-toh}
overlap, represented by \tdim{paitoh}; direction/side/manner \tdim{-lam}
boundary material, represented by \tdim{tawplam}; cautious toward
\tdim{-sawn} evidence, represented by \tdim{piasawn}; downward
\tdim{-suk} evidence, represented by \tdim{paisuk}; and deferred
\tdim{-lut}, \tdim{-phei}, \tdim{-cip}, plus \tdim{-tang}.

It does not yet attempt a full VP-slot chapter, a full TAM or aspect
account, a full inventory from raw suffix counts, or a treatment of all
lexicalized directional-looking forms. Dictionary and review-note slices
have not yet begun.

\subsubsection{Directionals in outline}\label{directionals-in-outline}

The current candidate-backed packet supports a narrow generalization.
Directionals are visible here as suffixed verbal or verbal-derived forms
in the candidate layer. The clearest first anchors are \tdim{pokhia},
\tdim{nawhkhiat}, \tdim{kilaktoh}, \tdim{piasawn}, and \tdim{paisuk}.
Nominalized forms such as \tdim{hotkhiatna} and \tdim{kahtohna} are
useful boundary evidence but are not identical to finite directional
verbs. \tdim{-toh} requires an explicit comitative/accompany warning
because \tdim{paitoh} is blocked as lexicalized \tdim{go-accompany}.
\tdim{-lam} remains direction/side/manner boundary material rather than
a clean simple verbal suffix in the current packet. \tdim{-lut},
\tdim{-phei}, \tdim{-cip}, and \tdim{-tang} remain deferred or not
print-ready.

\subsubsection{\texorpdfstring{Outward
\tdim{-khia}}{Outward }}\label{outward}

Genesis 2:5 supplies the cleanest current outward anchor:

\begin{exe}
\ex \label{ex:dir-khia-pokhia}
\gll \tdimword{po-khìa} \\
     grow-out \\
\glt 'grew'
\end{exe}

\tdim{Pokhia} is the cleanest current outward \tdim{-khia} anchor. It
supports a modest print claim that \tdim{-khia} can mark outward motion
or direction.

That claim must stay narrow. The row does \textbf{not} license raw
\tdim{khia} harvesting or a claim that every orthographic \tdim{khia}
sequence is directional evidence.

\subsubsection{\texorpdfstring{Away \tdim{-khiat} and nominalized
\tdim{-khiat-na}}{Away  and nominalized }}\label{away-and-nominalized}

Deuteronomy 9:4 supplies the current away anchor:

\begin{exe}
\ex \label{ex:dir-khiat-nawhkhiat}
\gll \tdimword{nawh-khìat} \\
     hurry-away \\
\glt 'cast them out'
\end{exe}

\tdim{Nawhkhiat} is the current compact away \tdim{-khiat} anchor. It is
accepted with caveat because the export still labels the selected
lemma/POS as \tdim{nawh} / \tdim{N}. The example is therefore usable as
construction-backed away evidence with an analyzer label caution, not as
an uncomplicated finite-verb showcase.

Future prose should keep \tdim{-khiat} construction-controlled rather
than treating every \tdim{...khiat} token as the same directional claim.

Exodus 14:13 then keeps nominalized boundary material visible:

\begin{exe}
\ex \label{ex:dir-khiatna-hotkhiatna}
\gll \tdimword{hot-khìat-ná} \\
     save-away-\textsc{nmlz} \\
\glt 'salvation'
\end{exe}

\tdim{Hotkhiatna} keeps \tdim{-khiat-na} visible in the packet. It is
nominalized boundary material and should not be treated as identical to
a finite directional predicate.

\subsubsection{\texorpdfstring{Upward \tdim{-toh} and
comitative/accompany
overlap}{Upward  and comitative/accompany overlap}}\label{upward-and-comitativeaccompany-overlap}

Numbers 9:17 supplies the current upward anchor:

\begin{exe}
\ex \label{ex:dir-toh-kilaktoh}
\gll \tdimword{kī-làk-tòh} \\
     \textsc{refl}-take-\textsc{up} \\
\glt 'was taken up'
\end{exe}

\tdim{Kilaktoh} is the current upward \tdim{-toh} anchor. It is usable
only with the packet's polysemy caveat: this row supports upward
\tdim{-toh}, but it does not license a raw equation of \tdim{-toh = UP}.

Blocked control:

\begin{quote}
\tdim{paitoh}
\end{quote}

\tdim{Paitoh} is blocked as comitative/accompany material. The analyzer
tests treat it as lexicalized \tdim{go-accompany}, so the present
grammar slice must not imply that every \tdim{-toh} token is
upward-directional evidence. Upward \tdim{-toh} examples should
therefore be paired with this explicit caveat.

Deuteronomy 32:50 also keeps nominalized \tdim{-toh-na} material
visible:

\begin{exe}
\ex \label{ex:dir-tohna-kahtohna}
\gll \tdimword{kàh-tòh-ná} \\
     climb-up-\textsc{nmlz} \\
\glt 'the mount whither thou goest up' / 'going up'
\end{exe}

\tdim{Kahtohna} keeps \tdim{-toh-na} visible, but it is nominalized
boundary material rather than a simple finite directional predicate.

\subsubsection{\texorpdfstring{Direction/side/manner
\tdim{-lam}}{Direction/side/manner }}\label{directionsidemanner}

Genesis 30:9 keeps \tdim{-lam} visible:

\begin{exe}
\ex \label{ex:dir-lam-tawplam}
\gll \tdimword{tawp-lam} \\
     end-\textsc{toward} \\
\glt 'left off' / 'at the end' (Genesis 30:9)
\end{exe}

\tdim{Tawplam} keeps \tdim{-lam} visible, but the export profile remains
nominal (\tdim{tawp} / \tdim{N}). The safest present analysis is
therefore direction/side/manner boundary material rather than a clean
verbal directional suffix.

This row should not be used to turn every \tdim{-lam} form into a simple
verbal directional.

\subsubsection{\texorpdfstring{Toward
\tdim{-sawn}}{Toward }}\label{toward}

Ezra 9:9 supplies the current cautious toward row:

\begin{exe}
\ex \label{ex:dir-sawn-piasawn}
\gll \tdimword{pía-sàwn} \\
     give-toward \\
\glt 'give us' / 'extend to us'
\end{exe}

\tdim{Piasawn} is the current cautious toward \tdim{-sawn} row. It is
more useful than kinship-heavy or nominal-looking rows for this packet,
but it still should stay construction-controlled. The current slice
should not generalize from lexicalized, continuative-looking, or
kinship-heavy \tdim{-sawn} material.

\subsubsection{\texorpdfstring{Downward
\tdim{-suk}}{Downward }}\label{downward}

Genesis 11:5 supplies the first corpus-backed downward row:

\begin{exe}
\ex \label{ex:dir-suk-paisuk}
\gll \tdimword{pái-sùk} \\
     go-\textsc{down} \\
\glt 'came down'
\end{exe}

\tdim{Paisuk} gives the packet a corpus-backed downward \tdim{-suk} row.
Analyzer tests also support \tdim{-suk}, but print prose should rely on
corpus-backed candidate rows rather than analyzer inventory alone.

\subsubsection{\texorpdfstring{Deferred forms: \tdim{-lut},
\tdim{-phei}, \tdim{-cip},
\tdim{-tang}}{Deferred forms: , , , }}\label{deferred-forms}

Several forms remain visible only as deferred or not print-ready
controls in this first slice:

\begin{itemize}
\tightlist
\item
  \tdim{uilut} keeps \tdim{-lut} visible but is not yet a clean
  print-safe inward anchor;
\item
  \tdim{paiphei} keeps \tdim{-phei} visible, but the current export
  gloss \tdim{go-enter} does not justify a clean horizontal claim;
\item
  \tdim{cip} remains lexical or analyzer-noise material rather than
  directional down evidence;
\item
  \tdim{tang} remains lexical or analyzer-noise material rather than
  endpoint-directional evidence.
\end{itemize}

These forms remain deferred or not print-ready in this slice. They
should not be promoted without cleaner analyzer-backed corpus rows.

\subsubsection{Editorial summary}\label{editorial-summary-5}

This slice safely supports seven modest claims: \tdim{-khia} is usable
as outward evidence, anchored by \tdim{pokhia}; \tdim{-khiat} is usable
as away evidence, anchored by \tdim{nawhkhiat}, with \tdim{-khiat-na}
boundary material from \tdim{hotkhiatna}; \tdim{-toh} is usable as
upward evidence, anchored by \tdim{kilaktoh}, but only with the
\tdim{paitoh} comitative/accompany caveat; \tdim{-toh-na} remains
visible through \tdim{kahtohna} as nominalized boundary material;
\tdim{-lam} remains direction/side/manner boundary material;
\tdim{-sawn} remains cautious toward evidence through \tdim{piasawn};
and \tdim{-suk} is now corpus-backed downward evidence through
\tdim{paisuk}.

What remains deferred is equally important: raw generated-report counts,
raw suffix harvesting, \tdim{paitoh} as upward \tdim{-toh}, nominalized
\tdim{-na} forms as equivalent to finite directional verbs, \tdim{-lut},
\tdim{-phei}, \tdim{-cip}, and \tdim{-tang} as print-ready directionals,
and broad TAM or VP-slot prose. Broad TAM, chrestomathy, Mizo/lus, and
other Kuki-Chin language work remain deferred.

The next step after this grammar slice is the dictionary print slice at
\texttt{output/publication\_review/dictionary\_directionals\_print\_slice.md}.
Dictionary and review-note work have not yet begun.

\subsection{Derivation / valency}\label{derivation-valency}

\emph{Source slice:
\texttt{output/publication\_review/grammar\_derivation\_valency\_print\_slice.md}}

\subsubsection{Editorial scope}\label{editorial-scope-7}

This is the first narrow derivation / valency grammar slice for Tedim.
It is controlled by
\texttt{output/publication\_review/candidates\_derivation\_valency.tsv}
and
\texttt{output/publication\_review/dossier\_derivation\_valency\_scope.md}.
Supporting/background evidence comes from
\texttt{docs/grammar/reports/05-verb-08-derivational.md},
\texttt{docs/grammar/reports/05-verb-09-valency.md},
\texttt{docs/grammar/morphemes/06-derivational.md},
\texttt{docs/grammar/lit-reviews/05-verb-09-valency-lit.md}, and the
regression evidence in \texttt{tests/test\_sak\_caus\_benf.py}.

This is not a full derivation chapter, not a full valency chapter, and
not a full verbal morphology chapter. It also does not reopen adjacent
packet domains already controlled through
\texttt{output/publication\_review/review\_notes\_vp\_structure\_stacking.md},
\texttt{output/publication\_review/review\_notes\_tam.md},
\texttt{output/publication\_review/review\_notes\_directionals.md},
\texttt{output/publication\_review/review\_notes\_pronouns.md},
\texttt{tests/test\_vp\_slots.py}, and
\texttt{tests/test\_prefix\_agr\_poss.py}.

The present slice therefore covers only the first safe \tdim{-sak}
claim: \tdim{paisak} as the clearest causative anchor and \tdim{muhsak}
as the clearest benefactive or applicative-like split row. No dictionary
slice exists yet for derivation/valency, because the packet still leaves
the \tdim{-sak} lexical split open for later editorial review.

\subsubsection{\texorpdfstring{Causative
\tdim{-sak}}{Causative }}\label{causative}

\tdim{paisak} is the safest causative anchor for the first print-facing
derivation / valency claim. The candidate layer marks it as the main
future \tdim{-sak} causative anchor, and
\texttt{docs/grammar/reports/05-verb-09-valency.md} already lists it
among the common \tdim{-sak} forms.

The safe grammar claim is deliberately narrow. The current evidence
supports a productive Form I plus \tdim{-sak} causative pattern, with
\tdim{paisak} as the candidate-controlled row that shows the pattern
most cleanly. \texttt{tests/test\_sak\_caus\_benf.py} protects that
contrast explicitly by requiring \tdim{paisak} to gloss as
\tdim{go-CAUS} rather than as a benefactive or lexicalized exception.

That is enough for the first core claim. The slice does not need to
generalize over every \tdim{-sak} form in the reports before saying that
a productive causative use is visible.

\subsubsection{\texorpdfstring{Benefactive / applicative-like
\tdim{-sak}}{Benefactive / applicative-like }}\label{benefactive-applicative-like}

\tdim{muhsak} is the safest benefactive or applicative-like split row in
the current packet. The candidate layer keeps it distinct from the plain
causative line, and \texttt{tests/test\_sak\_caus\_benf.py} protects
that distinction by requiring Form II plus \tdim{-sak} rows to keep the
\tdim{.II} marker and a \tdim{BENF} gloss.

The print claim here also stays narrow. The current evidence supports
keeping Form II plus \tdim{-sak} distinct from the plain causative line
in the candidate layer, with \tdim{muhsak} as the clearest first-row
anchor. The literature and morpheme files justify describing this as
benefactive or applicative-like material, but the slice should not
pretend that every higher-level label choice is already settled.

This is why \tdim{muhsak} belongs in the first slice even though it is
more caveated than \tdim{paisak}. It is the clearest compact row showing
that the \tdim{-sak} domain is not exhausted by a simple causative
paraphrase.

\subsubsection{\texorpdfstring{Editorial treatment of the \tdim{-sak}
split}{Editorial treatment of the  split}}\label{editorial-treatment-of-the-split}

The grammar can therefore treat the \tdim{paisak} versus \tdim{muhsak}
contrast as a controlled split in the first print slice. The evidence is
strong enough to keep Form I plus \tdim{-sak} and Form II plus
\tdim{-sak} apart in the editorial layer.

At the same time, the slice should keep open the higher-level question
of whether this contrast is best described as two readings of one suffix
or two editorial subsections of the same suffixal domain.
\texttt{docs/grammar/lit-reviews/05-verb-09-valency-lit.md} and
\texttt{docs/grammar/morphemes/06-derivational.md} both support cautious
wording here: the split is real enough for candidate control, but the
final theoretical framing should stay smaller than a full chapter claim.

This first slice therefore adopts a practical editorial solution rather
than a maximal theoretical one. It prints the causative and
benefactive/applicative-like uses as separate controlled subsections
while leaving open whether a later chapter should collapse them back
into one suffix with two readings.

\subsubsection{Boundary material}\label{boundary-material-4}

The rest of the candidate packet remains outside the first grammar slice
because each row is still dominated by another unresolved boundary.

\tdim{paipih} stays outside the first grammar slice because \tdim{-pih}
still carries applicative, comitative, associative, and benefactive
uncertainty. \tdim{mipihte} stays outside because it is nominal or
lexicalized \tdim{pih} boundary material rather than a clean verbal
anchor.

\tdim{kisep} and \tdim{kigen} stay outside because \tdim{ki-} still
needs a separate reflexive, middle, and passive-like treatment, and
because the packet must keep prefix/agreement boundary control explicit
through \texttt{output/publication\_review/review\_notes\_pronouns.md}
and \texttt{tests/test\_prefix\_agr\_poss.py}.

\tdim{ciahsakkik}, \tdim{bawlsakthei}, and \tdim{paikhiatsak} stay
outside because they are derivation-heavy stacks interacting
respectively with aspect, modal, and directional material. Those rows
remain visible through
\texttt{output/publication\_review/review\_notes\_vp\_structure\_stacking.md},
\texttt{output/publication\_review/review\_notes\_tam.md},
\texttt{output/publication\_review/review\_notes\_directionals.md}, and
\texttt{tests/test\_vp\_slots.py}, but they are not the first core
derivation anchor.

\tdim{piangsak} stays outside because it is lexicalized or
transitivity-adjacent rather than a clean productive \tdim{-sak} anchor.

\subsubsection{Safe first-slice claim}\label{safe-first-slice-claim-5}

At the current slice maturity level, the safest derivation / valency
claim is that Tedim has candidate-controlled evidence for a productive
\tdim{-sak} domain, with \tdim{paisak} supporting a causative use and
\tdim{muhsak} supporting a distinct benefactive or applicative-like use.

That claim is deliberately smaller than a full derivation chapter,
smaller than a full valency chapter, and smaller than a full verbal
morphology chapter. It does not settle the whole \tdim{-pih} system, the
whole \tdim{ki-} system, derivation-heavy stacking, or transitivity as a
separate domain.

\subsubsection{Recommended next step}\label{recommended-next-step-7}

After this grammar slice, the next step should be review notes rather
than a dictionary layer, because a dictionary layer would risk
overclaiming beyond the current candidate-controlled evidence before the
\tdim{-sak} lexical split is reviewed.

\subsection{Nominalization}\label{nominalization}

\emph{Source slice:
\texttt{output/publication\_review/grammar\_nominalization\_print\_slice.md}}

\subsubsection{Editorial scope}\label{editorial-scope-8}

This is the first narrow nominalization grammar slice. It is controlled
by \texttt{output/publication\_review/candidates\_nominalization.tsv}
and
\texttt{output/publication\_review/dossier\_nominalization\_scope.md}.
Supporting/background evidence comes from
\texttt{docs/grammar/reports/07-nmlz-01-deverbal.md},
\texttt{docs/grammar/morphemes/06-derivational.md},
\texttt{docs/grammar/grammar\_source\_map.json}, and
\texttt{docs/SKELETON\_GRAMMAR.md}.

This is not a full nominalization chapter, not a full derivation
chapter, not a full relative-clause chapter, and not a full case-routing
chapter. It also stays narrow against
\texttt{output/publication\_review/candidates\_clause\_linkage.tsv},
\texttt{output/publication\_review/review\_notes\_clause\_linkage.md},
\texttt{output/publication\_review/review\_notes\_case\_marking.md},
\texttt{output/publication\_review/review\_notes\_derivation\_valency.md},
\texttt{output/publication\_review/review\_notes\_prefix\_agreement.md},
and \texttt{output/publication\_review/review\_notes\_pronouns.md}.

The present slice therefore covers only the first safe nominalization
claim. \tdim{-na} is the clearest productive deverbal nominalizer, with
\tdim{bawlna} as the controlled anchor form, \tdim{bawl-na} as the
segmentation, and \texttt{make-NMLZ\ /\ making,\ creation} as the
controlled gloss and function. No dictionary slice exists yet for
nominalization, because this packet is still establishing a controlled
constructional or morphological claim rather than a settled lexical
layer. The packet now properly proceeds through nominalization review
notes rather than a dictionary slice.

\subsubsection{\texorpdfstring{Deverbal nominalization with
\tdim{-na}}{Deverbal nominalization with }}\label{deverbal-nominalization-with}

\tdim{-na} is the safest current nominalization anchor in the packet.
The candidate TSV marks it as the main future print-facing row, and the
report, source map, and skeleton all converge on it as the primary
productive deverbal nominalizer.

The grammar claim here is deliberately limited. At the current slice
maturity level, the safe print-facing row is \tdim{bawlna}, with the
segmentation \tdim{bawl-na} and the controlled gloss
\texttt{make-NMLZ\ /\ making,\ creation}. That is enough to support a
narrow claim for productive deverbal action or result nominalization
without broadening into a full nominalization chapter.

\subsubsection{\texorpdfstring{Why \tdim{-pa} and \tdim{-mi} are not yet
the first
slice}{Why  and  are not yet the first slice}}\label{why-and-are-not-yet-the-first-slice}

Agentive or person-head nominalization evidence also exists in the
candidate layer, especially through \tdim{bawlpa} and
\tdim{hong pai mi}.

Those rows stay outside the first print-facing claim because \tdim{-pa}
has lexicalized or title-like boundary rows such as \tdim{kumpipa} and
\tdim{Topa}, while \tdim{-mi} still overlaps with person-head and
relative-clause analysis. The packet is stronger if it keeps \tdim{-pa}
and \tdim{-mi} visible but secondary rather than forcing them into the
first nominalization slice before those lexical and clausal boundaries
are better settled.

\subsubsection{Why nominalized relatives are not yet the first
slice}\label{why-nominalized-relatives-are-not-yet-the-first-slice}

Nominalized relative or clause-derived nominal evidence also exists in
the candidate layer, especially through \tdim{omna}.

That row stays outside the first print-facing claim because it still
belongs partly to clause linkage and relative-clause analysis. The
current packet can safely preserve \tdim{omna} as candidate-layer
evidence without pretending that the first nominalization slice should
already resolve the nominalization versus relative-clause boundary.

\subsubsection{Why nominalization plus case is not yet the first
slice}\label{why-nominalization-plus-case-is-not-yet-the-first-slice}

Nominalization-plus-case evidence also exists in the boundary layer,
especially through \tdim{muhna-ah}.

That row stays outside the first print-facing claim because it belongs
partly to case marking and clause-linkage boundaries. It is important
evidence for later packet coordination, but it should not lead the first
nominalization slice.

\subsubsection{Boundary material}\label{boundary-material-5}

The rest of the nominalization packet stays outside the first grammar
slice because each row is still dominated by another unresolved
boundary.

\tdim{bawlpa} stays outside because agentive \tdim{-pa} still has to be
separated from lexicalized and title-like rows.

\tdim{hong pai mi} stays outside because person-head \tdim{-mi} still
overlaps with relative-clause analysis.

\tdim{omna} stays outside because nominalized relatives remain shared
with clause linkage.

\tdim{muhna-ah} stays outside because nominalization plus case routing
remains shared with case marking.

\tdim{kumpipa} and \tdim{Topa} stay outside because lexicalized or
title-like \tdim{-pa} rows should not be mistaken for the cleanest
productive anchor.

\tdim{a bawl mi} and similar human-head relative rows stay outside
because they still sit between nominalization, relative clauses, and
prefix/agreement questions.

bare \tdim{na} stays outside because it is an analyzer-noisy surface
form.

report-only counts stay outside because attestation by itself does not
make a row safe for the first print-facing claim.

Any broad nominalization chapter claim stays outside because this packet
is not yet ready to generalize from one safe deverbal nominalizer to the
whole nominalization system.

\subsubsection{Safe first-slice claim}\label{safe-first-slice-claim-6}

At the current slice maturity level, the safest nominalization claim is
that Tedim has candidate-controlled evidence for productive deverbal
nominalization with \tdim{-na}, with \tdim{bawlna} / \tdim{bawl-na} as
the clearest current anchor. Agentive \tdim{-pa}, person-head
\tdim{-mi}, nominalized relatives, and nominalization-plus-case rows
remain candidate-layer or boundary material.

That claim is deliberately smaller than a full nominalization chapter,
smaller than a full derivation chapter, smaller than a full
relative-clause chapter, and smaller than a full case-routing chapter.

\subsubsection{Recommended next step}\label{recommended-next-step-8}

This packet now properly proceeds to nominalization review notes rather
than to a dictionary slice, because this is a constructional or
morphological packet and the lexical treatment of \tdim{-pa},
\tdim{-mi}, and lexicalized forms remains unsettled.

If the project later wants more nominalization work after review notes,
the next sub-scope should be a separate \tdim{-pa} or \tdim{-mi}
agentive candidate expansion rather than a dictionary layer, but not in
this commit.

\subsection{Clause linkage}\label{clause-linkage}

\emph{Source slice:
\texttt{output/publication\_review/grammar\_clause\_linkage\_print\_slice.md}}

\subsubsection{Editorial scope}\label{editorial-scope-9}

This is the first narrow clause-linkage grammar slice for Tedim. It is
controlled by
\texttt{output/publication\_review/candidates\_clause\_linkage.tsv} and
\texttt{output/publication\_review/dossier\_clause\_linkage\_scope.md}.
Supporting/background evidence comes from
\texttt{docs/grammar/reports/08-clause-01-subordination.md},
\texttt{docs/grammar/reports/08-clause-02-switch-reference.md},
\texttt{docs/grammar/reports/08-clause-03-relatives.md}, and
\texttt{docs/grammar/lit-reviews/08-clause-03-subordination-lit.md}.

This is not a full complex-sentence chapter, not a full switch-reference
chapter, and not a full relative-clause chapter. It also stays narrow
against
\texttt{output/publication\_review/review\_notes\_sentence\_final\_particles.md},
\texttt{output/publication\_review/review\_notes\_tam.md},
\texttt{output/publication\_review/review\_notes\_vp\_structure\_stacking.md},
\texttt{output/publication\_review/review\_notes\_prefix\_agreement.md},
and \texttt{output/publication\_review/review\_notes\_pronouns.md}.

The present slice therefore covers only the first safe subordination
claim. \tdim{Ciangin} is the clearest temporal subordination anchor,
while \tdim{dingin} remains visible only as a caveated purposive or
clause-bound irrealis overlap row. No dictionary slice exists yet for
clause linkage, because this packet is still establishing a controlled
clausal claim rather than a lexical headword layer. The packet now
properly proceeds through clause-linkage review notes rather than a
dictionary slice.

\subsubsection{Temporal subordination:
ciangin}\label{temporal-subordination-ciangin}

\tdim{Ciangin} is the safest current subordination anchor in the packet.
The candidate TSV marks it as the main future print-facing row, and the
reports plus literature review converge on it as the clearest temporal
subordinator.

The grammar claim here is deliberately limited. At the current slice
maturity level, the safe print-facing row is \tdim{tua ciangin}, with
the segmentation \tdim{ciang-in} kept visible where useful. That is
enough to support a narrow temporal subordination claim without
broadening into a full complex-sentence chapter or a full inventory of
all clause-linkage constructions.

\subsubsection{Clause-bound purposive / irrealis overlap:
dingin}\label{clause-bound-purposive-irrealis-overlap-dingin}

\tdim{Dingin} is relevant to clause linkage because the reports support
purposive or clause-bound irrealis use, but it remains partly shared
with TAM. The candidate layer is therefore right to keep it explicit
while refusing to let it lead the first slice.

The grammar claim stays cautious on purpose. At the current slice
maturity level, \tdim{dingin} is visible only as a caveated overlap row:
it helps show that clause linkage interacts with purposive or irrealis
marking, but it does not yet justify a broader clause-linkage or TAM
rewrite.

\subsubsection{Why switch reference is not yet the first
slice}\label{why-switch-reference-is-not-yet-the-first-slice}

Switch-reference evidence exists in the candidate layer, especially
through \tdim{VERB-in} with \tdim{ngenin} as the anchor example and
through \tdim{ahih ciangin} as the clearest different-subject
construction.

Those rows stay outside the first print-facing claim because the
converb, subordinator, and switch-reference analysis remains
theoretically unsettled. The current packet can safely preserve them as
candidate-layer evidence without pretending that the first
clause-linkage slice should already be a switch-reference chapter.

\subsubsection{Why relative clauses are not yet the first
slice}\label{why-relative-clauses-are-not-yet-the-first-slice}

Relative-clause evidence also exists in the candidate layer, especially
through \tdim{a bawl mi} and \tdim{omna}.

Those rows stay outside the first print-facing claim because relative
clauses still interact with prefix/agreement, nominalization, and
case-routing questions. The packet is stronger if it keeps relative
clauses visible but secondary rather than forcing them into the first
clause-linkage slice before those boundaries are better settled.

\subsubsection{Boundary material}\label{boundary-material-6}

The rest of the clause-linkage packet stays outside the first grammar
slice because each row is still dominated by another unresolved
boundary.

\tdim{VERB-in} and \tdim{ngenin} stay outside because same-subject
chaining still sits on the boundary between converbial chaining,
subordination, and switch-reference analysis.

\tdim{ahih ciangin} stays outside because different-subject linkage
still overlaps with temporal subordination and pronominal boundary
questions.

\tdim{a bawl mi} and \tdim{omna} stay outside because relative clauses
still interact with prefix/agreement and nominalization.

\tdim{muhna-ah} stays outside because nominalized relative-plus-case
material belongs to a later nominalization and case-routing treatment.

\tdim{leh} stays outside because sentence-final-particle, coordinator,
and subordinator analyses still overlap there.

\tdim{hangin} and \tdim{bangin} stay outside because they remain broader
report-inventory rows rather than the first safe print anchor.

report-only relative-clause counts involving \tdim{a-} also stay outside
because they do not yet separate relativizer behavior cleanly from
broader prefix/agreement material.

Any broad complex-sentence chapter claim stays outside because this
packet is not yet ready to generalize from one safe subordination anchor
to the whole clause-linkage system.

\subsubsection{Safe first-slice claim}\label{safe-first-slice-claim-7}

At the current slice maturity level, the safest clause-linkage claim is
that Tedim has candidate-controlled evidence for temporal subordination,
with \tdim{ciangin} as the clearest current anchor. \tdim{Dingin} is
visible as a caveated purposive or clause-bound irrealis overlap row,
but switch reference and relative clauses remain candidate-layer
material.

That claim is deliberately smaller than a full complex-sentence chapter,
smaller than a full switch-reference chapter, and smaller than a full
relative-clause chapter.

\subsubsection{Recommended next step}\label{recommended-next-step-9}

This packet now properly proceeds to clause-linkage review notes rather
than to a dictionary slice, because it is a constructional or clausal
packet rather than a lexical-headword packet.

If the project later wants one more clause-linkage step before review
notes and human review, the next sub-scope should be a separate
switch-reference or relative-clause candidate expansion rather than a
dictionary layer, but not in this commit.

\reviewchapterbreak

\section{Clause type, discourse-facing material, and expressive
morphology}\label{clause-type-discourse-facing-material-and-expressive-morphology}

\subsection{Negation}\label{negation}

\emph{Source slice:
\texttt{output/publication\_review/grammar\_negation\_print\_slice.md}}

\subsubsection{Scope}\label{scope-6}

This account offers a short treatment of Tedim negation. It is
intentionally narrow. It focuses on a small set of manually checked
Bible examples and on the interaction between \tdim{lo}, \tdim{loh}, and
\tdim{kei}. It does not attempt a full TAM chapter, a treatment of
directionals, or a complete account of every negative construction in
the corpus.

\subsubsection{Overview of the negation
system}\label{overview-of-the-negation-system}

The literature and the Bible corpus agree on one central point: Tedim
negation cannot be reduced to \tdim{lo} alone. Henderson foregrounds
\tdim{-kei}, including prohibitive uses, while Zam Ngaih Cing explicitly
describes a contrast between \tdim{-kei} and \texttt{-lou/-louh} and
links the latter to stem-sensitive or dependent environments
\citep{henderson1965, zamngaihcing2017}. The safest summary is therefore
not ``Tedim negation = \tdim{lo}'', but rather that the present corpus
supports a three-part treatment: ordinary clause-level negation with
\tdim{lo}, dependent or derived negation with \tdim{loh}, and \tdim{kei}
as a central negator in prohibitives and many irrealis-heavy or
directive contexts.

\subsubsection{\texorpdfstring{Clause-level negation with
\tdim{lo}}{Clause-level negation with }}\label{clause-level-negation-with}

\tdim{Lo} is the safest place to start because it is the clearest
ordinary clause-level negator in the Bible corpus. It appears in
straightforward negative predicates such as \tdim{thusim lo hi},
\tdim{nei lo hi}, \tdim{om lo hi}, and in many future or modal
combinations such as \tdim{V lo ding}. That does not make it the whole
negation system, but it does make it the most transparent starting
point.

\begin{exe}
\ex \label{ex:neg-lo}
\gll \tdimword{àhìh} \tdimword{hàng-ìn} \tdimword{Kain} \tdimword{lè} \tdimword{àmà} \tdimword{pìak-ná} \tdimword{thusim} \tdimword{lò} \tdimword{hì.} \tdimword{Túa} \tdimword{àhìh} \tdimword{cìang-ìn} \tdimword{Kain} \tdimword{hèh} \tdimword{màhmàh} \tdimword{à,} \tdimword{à} \tdimword{mái} \tdimword{sía} \tdimword{hì.} \\
     be.\textsc{3sg}.\textsc{rel} because-\textsc{erg} \textsc{kain} and \textsc{3sg}.\textsc{poss} give.to-\textsc{nmlz} parable \textsc{neg} \textsc{decl} \textsc{dist} be.\textsc{3sg}.\textsc{rel} then-\textsc{erg} \textsc{kain} anger very \textsc{3sg} \textsc{3sg} face evil \textsc{decl} \\
\glt 'but unto Cain and to his offering he had not respect. And Cain was very wroth, and his countenance fell.' (Genesis 4:5)
\end{exe}

Genesis 4:5 is a good central example because it shows exactly the
clause type a reader expects from a first negation section: a verbal
predicate followed by \tdim{lo} and the clause-final declarative
material. Genesis 11:30 \tdim{ta nei lo hi} is equally useful in a more
compact predicative clause, and existential patterns such as
\tdim{om lo hi} belong in the same zone.

\tdim{V lo uh} is a real and common surface string, but it is not
automatically prohibitive. Genesis 2:25 \tdim{maizum lo uh hi} means
``they were not ashamed'' and is simply an ordinary plural negative
predicate. Similar plural negatives occur elsewhere with no imperative
force. It is better treated as an ordinary clausal negative pattern than
as the defining prohibitive construction.

\subsubsection{\texorpdfstring{Dependent and derived negation with
\tdim{loh}}{Dependent and derived negation with }}\label{dependent-and-derived-negation-with}

\tdim{Loh} should not be treated as a random spelling variant of
\tdim{lo}. In the current Bible corpus it clusters in dependent,
purposive, and derived environments such as \tdim{... loh dinga},
\tdim{... loh dingin}, and \tdim{... loh nadingin}. The safest label is
therefore something like ``dependent or derived negative form'', not
just ``another way to spell \tdim{lo}''.

\begin{exe}
\ex \label{ex:neg-loh}
\gll \tdimword{Àmàh} \tdimword{ìn,} \tdimword{nà} \tdimword{guak-tāng-à} \tdimword{nà} \tdimword{ōm-ná} \tdimword{kua} \tdimword{ìn} \tdimword{hòng} \tdimword{gēn} \tdimword{àhì} \tdimword{hìam?} \tdimword{Nà} \tdimword{nèk} \tdimword{lòh} \tdimword{dìng-à} \tdimword{kòng} \tdimword{thupiak} \tdimword{singgah} \tdimword{né} \tdimword{khá} \tdimword{nà} \tdimword{hì} \tdimword{hìam} \tdimword{cí} \tdimword{hì.} \\
     \textsc{3sg}.\textsc{pro} \textsc{erg} \textsc{2sg} bare-body-\textsc{loc} \textsc{2sg} exist-\textsc{nmlz} who \textsc{erg} \textsc{3→1} speak be.\textsc{3sg} \textsc{q} \textsc{2sg} eat.\textsc{ii} \textsc{neg}-\textsc{nom} \textsc{irr}-\textsc{loc} \textsc{1sg→3} commandment hang eat \textsc{neg}.\textsc{perf} \textsc{2sg} \textsc{decl} \textsc{q} say \textsc{decl} \\
\glt 'And he said, Who told thee that thou wast naked? Hast thou eaten of the tree, whereof I commanded thee that thou shouldest not eat?' (Genesis 3:11)
\end{exe}

Genesis 3:11 is especially useful because \tdim{loh} is not just
attached to a bare verb in isolation. It appears inside a dependent
complex, \tdim{na nek loh dinga kong thupiak}, where the whole point is
the commanded non-occurrence of the event. Isaiah 6:10 strengthens the
same pattern by repeating \tdim{muh loh nading}, \tdim{zak loh nading},
and \tdim{theihtel loh nading} inside clearly purposive or
clause-linking material. This is exactly the distribution seen again and
again in the corpus: \tdim{loh} is where the Bible orthography most
visibly tracks non-finite, dependent, and derived negation.

\subsubsection{\texorpdfstring{\tdim{kei} in prohibitives and
irrealis-heavy
negation}{ in prohibitives and irrealis-heavy negation}}\label{in-prohibitives-and-irrealis-heavy-negation}

\tdim{Kei} is central to real prohibitives in the Bible corpus, and
those prohibitives are among the cleanest manually checked negative
examples available. This alone is enough to show why the negation system
cannot be flattened to \tdim{lo}.

\begin{exe}
\ex \label{ex:neg-kei}
\gll \tdimword{Àmàh} \tdimword{ìn,} \tdimword{tangval-pà} \tdimword{sù} \tdimword{kèi} \tdimword{ìn.} \tdimword{Àmà} \tdimword{tùngàh} \tdimword{bangmah} \tdimword{híh} \tdimword{kèi} \tdimword{ìn.} \tdimword{Bàng} \tdimword{hàng} \tdimword{hìam} \tdimword{cìh} \tdimword{lèh} \tdimword{nàng} \tdimword{ìn} \tdimword{Pasian} \tdimword{nà} \tdimword{zah-tàk-ná} \tdimword{tù-ìn} \tdimword{kà} \tdimword{théi} \tdimword{à,} \tdimword{nà} \tdimword{tá-pà} \tdimword{khàt} \tdimword{nèih-sun} \tdimword{kèima} \tdimword{à} \tdimword{dìng-ìn} \tdimword{nà} \tdimword{hum-cìp} \tdimword{lòh-ná} \tdimword{kà} \tdimword{mú} \tdimword{hì} \tdimword{à} \tdimword{cí} \tdimword{hì.} \\
     \textsc{3sg}.\textsc{pro} \textsc{erg} \textsc{tangvalpa} tear \textsc{neg} \textsc{erg} \textsc{3sg}.\textsc{poss} on-\textsc{loc} nothing \textsc{prox} \textsc{neg} \textsc{erg} like reason \textsc{q} say.\textsc{nom} and \textsc{2sg}.\textsc{pro} \textsc{erg} God \textsc{2sg} fear-exact-\textsc{nmlz} now-\textsc{erg} \textsc{1sg} know \textsc{3sg} \textsc{2sg} child-male one have.\textsc{ii}-long \textsc{1sg}.self \textsc{3sg} \textsc{irr}-\textsc{erg} \textsc{2sg} fear-grip \textsc{neg}-\textsc{nmlz} \textsc{1sg} see.\textsc{i} \textsc{decl} \textsc{3sg} say \textsc{decl} \\
\glt 'And he said, Lay not thine hand upon the lad, neither do thou any thing unto him: for now I know that thou fearest God, seeing thou hast not withheld thy son, thine only son from me.' (Genesis 22:12)
\end{exe}

Genesis 22:12 gives two clear negative imperatives in the same speech
turn. Genesis 15:1 \tdim{Lau kei in}, Genesis 19:17
\tdim{Nunghei kei unla, ... khawl kei un}, Leviticus 10:9
\tdim{ne kei un}, and Numbers 14:42 \tdim{kuanto kei un ... do kei un}
all point in the same direction: \tdim{kei} is the core prohibitive
negator in the current Bible data.

The corpus also shows that \tdim{kei} is not limited to prohibitives or
to first-person realis. Exodus 5:2 has both \tdim{ka thei kei hi} and
\tdim{ka paisak kei ding hi}, which shows \tdim{kei} in ordinary
negation and in irrealis-heavy future material. The safest print summary
is therefore that \tdim{kei} is central in prohibitives and common in
directive, quoted, or otherwise irrealis-heavy negation, without forcing
the whole system into a rigid person-based rule.

\subsubsection{Negative existence and cessative
negation}\label{negative-existence-and-cessative-negation}

Negative existence belongs with ordinary \tdim{lo}-based negation, but
it deserves its own short section because it is extremely common and
very readable in Bible prose. Clauses such as Genesis 11:30
\tdim{ta nei lo hi}, Genesis 37:24 \tdim{a sungah tui om lo hi}, and
Genesis 39:9 \tdim{a kep tuam bangmah om lo hi} show the ordinary Tedim
way of saying that something is absent, unavailable, or not possessed.

The corpus also supports a neat constructional entry for \tdim{nawn lo}
``no longer, not again''. Genesis 8:12 \tdim{hong ciahkik nawn lo hi} is
a clean cessative example, and Genesis 17:5
\tdim{Na min Abram hi nawn lo ding a} shows the same construction in a
renamed-future environment. This material is better treated as a
construction than as a new basic negator.

\subsubsection{Ability and inability}\label{ability-and-inability}

The clearest treatment here is not a bare negator but the constructional
pair \tdim{thei lo / theih loh}. The Bible shows common clause-level
inability with \tdim{thei lo}, while dependent or derived negative
ability is better represented by \tdim{theih loh} patterns than by the
much rarer exact string \tdim{theih lo}.

\begin{exe}
\ex \label{ex:neg-thei-lo}
\gll \tdimword{À} \tdimword{khùt-tè} \tdimword{ìn} \tdimword{à} \tdimword{sanggampa} \tdimword{Esau'} \tdimword{khùt} \tdimword{bàng-ìn} \tdimword{múl} \tdimword{néi} \tdimword{àhìh} \tdimword{màn-ìn} \tdimword{àmàh} \tdimword{ìn} \tdimword{théi} \tdimword{lò} \tdimword{hì.} \tdimword{Túa} \tdimword{àhìh} \tdimword{cìang-ìn} \tdimword{àmàh} \tdimword{ìn} \tdimword{thù-phà} \tdimword{pía} \tdimword{dìng-ìn} \tdimword{kithawi} \tdimword{à,} \\
     \textsc{3sg} hand-\textsc{pl} \textsc{erg} \textsc{3sg} brother \textsc{esau}.\textsc{poss} hand like-\textsc{erg} hair have be.\textsc{3sg}.\textsc{rel} reason-\textsc{erg} \textsc{3sg}.\textsc{pro} \textsc{erg} know \textsc{neg} \textsc{decl} \textsc{dist} be.\textsc{3sg}.\textsc{rel} then-\textsc{erg} \textsc{3sg}.\textsc{pro} \textsc{erg} word-good give \textsc{irr}-\textsc{erg} hot.displeasure \textsc{3sg} \\
\glt 'And he discerned him not, because his hands were hairy, as his brother Esau's hands: so he blessed him.' (Genesis 27:23)
\end{exe}

Genesis 27:23 is a compact ordinary example: \tdim{thei lo hi} means
that recognition or knowledge failed to occur. Genesis 37:4
\tdim{amah hopih thei lo uh hi} shows the same pattern with an ability
reading, while Exodus 33:20 \tdim{na mu thei kei ding hi} shows
inability under \tdim{kei}. On the dependent side, Exodus 10:5
\tdim{muh theih loh nadingin} and Isaiah 6:10 \tdim{theihtel loh nading}
show why \tdim{theih loh} must be represented directly rather than
treating exact \tdim{theih lo} as the dominant form in the current
corpus.

\subsubsection{Negative polarity items}\label{negative-polarity-items}

The Bible corpus gives clean negative-polarity material, but only after
filtering. Raw exact-string counts overgenerate because \tdim{kuamah}
can occur in non-pronominal strings and \tdim{bangmah} has non-NPI uses
such as \tdim{tua bangmah hi-in} ``likewise''. The discussion here
therefore uses only manually checked examples.

\begin{exe}
\ex \label{ex:neg-kuamah}
\gll \tdimword{Àmàh} \tdimword{khùa-dāk} \tdimword{kawi-kawi} \tdimword{à,} \tdimword{kuamah} \tdimword{mú} \tdimword{lò} \tdimword{àhìh} \tdimword{màn-ìn} \tdimword{Egypt} \tdimword{mipa} \tdimword{thàt-ìn} \tdimword{sehnel} \tdimword{sùng-àh} \tdimword{à} \tdimword{seel} \tdimword{hì.} \\
     \textsc{3sg}.\textsc{pro} town-near crooked\textasciitilde{}REDUP \textsc{3sg} nobody see.\textsc{i} \textsc{neg} be.\textsc{3sg}.\textsc{rel} reason-\textsc{erg} \textsc{egypt} man kill-\textsc{cvb} wilderness inside-\textsc{loc} \textsc{3sg} hide \textsc{decl} \\
\glt 'And he looked this way and that way, and when he saw that there was no man, he slew the Egyptian, and hid him in the sand.' (Exodus 2:12)
\end{exe}

Exodus 2:12 is a good \tdim{kuamah} example because the NPI is clearly
licensed by the negative predicate. \tdim{Bangmah} behaves similarly in
real NPI contexts such as Genesis 22:12
\tdim{Ama tungah bangmah hih kei in} and Genesis 39:9
\tdim{bangmah om lo hi}. Those are the kinds of examples the dictionary
can safely quote. What should stay out of the main prose is any claim
based on raw counts alone.

\subsubsection{Summary}\label{summary}

The current evidence is strong enough for a stable negation chapter, but
only if it keeps the system narrow and honest. \tdim{Lo} is the safest
starting point for ordinary clause-level negation. \tdim{Loh} is a real
dependent or derived negative form and should not be collapsed into
accidental spelling variation. \tdim{Kei} is central in prohibitives and
in many irrealis-heavy negatives, so the chapter must represent it
directly rather than relegating it to a footnote.

Just as important, the description must avoid false simplicity.
\tdim{V lo uh} is not the prohibitive construction, Genesis 2:25 is not
a prohibitive example, and the NPI section must use filtered examples
rather than raw exact-string totals. The result is a compact but
linguistically honest account of the current Bible evidence.

\subsection{Interrogatives}\label{interrogatives}

\emph{Source slice:
\texttt{output/publication\_review/grammar\_interrogatives\_print\_slice.md}}

\subsubsection{Scope}\label{scope-7}

This short print-facing draft section on interrogatives in Tedim Chin is
controlled by \texttt{candidates\_interrogatives.tsv} and
\texttt{dossier\_interrogatives.md}. It covers clause-final \tdim{hiam}
and selected WH + \tdim{hiam} content questions. It does not yet attempt
a full treatment of the sentence-final particle system, and it leaves
comparison particles such as \tdim{maw}, \tdim{ham}, and \tdim{em} for
later review.

\subsubsection{Interrogatives in
outline}\label{interrogatives-in-outline}

The current candidate-backed packet supports a narrow generalization.
Yes/no questions can be marked by clause-final \tdim{hiam}, and content
questions use a WH element plus \tdim{hiam}. The present WH inventory
represented in the packet is \tdim{bang}, \tdim{kua}, \tdim{bangci}, and
\tdim{banghangin}. Embedded question complements are visible in the
dossier, but they are not yet ready for the first printed analysis.

Because this slice is controlled by the candidate-first packet rather
than by raw report counts, it should stay narrow. The point here is not
to describe every sentence-final question marker or every surface
occurrence of \tdim{hiam}, but to print only the examples that the
current candidate and dossier layer already support.

\subsubsection{\texorpdfstring{Clause-final
\tdim{hiam}}{Clause-final }}\label{clause-final}

The current candidate-backed examples support clause-final \tdim{hiam}
as the core yes/no pattern represented in this packet. The best print
anchor is Genesis 24:23, using the attested analyzer-backed clause
rather than the older generated-report paraphrase.

\begin{exe}
\ex \label{ex:hiam-yes-no}
\gll \tdimword{Nà} \tdimword{pà} \tdimword{ínn-àh} \tdimword{kòtè} \tdimword{giah} \tdimword{nàdìng} \tdimword{à} \tdimword{áwng} \tdimword{dìng} \tdimword{hìam} \\
     \textsc{2sg} father house-\textsc{loc} \textsc{1pl}.\textsc{pro} camp \textsc{purp} \textsc{3sg} open \textsc{irr} \textsc{q} \\
\glt 'Is there room in thy father's house for us to lodge?'
\end{exe}

This printed example should replace the older report paraphrase
\tdim{Inn-ah hong tum theih na hiam}. The dossier is explicit that the
exported window for Genesis 24:23 is
\tdim{Na pa inn-ah kote giah nading a awng ding hiam}, and that attested
clause is what the print slice should use. The candidate TSV also
records one export caveat: punctuation is attached to the final token in
the raw analyzer window, but that punctuation artifact should not be
printed as if it were part of \tdim{hiam} itself.

\subsubsection{\texorpdfstring{WH + \tdim{hiam} content
questions}{WH +  content questions}}\label{wh-content-questions}

The content-question pattern in the current packet is WH + \tdim{hiam}.
The safest print-facing claim is therefore not that every WH form
behaves identically, but that the present candidate-backed examples show
\tdim{bang}, \tdim{kua}, \tdim{bangci}, and \tdim{banghangin} in this
construction.

\begin{exe}
\ex \label{ex:hiam-bang}
\gll \tdimword{Bàng} \tdimword{àhì} \tdimword{hìam?} \\
     like be.\textsc{3sg} \textsc{q} \\
\glt 'What is it?'
\end{exe}

Exodus 16:15 is the current print anchor for \tdim{bang}. The dossier
keeps one export caveat explicit: the analyzer glosses \tdim{bang} here
as \tdim{like}. The clause-level evidence and KJV alignment still
support the ordinary what-question reading, so the example remains
usable in print with that caveat recorded in the packet rather than
foregrounded in the example line.

\begin{exe}
\ex \label{ex:hiam-kua}
\gll \tdimword{Híh-tè} \tdimword{kua} \tdimword{àhì} \tdimword{hìam?} \\
     \textsc{prox}-\textsc{pl} who be.\textsc{3sg} \textsc{q} \\
\glt 'Who are these?'
\end{exe}

Genesis 48:8 is the current print anchor for \tdim{kua ... hiam}, and 2
Samuel 22:32 gives additional support for the same pattern
(\tdim{Topa longal Pasian kua hiam}). The dossier again keeps the
relevant export caveat explicit: \tdim{kua} is tagged as \tdim{NUM} in
the analyzer output. That label should be treated as an export
limitation rather than as evidence against the interrogative reading.

\begin{exe}
\ex \label{ex:hiam-bangci}
\gll \tdimword{Bangci} \tdimword{à} \tdimword{hi-cí} \tdimword{gàm-tat} \tdimword{nà} \tdimword{hì} \tdimword{hìam?} \\
     how \textsc{3sg} \textsc{prox}-say land-rule \textsc{2sg} \textsc{decl} \textsc{q} \\
\glt 'How have you acted thus?' (Genesis 3:13)
\end{exe}

Genesis 3:13 is the cleanest current \tdim{bangci} row, and it matters
methodologically because it should remain visible as \tdim{bangci}, not
be flattened into a generic \tdim{bang} example.

The present packet also keeps Genesis 4:6 as reason-question evidence:

\begin{quote}
bang hangin na mai sia ahi hiam
\end{quote}

This row is important because it keeps \tdim{banghangin} in the slice
while also recording the segmentation caveat from the dossier. The
analyzer exports the sequence as
\tdim{bang | hang-in | na | mai | sia | ahi | hiam}, but the current
packet still treats the clause as curated reason-question evidence
rather than as a raw \tdim{bang} hit.

\subsubsection{Embedded questions}\label{embedded-questions}

The dossier keeps Exodus 16:15 \tdim{bang hiam cih thei lo uh hi}
visible as embedded-question material, but it is not promoted here as a
core independent-clause \tdim{hiam} example. It should remain deferred
for a later treatment of embedded questions or interrogative
complementation, not folded into the first printed outline as if it were
interchangeable with the main clause-final examples above.

\subsubsection{Blocked false friends}\label{blocked-false-friends}

The packet also shows why raw \tdim{hiam} and \tdim{bang} matching would
overgenerate. \tdim{Bang hang hiam cih leh} in Genesis 3:20 is a
formulaic explanatory frame, not an ordinary question example.
Revelation 1:16 \tdim{langnih a hiam namsau} and 2 Kings 11:11
\tdim{a hiam ciat uh} are lexical or non-interrogative \tdim{hiam}
controls rather than question-particle evidence. Likewise,
\tdim{bangmah} and \tdim{bangin} are bang-family false friends and
should not be treated as ordinary \tdim{bang} interrogatives.

These blocked rows matter because they keep the printed claim small and
accurate: the slice prints candidate-backed interrogative examples, not
every surface string that happens to contain \tdim{hiam} or \tdim{bang}.

\subsubsection{Deferred particles}\label{deferred-particles}

Existing generated reports mention \tdim{maw}, \tdim{ham}, and
\tdim{em}, but they remain deferred in this first print slice. They are
not part of the core \tdim{hiam} evidence printed here, and this section
does not yet build a comparison-particle analysis.

\subsubsection{Editorial summary}\label{editorial-summary-6}

This slice can now safely support four modest claims. First, \tdim{hiam}
is the candidate-backed question marker in the current print packet,
with clause-final \tdim{hiam} as the core pattern represented in the
accepted examples. Second, WH + \tdim{hiam} is the current
content-question pattern. Third, \tdim{bang}, \tdim{kua}, \tdim{bangci},
and \tdim{banghangin} are the current candidate-backed WH forms. Fourth,
embedded questions and deferred comparison particles remain outside the
scope of this first printed analysis.

\subsection{Sentence-final particles}\label{sentence-final-particles}

\emph{Source slice:
\texttt{output/publication\_review/grammar\_sentence\_final\_particles\_print\_slice.md}}

\subsubsection{Scope}\label{scope-8}

This is a short print-facing draft section on sentence-final particles
in Tedim Chin, controlled by
\texttt{candidates\_sentence\_final\_particles.tsv} and
\texttt{dossier\_sentence\_final\_particles.md}.

It covers only a small candidate-backed set: caveated declarative
evidence with \tdim{ahi hi}; caveated negative-plus-declarative evidence
with \tdim{thusim lo hi}; optative \tdim{hen} in \tdim{Khuavak om hen};
singular imperative \tdim{in} in \tdim{teembaw khat bawl in}, with
case-marker overlap caveat; and plural imperative \tdim{un} in
\tdim{gingsak un}, currently the cleanest imperative anchor.

It does not yet attempt a full sentence-final particle system, a full
mood or aspect chapter, a full TAM account, a new treatment of
\tdim{hiam}, a settled analysis of \tdim{tahen}, \tdim{aw}, \tdim{ta},
or \tdim{zo}, or generated-report frequency tables. Dictionary and
review-note slices have not yet begun.

\subsubsection{Sentence-final particles in
outline}\label{sentence-final-particles-in-outline}

The current candidate-backed packet supports a narrow generalization.
\tdim{Hi} is visible only through constructionally controlled
\tdim{ahi hi} and \tdim{lo hi} rows. \tdim{Hiam} is visible only as
overlap control because interrogatives are already stabilized.
\tdim{Hen} has one usable optative row. \tdim{In} and \tdim{un}
represent imperative material, but \tdim{in} has case-marker overlap
while \tdim{un} is cleaner. \tdim{Aw} is only vocative or exclamative
boundary material and is not print-ready. \tdim{Tahen}, \tdim{ta}, and
\tdim{zo} remain deferred or needs-review because the export is noisy.
Broad TAM remains deferred.

\subsubsection{\texorpdfstring{Declarative \tdim{hi} with copula
overlap}{Declarative  with copula overlap}}\label{declarative-with-copula-overlap}

Genesis 1:13 supplies the current declarative anchor:

\begin{exe}
\ex \label{ex:sfp-ahi-hi}
\gll \tdimword{àhì} \tdimword{hì} \\
     be.\textsc{3sg} \textsc{decl} \\
\glt 'it was'
\end{exe}

This is the current copula-plus-declarative evidence, but only with
caveat. It is not a bare \tdim{hi} example. The row bundles copular
\tdim{ahi} with final \tdim{hi}, so it does not license raw \tdim{hi}
harvesting or a claim that every \tdim{hi} token is sentence-final
declarative.

\subsubsection{\texorpdfstring{Negative-plus-declarative
\tdim{lo hi}}{Negative-plus-declarative }}\label{negative-plus-declarative}

Genesis 4:5 keeps one negative-plus-declarative environment visible:

\begin{exe}
\ex \label{ex:sfp-lo-hi}
\gll \tdimword{thusim} \tdimword{lò} \tdimword{hì} \\
     parable \textsc{neg} \textsc{decl} \\
\glt 'he had not respect' / 'it was not accepted' (Genesis 4:5)
\end{exe}

This row keeps \tdim{lo hi} visible as a sentence-final environment, but
it overlaps the stabilized negation packet. The present slice should not
reopen \tdim{lo} or build a new negation section here. It is treated as
negation-overlap evidence only.

\subsubsection{\texorpdfstring{\tdim{Hiam} as interrogatives
overlap}{ as interrogatives overlap}}\label{as-interrogatives-overlap}

Deferred / overlap control:

\begin{quote}
\tdim{Hihte kua ahi hiam?}
\end{quote}

Genesis 48:8 is mentioned here only as cross-reference material.
\tdim{Hiam} belongs to the stabilized interrogatives packet, so this
sentence-final particle slice should not reopen or duplicate \tdim{hiam}
analysis. The row is not a print anchor for this slice, and the
quotation or punctuation noise plus the \tdim{kua = NUM} export caveat
should remain visible rather than normalized away.

\subsubsection{\texorpdfstring{Optative
\tdim{hen}}{Optative }}\label{optative}

Genesis 1:3 supplies the current usable optative row:

\begin{exe}
\ex \label{ex:sfp-hen}
\gll \tdimword{Khuavak} \tdimword{ōm} \tdimword{hèn} \\
     light exist \textsc{juss} \\
\glt 'Let there be light' (Genesis 1:3)
\end{exe}

This is the current usable optative evidence with caveat. The candidate
is analyzer-backed as \tdim{om hen}, so the slice should not replace it
with report-style \tdim{ta hen} wording. It also should not turn this
narrow row into a broad mood chapter.

\subsubsection{\texorpdfstring{Imperative \tdim{in} and
\tdim{un}}{Imperative  and }}\label{imperative-and}

Genesis 6:14 keeps one singular imperative row visible:

\begin{exe}
\ex \label{ex:sfp-in}
\gll \tdimword{teembaw} \tdimword{khàt} \tdimword{bāwl} \tdimword{ìn} \\
     ark one make \textsc{erg} \\
\glt 'Make thee an ark' (Genesis 6:14)
\end{exe}

This keeps singular imperative \tdim{in} visible, but only with a
serious case-marker and export caveat: \tdim{in} is exported as
\tdim{ERG} / \tdim{FUNC}. The candidate also has \tdim{teembaw}, not
report-style \tdim{lawng}. The slice should not harvest raw \tdim{in}
hits and should not reopen case marking.

Psalms 100:1 supplies the cleanest current imperative anchor:

\begin{exe}
\ex \label{ex:sfp-un}
\gll \tdimword{ging-sàk} \tdimword{ùn} \\
     sound-\textsc{caus} \textsc{imp}.\textsc{pl} \\
\glt 'Make a joyful noise'
\end{exe}

This is the clean plural-imperative anchor. The span should stay tight
so nearby \tdim{aw} material in the same verse is not absorbed, and the
row should not be used to launch a full imperative paradigm.

\subsubsection{\texorpdfstring{\tdim{Aw} as vocative/exclamative
boundary
material}{ as vocative/exclamative boundary material}}\label{as-vocativeexclamative-boundary-material}

Boundary material:

\begin{quote}
\tdim{Gam khempeuh aw}
\end{quote}

Psalms 100:1 keeps \tdim{aw} visible, but only as vocative or
exclamative boundary material. The analyzer glosses \tdim{aw} as
\tdim{voice} and gives \tdim{N} in \tdim{pos\_span}, and the same verse
also has another \tdim{aw} in \tdim{lungdamna aw}. This is therefore not
settled sentence-final mood evidence, and raw \tdim{aw} hits should not
be treated as sentence-final particles.

\subsubsection{\texorpdfstring{\tdim{Tahen}, \tdim{ta}, and \tdim{zo}
deferred}{, , and  deferred}}\label{and-deferred}

Several report-visible rows remain too noisy to promote:

\begin{itemize}
\tightlist
\item
  \tdim{hi tahen}: deferred because \tdim{tahen} is exported as
  \tdim{army} / \tdim{N}; the slice should not normalize fused
  \tdim{tahen} or split \tdim{ta hen} into a clean jussive example.
\item
  \tdim{mangngilh ta hi}: needs-review TAM-overlap material because
  \tdim{ta} is exported as \tdim{child} / \tdim{FUNC}.
\item
  \tdim{zo}: deferred because the export glosses \tdim{zo} as
  \tdim{south} / \tdim{N} rather than as clean completive evidence.
\end{itemize}

These rows keep report-visible material in view, but they are not
print-ready as jussive, perfective, or completive evidence. Broad TAM
remains deferred.

\subsubsection{Editorial summary}\label{editorial-summary-7}

This slice safely supports five modest claims: \tdim{ahi hi} is usable
as copula-plus-declarative evidence with caveat; \tdim{thusim lo hi} is
usable as negation-plus-declarative evidence with caveat;
\tdim{Khuavak om hen} is usable as optative evidence with caveat;
\tdim{teembaw khat bawl in} is usable as singular imperative evidence
only with a case-overlap caveat; and \tdim{gingsak un} is the clean
plural-imperative anchor.

What remains deferred is equally important: bare \tdim{hi} as a general
declarative particle, \tdim{hiam} as new sentence-final evidence,
\tdim{tahen} as a settled jussive, \tdim{aw} as a settled exclamative or
mood particle, \tdim{ta} and \tdim{zo} as settled aspectual particles,
broad TAM and full mood or aspect chapters, and raw generated-report
counts. Broad TAM, directionals, chrestomathy, Mizo/lus, and other
Kuki-Chin language work remain deferred.

The next step after this grammar slice is the dictionary print slice at
\texttt{output/publication\_review/dictionary\_sentence\_final\_particles\_print\_slice.md}.
Dictionary and review-note slices have not yet begun.

\subsection{Coordinators}\label{coordinators}

\emph{Source slice:
\texttt{output/publication\_review/grammar\_coordinators\_print\_slice.md}}

\subsubsection{Scope}\label{scope-9}

This is a short print-facing draft section on coordinators in Tedim
Chin, controlled by \texttt{candidates\_coordinators.tsv} and
\texttt{dossier\_coordinators.md}.

It covers only a small candidate-backed set: NP coordination with
\tdim{le}, anchored by Genesis 1:1 \tdim{vantung le leitung};
conditional or boundary \tdim{leh}, not clean clause-conjunction
\tdim{leh}; sequential \tdim{a} only as caveated boundary material,
paired with a blocked agreement-\tdim{a} false friend; \tdim{mawh} only
as deferred lexical or analyzer-noise material; \tdim{Ahih hangin} as a
caveated adversative connector; and \tdim{ahih kei leh} as
conditional-adversative boundary material.

It does not yet attempt a full coordination system, a full
clause-linking or converb system, a full temporal or causal subordinator
treatment, a full sentence-final particle treatment, or generated-report
frequency tables. Dictionary and review-note slices have not yet begun.

\subsubsection{Coordinators in outline}\label{coordinators-in-outline}

The current candidate-backed packet supports a narrow generalization.
\tdim{Le} currently provides the safe NP-conjunction anchor. \tdim{Leh}
is visible, but only as conditional or boundary evidence rather than as
print-ready simple clause conjunction. \tdim{A} is visible only as
caveated sequential-linkage boundary evidence and must not be harvested
raw. \tdim{Mawh} remains deferred until a clean disjunction or
alternative-question example is located. \tdim{Ahih hangin} is the
current adversative connector anchor, with an internal-analysis caveat.
\tdim{Ahih kei leh} is conditional-adversative boundary material, not a
simple coordinator.

\subsubsection{\texorpdfstring{NP coordination with
\tdim{le}}{NP coordination with }}\label{np-coordination-with}

Genesis 1:1 supplies the clean print-ready anchor:

\begin{exe}
\ex \label{ex:coord-le-np}
\gll \tdimword{vàn-tùng} \tdimword{lè} \tdimword{léi-tùng} \\
     sky-on and land-on \\
\glt 'heaven and earth'
\end{exe}

This is the current safe NP-conjunction anchor. It supports a modest
print claim that \tdim{le} joins noun phrases.

That claim must remain narrow. The row does \textbf{not} license broad
raw \tdim{le} harvesting, and the current packet does not treat every
\tdim{le} token as coordinator evidence. A separate blocked \tdim{le}
row was not created in the first pass; overgeneration is controlled here
by curated selection rather than by a broad token sweep.

\subsubsection{\texorpdfstring{\tdim{Leh} as conditional or boundary
material}{ as conditional or boundary material}}\label{as-conditional-or-boundary-material}

Genesis 13:9 keeps \tdim{leh} visible only as warning or boundary
evidence:

\begin{exe}
\ex \label{ex:coord-leh-boundary}
\gll \tdimword{véi-lam} \tdimword{nà} \tdimword{lāk} \tdimword{lèh} \tdimword{kèi} \tdimword{taklam-àh} \tdimword{kà} \tdimword{pái} \tdimword{dìng} \tdimword{hì} \\
     faint-direction \textsc{2sg} among and \textsc{neg} right-\textsc{loc} \textsc{1sg} go \textsc{irr} \textsc{decl} \\
\glt 'if thou wilt take the left hand, then I will go to the right' (Genesis 13:9)
\end{exe}

This row keeps \tdim{leh} visible, but not as a clean simple
clause-conjunction anchor. The wider English context is conditional, so
future coordinator prose must not flatten conditional \tdim{leh} into a
simple printed ``\tdim{and}'' analysis.

The row also carries an analyzer or export caveat: \tdim{kei} is glossed
as \tdim{NEG} in the selected window even though the wider English
context is ``I will go.'' That should be recorded as an export caveat,
not used to reopen the pronouns or negation packets.

\subsubsection{\texorpdfstring{Sequential \tdim{a} and
agreement-\tdim{a} false
friends}{Sequential  and agreement- false friends}}\label{sequential-and-agreement--false-friends}

Genesis 2:10 keeps one possible sequential-linkage window visible:

\begin{exe}
\ex \label{ex:coord-a-sequential-boundary}
\gll \tdimword{luang} \tdimword{à} \tdimword{túa} \tdimword{mùn} \tdimword{pàn-ìn} \tdimword{gun} \tdimword{hòng} \tdimword{kikhen-ìn} \\
     flow \textsc{3sg} \textsc{dist} place \textsc{abl}-\textsc{erg} river \textsc{3→1} separate-\textsc{cvb} \\
\glt 'and from thence it was parted' (Genesis 2:10)
\end{exe}

This row keeps possible sequential linkage visible, but it is not
print-ready as a core coordinator example. The analyzer still exports
the relevant \tdim{a} as \tdim{3SG} / \tdim{FUNC}, so the present slice
treats it as warning or boundary evidence only.

Blocked control:

\begin{quote}
\tdim{a piangsak}
\end{quote}

Genesis 1:1 \tdim{a piangsak} is blocked agreement or function material,
not coordinator evidence. This false-friend control prevents raw
\tdim{a} harvesting from flooding the coordinator packet.

\subsubsection{\texorpdfstring{\tdim{Mawh}
deferred}{ deferred}}\label{deferred}

The present slice does not print a positive \tdim{mawh} example.

Deferred control:

\begin{quote}
\tdim{mawh}
\end{quote}

The generated report mentions \tdim{mawh} as disjunction or
alternative-question material, but the current candidate layer only has
Genesis 6:3, where \tdim{mawh} is lexical or analyzer-noise material
glossed as \tdim{sin} / \tdim{V} rather than as disjunction. \tdim{Mawh}
therefore remains not print-ready in this slice.

Report-only schematic examples such as \tdim{mi mawh ganhing mawh} or
\tdim{pai ding mawh om ding mawh} should not be printed unless a later
analyzer-backed row is actually located.

\subsubsection{\texorpdfstring{Adversative
\tdim{Ahih hangin}}{Adversative }}\label{adversative}

Genesis 3:4 supplies the current adversative connector anchor:

\begin{exe}
\ex \label{ex:coord-ahih-hangin}
\gll \tdimword{Àhìh} \tdimword{hàng-ìn} \\
     be.\textsc{3sg}.\textsc{rel} because-\textsc{erg} \\
\glt 'but'
\end{exe}

This row is usable with caveat as the current adversative connector. It
remains internally analyzable as \tdim{ahih} + \tdim{hang-in}, so the
present slice should not open a full causal or subordinator analysis
from this row.

\subsubsection{\texorpdfstring{Conditional-adversative
\tdim{ahih kei leh}}{Conditional-adversative }}\label{conditional-adversative}

Exodus 12:3 keeps one compact boundary row visible:

\begin{exe}
\ex \label{ex:coord-ahih-kei-leh}
\gll \tdimword{àhìh} \tdimword{kèi} \tdimword{lèh} \\
     be.\textsc{3sg}.\textsc{rel} \textsc{neg} and \\
\glt 'otherwise' / 'if not'
\end{exe}

This row is useful as conditional-adversative boundary material, but it
is not a simple coordinator. It overlaps with negation and conditional
\tdim{leh}, yet the slice should not reopen the stabilized negation
packet.

\subsubsection{Deferred material}\label{deferred-material-1}

Several limits remain important at this stage:

\begin{itemize}
\tightlist
\item
  clean simple clause-conjunction \tdim{leh} remains deferred;
\item
  clean \tdim{mawh} disjunction or alternative-question evidence remains
  deferred;
\item
  full sequential-\tdim{a} analysis remains deferred;
\item
  \tdim{ciangin} and broader \tdim{hangin} temporal or causal
  subordinator material remain deferred;
\item
  full clause-chaining or converb coordination remains deferred;
\item
  sentence-final particles remain outside the scope of this slice;
\item
  generated-report raw counts remain excluded.
\end{itemize}

Broad TAM, directionals, chrestomathy, Mizo/lus, and other Kuki-Chin
language work likewise remain deferred.

\subsubsection{Editorial summary}\label{editorial-summary-8}

This slice now safely supports four modest claims: \tdim{le} is the
current NP conjunction anchor, anchored by \tdim{vantung le leitung};
\tdim{Ahih hangin} is usable as an adversative connector with caveat;
\tdim{ahih kei leh} is useful as conditional-adversative boundary
evidence with caveat; and \tdim{leh}, \tdim{a}, plus \tdim{mawh} remain
visible only as warning, deferred, or boundary material.

What remains deferred is equally important: raw frequency counts,
treating every \tdim{le}, \tdim{leh}, or \tdim{a} token as coordinator
evidence, simple clause-conjunction \tdim{leh}, \tdim{mawh} as accepted
disjunction, a full sequential-\tdim{a} analysis, a fuller temporal or
causal subordinator treatment, and sentence-final particles. The next
step after this grammar slice is the dictionary print slice, while
dictionary and review-note work have not yet begun.

\subsection{Reduplication}\label{reduplication}

\emph{Source slice:
\texttt{output/publication\_review/grammar\_reduplication\_print\_slice.md}}

\subsubsection{Editorial scope}\label{editorial-scope-10}

This is the first narrow reduplication grammar slice. It is controlled
by \texttt{output/publication\_review/candidates\_reduplication.tsv} and
\texttt{output/publication\_review/dossier\_reduplication\_scope.md}.
Supporting/background evidence comes from
\texttt{docs/grammar/reports/07-deriv-02-reduplication.md}.

This is not a full derivation chapter, not a full reduplication chapter,
not a dictionary slice, and not a TAM/aspect or VP-structure slice. It
also stays narrow against
\texttt{output/publication\_review/review\_notes\_derivation\_valency.md},
\texttt{output/publication\_review/review\_notes\_nominalization.md},
\texttt{output/publication\_review/review\_notes\_vp\_structure\_stacking.md},
\texttt{output/publication\_review/review\_notes\_noun\_domain.md}, and
\texttt{output/publication\_review/review\_notes\_tam.md}.

The present slice therefore covers only the first safe
full-reduplication claim. No dictionary slice exists for reduplication,
because this packet is grammar-facing and constructional rather than
lexical.

\subsubsection{Full reduplication as
intensification}\label{full-reduplication-as-intensification}

\tdim{mahmah} is the main full-reduplication intensifier anchor.

The controlled segmentation and gloss/function are
\tdim{mah\textasciitilde{}mah} and
\tdim{EMPH\textasciitilde{}EMPH / very, truly}.

The report also supports the worked example \tdim{pha mahmah hi},
glossed there as \tdim{good very DECL}.

\tdim{taktak} is the closest support row for the same small intensifying
pattern.

Its controlled segmentation and gloss/function are
\tdim{tak\textasciitilde{}tak} and
\tdim{TRUE\textasciitilde{}TRUE / truly, certainly}.

Taken together, \tdim{mahmah} and \tdim{taktak} support a narrow
full-reduplication intensifier pattern, not the whole reduplication
system.

\subsubsection{Secondary distributive
evidence}\label{secondary-distributive-evidence}

\tdim{peuhpeuh} remains visible as secondary distributive evidence for
full reduplication.

Its controlled segmentation and gloss/function are
\tdim{peuh\textasciitilde{}peuh} and
\tdim{each\textasciitilde{}each / every, each}.

It stays outside the leading claim because it is more quantifier-like or
noun-modifying than the clean intensifier anchors.

\subsubsection{Boundary material}\label{boundary-material-7}

\tdim{ni ni} stays outside the first grammar slice because syntactic X X
repetition overlaps with temporal/adverbial and VP-structure
interpretation.

\tdim{leuleu} stays outside because iterative or continuative
reduplication overlaps with TAM/aspect and VP structure.

\tdim{gengen} stays outside because the verbal reduplication evidence is
too thin for the first print-facing claim.

\tdim{kawikawi} stays outside because it is expressive,
spatial-totality, or lexicalized-looking material.

\tdim{theithei} stays outside because it is report-only and overlaps
adverbial, modal, or ability-related territory.

\tdim{bangbang}, \tdim{bekbek}, \tdim{zenzen}, \tdim{tuamtuam}, or
similar report-only analysis-table rows stay outside because they are
not yet controlled enough for the first slice.

analyzer-noisy, count-only, or theory-heavy whole-system claims stay
outside because they do not yet produce a safe first print-facing
pattern.

Any broad derivation chapter claim stays outside because this packet is
not yet ready to widen one narrow full-reduplication slice into a full
derivational account.

Any dictionary-entry claim stays outside because this packet is
constructional rather than lexical.

\subsubsection{Safe first-slice claim}\label{safe-first-slice-claim-8}

At the current slice maturity level, the safest reduplication claim is
that Tedim has candidate-controlled evidence for full reduplication used
in intensification, with \tdim{mahmah} / \tdim{mah\textasciitilde{}mah}
as the main anchor and \tdim{taktak} / \tdim{tak\textasciitilde{}tak} as
the closest support row. Distributive \tdim{peuhpeuh} remains visible as
secondary evidence, while syntactic, aspectual, verbal,
lexicalized-looking, and report-only reduplication rows remain
candidate-layer or boundary material.

That claim is deliberately smaller than a full derivation chapter,
smaller than a full reduplication chapter, smaller than a TAM/aspect or
VP-structure slice, and smaller than a dictionary slice.

\subsubsection{Recommended next step}\label{recommended-next-step-10}

This grammar slice is now paired with
\texttt{output/publication\_review/review\_notes\_reduplication.md}, so
the packet is ready for human review at its current
full-reduplication-intensifier slice maturity level.

The next editorial step should be a whole-grammar coverage checkpoint
rather than starting another new packet immediately. If more
reduplication work is chosen after that checkpoint, the next sub-scope
should be distributive \tdim{peuhpeuh} or syntactic \tdim{ni ni}, not
lexicalized-looking or aspect-heavy rows.

\subsection{Broader discourse}\label{broader-discourse}

{[}MAJOR GAP: broader discourse remains partly surfaced and
boundary-heavy.{]}

Current packetized material reaches clause type and sentence-final
behavior, but a broader discourse packet is still only partly surfaced
and remains boundary-heavy.

\subsection{Analyzer-gap caution}\label{analyzer-gap-caution}

{[}MAJOR GAP: analyzer-gap topics remain cross-cutting blockers.{]}

Analyzer-gap topics still cut across tone in \texttt{-a}, conditioned
variants, hong-/kong-, \texttt{-sak}, \texttt{-pih}, and related
cross-packet boundaries, so they remain visible blockers rather than
assembled review prose.

\frontmattersection{End state of this preview}

This assembled review preview contains the actual prose of the current
first-pass publication-review grammar slices in a single ordered draft.
It does not claim that the whole grammar is finished, and the generated
PDF is a review preview PDF rather than a final publication PDF.

\bibliography{../../literature/bibliography.bib}

\end{document}
